% La période allemande — Histoire 1re Tech — Cours signé OnBuch+
% Programme officiel MINESEC. Compiler avec : tectonic N05-periode-allemande.tex
\newcommand{\DOCMATIERE}{Histoire}
\newcommand{\DOCNIVEAU}{1re Tech}
\newcommand{\DOCMODULE}{Histoire – classe de Première technique}
\newcommand{\DOCLECON}{N05}
\newcommand{\DOCTITRE}{La période allemande}
% OnBuch+ — préambule « cours signés » ENRICHI (compilé avec tectonic / XeLaTeX)
% Système visuel « Pop » d'OnBuch+ : fond crème (jamais blanc), bordures encre
% épaisses, ombres « dures » décalées, titres Archivo Black en capitales.
% Version MATHÉMATIQUES : graphes (pgfplots/tikz), boîtes pédagogiques
% (définition, propriété, méthode, attention, le savais-tu, exercice résolu,
% exercice, TP, bilan), page de garde et signature OnBuch+ ancrée en bas.
%
% Macros attendues dans main.tex AVANT \input{preamble} :
%   \DOCMATIERE  \DOCNIVEAU  \DOCMODULE  \DOCLECON (numéro)  \DOCTITRE
\documentclass[11pt]{article}

\usepackage{fontspec}
\usepackage[french]{babel}
\usepackage[a4paper,margin=2cm,top=3.4cm,bottom=2.8cm,headheight=1.4cm,footskip=1.3cm]{geometry}
\usepackage{amsmath,amssymb,amsfonts}
\usepackage{mathtools} % \xmapsto, \coloneqq... souvent utilisés par les modèles
\usepackage{cancel} % simplifications barrées (bilans de forces)
% \abs{z} (module, valeur absolue) et \norm{u} (norme), utilisés par les modèles
\providecommand{\abs}{}\let\abs\relax\DeclarePairedDelimiter{\abs}{\lvert}{\rvert}
\providecommand{\norm}{}\let\norm\relax\DeclarePairedDelimiter{\norm}{\lVert}{\rVert}
\usepackage[table]{xcolor} % [table] : fournit aussi \rowcolors (alternance de lignes)
\usepackage{pagecolor}
\usepackage{tikz}
\usetikzlibrary{arrows.meta,positioning,calc,shapes.geometric,shapes.misc,shapes.arrows,decorations.pathmorphing,decorations.pathreplacing,decorations.markings,patterns,shadows,mindmap,trees,fit,backgrounds,matrix,intersections,angles,quotes}
\usepackage{pgfplots}
\usepackage[version=4]{mhchem} % équations nucléaires \ce{^{235}_{92}U}
\usepackage[european,siunitx]{circuitikz} % schémas électriques
\pgfplotsset{compat=1.18}
\usepgfplotslibrary{fillbetween,groupplots} % aires entre courbes (calcul intégral)
\usepackage{enumitem}
\usepackage{fancyhdr}
% [most] charge toutes les bibliothèques tcolorbox (listings, minted, raster,
% documentation...) et gonfle inutilement la mémoire TeX sur un document long
% (~300 boîtes) : on ne charge que ce qui est réellement utilisé.
\usepackage[skins,breakable]{tcolorbox}
\usepackage{graphicx}
\usepackage{colortbl}
\usepackage{booktabs}
\usepackage{tabularx}
\usepackage{multirow}
\usepackage{diagbox} % case d'en-tête diagonale des tableaux à double entrée
% Colonnes extensibles centrée / gauche / droite (C, L, R), souvent utilisées par les modèles
\newcolumntype{C}{>{\centering\arraybackslash}X}
\newcolumntype{L}{>{\raggedright\arraybackslash}X}
\newcolumntype{R}{>{\raggedleft\arraybackslash}X}
\usepackage{array}
\usepackage{multicol}
\usepackage{siunitx}
% parse-numbers=false : accepte aussi \SI{2,5 \times 10^{-3}}{...}
\sisetup{locale=FR,output-decimal-marker={,},group-minimum-digits=4,parse-numbers=false}
\DeclareSIUnit\ohm{\ensuremath{\symup{\Omega}}} % U+2126 absent de la police
\DeclareSIUnit\division{div} % réglages d'oscilloscope (ms/div, V/div)
\DeclareSIUnit\tr{tr} % tours (vitesse de rotation, ex. \SI{1500}{\tr\per\minute})
\DeclareSIUnit\molar{M} % concentration molaire (ex. \SI{3}{\molar}, \SI{1}{\micro\molar})
\DeclareSIUnit\an{an} % année (le modèle confond parfois avec \year, un compteur TeX) — ex. \SI{5}{\kilo\meter\per\an}
\DeclareSIUnit\FCFA{FCFA} % franc CFA (ex. \SI{500}{\FCFA\per\kilo\gram})
\DeclareSIUnit\degreeBrix{\textdegree Brix} % teneur en sucre d'un sirop/jus (réfractométrie)
\DeclareSIUnit\month{mois} % le modèle utilise parfois \month, un compteur TeX, comme unité de durée
\DeclareSIUnit\microliter{\micro\liter} % le modèle écrit parfois \microliter en un seul token
\DeclareSIUnit\million{millions} % dénombrement (ex. \SI{15}{\million\per\mL} spermatozoïdes)
\usepackage{qrcode}
% \degree : commande fréquemment utilisée par les modèles pour un angle ou une
% température en dehors de \SI{}{} (non fournie par les paquets chargés).
\providecommand{\degree}{^{\circ}}
% \qed (fin de démonstration), utilisé par les modèles sans amsthm
\providecommand{\qed}{\hfill$\square$}
% \barre : double barre des valeurs interdites dans un tableau de variations (tkz-tab)
\providecommand{\barre}{\ensuremath{\|}}
% environnement proof (amsthm non chargé) utilisé par les modèles
\ifdefined\proof\else\newenvironment{proof}[1][Démonstration]{\par\noindent\textit{#1.}\ }{\qed\par}\fi
\usepackage{lastpage}
\usepackage{hyperref}
\hypersetup{hidelinks}

% ============================================================
% Palette « Pop » OnBuch+ — mêmes hex que la classe P (plus_ui.dart)
% ============================================================
\definecolor{poporange}{HTML}{F59321}
\definecolor{popdark}{HTML}{E07A0C}
\definecolor{popcream}{HTML}{FFF6E8}
\definecolor{popink}{HTML}{1C1714}
\definecolor{poppurple}{HTML}{7A5AE0}
\definecolor{popgreen}{HTML}{2E9E6A}
\definecolor{poppink}{HTML}{E0457B}
\definecolor{popblue}{HTML}{2F7DE1}
\definecolor{popgold}{HTML}{C98A12}
\definecolor{popmuted}{HTML}{8A7F76}
\definecolor{popline}{HTML}{EADFD2}
\definecolor{popgray}{HTML}{8A7F76}
\colorlet{popgrey}{popgray}
% Alias tolérants (le modèle nomme parfois les couleurs différemment)
\colorlet{popwhite}{white}
\colorlet{popblack}{popink}
\colorlet{popred}{poppink}
\colorlet{popyellow}{popgold}
% Teintes claires (fonds de tableaux/encadrés) — le modèle nomme parfois
% les couleurs avec un suffixe L (ex. popblueL) sans les avoir définies.
\colorlet{poporangeL}{poporange!20}
\colorlet{popdarkL}{popdark!20}
\colorlet{poppurpleL}{poppurple!20}
\colorlet{popgreenL}{popgreen!20}
\colorlet{poppinkL}{poppink!20}
\colorlet{popblueL}{popblue!20}
\colorlet{popgoldL}{popgold!20}
\colorlet{popredL}{popred!20}
\colorlet{popgrayL}{popgray!20}
% Teintes claires pour fonds de tableaux / zones
\colorlet{popblueL}{popblue!10!white}
\colorlet{poporangeL}{poporange!14!white}
\colorlet{popgreenL}{popgreen!12!white}
\colorlet{poppurpleL}{poppurple!12!white}
\colorlet{poppinkL}{poppink!12!white}
\colorlet{popgoldL}{popgold!14!white}

\pagecolor{popcream}

% ============================================================
% Typographie
% ============================================================
\defaultfontfeatures{Ligatures=TeX}
\setmainfont{PlusJakartaSans-Medium.ttf}[Path=../../../fonts/,BoldFont=PlusJakartaSans-Bold.ttf]
% Maths en sans-serif (Fira Math) avec chiffres et lettres droites de Plus
% Jakarta, pour que les formules se fondent dans le texte.
\usepackage[math-style=ISO,bold-style=ISO]{unicode-math}
\setmathfont{FiraMath-Regular.otf}[Path=../../../fonts/]
\setmathfont{PlusJakartaSans-Medium.ttf}[Path=../../../fonts/,range={up/num,up/{Latin,latin},\mathup}]
\usepackage{icomma} % virgule décimale sans espace en mode math
% Caractères mathématiques Unicode absents des polices de texte (Plus Jakarta,
% Archivo Black) : rendus par leur équivalent mathématique, en texte comme en
% formule — sinon ils s'affichent en carré vide (ex. « Arithmétique dans ℤ »).
\usepackage{newunicodechar}
\newunicodechar{ℝ}{\ensuremath{\mathbb{R}}}
\newunicodechar{ℤ}{\ensuremath{\mathbb{Z}}}
\newunicodechar{ℂ}{\ensuremath{\mathbb{C}}}
\newunicodechar{ℕ}{\ensuremath{\mathbb{N}}}
\newunicodechar{ℚ}{\ensuremath{\mathbb{Q}}}
\newunicodechar{𝔻}{\ensuremath{\mathbb{D}}}
\newunicodechar{α}{\ensuremath{\alpha}}
\newunicodechar{β}{\ensuremath{\beta}}
\newunicodechar{γ}{\ensuremath{\gamma}}
\newunicodechar{δ}{\ensuremath{\delta}}
\newunicodechar{ε}{\ensuremath{\varepsilon}}
\newunicodechar{θ}{\ensuremath{\theta}}
\newunicodechar{λ}{\ensuremath{\lambda}}
\newunicodechar{μ}{\ensuremath{\mu}}
\newunicodechar{σ}{\ensuremath{\sigma}}
\newunicodechar{τ}{\ensuremath{\tau}}
\newunicodechar{φ}{\ensuremath{\varphi}}
\newunicodechar{ψ}{\ensuremath{\psi}}
\newunicodechar{ω}{\ensuremath{\omega}}
\newunicodechar{Σ}{\ensuremath{\Sigma}}
% majuscules produites par \MakeUppercase dans les titres (α-aminés -> Α)
\newunicodechar{Α}{\ensuremath{\alpha}}
\newunicodechar{Β}{\ensuremath{\beta}}
\newunicodechar{Γ}{\ensuremath{\Gamma}}
\newunicodechar{Δ}{\ensuremath{\Delta}}
\newunicodechar{Ε}{\ensuremath{\varepsilon}}
\newunicodechar{Θ}{\ensuremath{\Theta}}
\newunicodechar{Λ}{\ensuremath{\Lambda}}
\newunicodechar{Μ}{\ensuremath{\mu}}
\newunicodechar{Τ}{\ensuremath{\tau}}
\newunicodechar{Φ}{\ensuremath{\Phi}}
\newunicodechar{Ψ}{\ensuremath{\Psi}}
\newunicodechar{Ω}{\ensuremath{\Omega}}
\newunicodechar{↦}{\ensuremath{\mapsto}}
\newunicodechar{∈}{\ensuremath{\in}}
\newunicodechar{∉}{\ensuremath{\notin}}
\newunicodechar{∧}{\ensuremath{\wedge}}
\newunicodechar{⇔}{\ensuremath{\Leftrightarrow}}
\newunicodechar{⟺}{\ensuremath{\Longleftrightarrow}}
\newunicodechar{⇒}{\ensuremath{\Rightarrow}}
\newunicodechar{⟹}{\ensuremath{\Longrightarrow}}
\newunicodechar{∘}{\ensuremath{\circ}}
\newunicodechar{∪}{\ensuremath{\cup}}
\newunicodechar{∩}{\ensuremath{\cap}}
\newunicodechar{≡}{\ensuremath{\equiv}}
\newunicodechar{⊂}{\ensuremath{\subset}}
\newunicodechar{⊥}{\ensuremath{\perp}}
\newunicodechar{∃}{\ensuremath{\exists}}
\newunicodechar{∀}{\ensuremath{\forall}}
\newunicodechar{∛}{\ensuremath{\sqrt[3]{\,}}}
\newunicodechar{∜}{\ensuremath{\sqrt[4]{\,}}}
\newunicodechar{⃗}{\ensuremath{{}^{\to}}}
\newunicodechar{ⁿ}{\ifmmode^{n}\else\textsuperscript{\ensuremath{n}}\fi}
\newunicodechar{ˣ}{\ifmmode^{x}\else\textsuperscript{\ensuremath{x}}\fi}
\newunicodechar{ᵇ}{\ifmmode^{b}\else\textsuperscript{\ensuremath{b}}\fi}
\newunicodechar{ʳ}{\ifmmode^{r}\else\textsuperscript{\ensuremath{r}}\fi}
\newunicodechar{ᵘ}{\ifmmode^{u}\else\textsuperscript{\ensuremath{u}}\fi}
\newunicodechar{ʸ}{\ifmmode^{y}\else\textsuperscript{\ensuremath{y}}\fi}
\newunicodechar{ᵃ}{\ifmmode^{a}\else\textsuperscript{\ensuremath{a}}\fi}
\newunicodechar{ᵉ}{\ifmmode^{e}\else\textsuperscript{\ensuremath{e}}\fi}
\newunicodechar{ᵏ}{\ifmmode^{k}\else\textsuperscript{\ensuremath{k}}\fi}
\newunicodechar{ᵖ}{\ifmmode^{p}\else\textsuperscript{\ensuremath{p}}\fi}
\newunicodechar{ᴺ}{\ifmmode^{N}\else\textsuperscript{\ensuremath{N}}\fi}
\newunicodechar{ᶜ}{\ifmmode^{c}\else\textsuperscript{\ensuremath{c}}\fi}
\newunicodechar{ʰ}{\ifmmode^{h}\else\textsuperscript{\ensuremath{h}}\fi}
\newunicodechar{ᵠ}{\ifmmode^{q}\else\textsuperscript{\ensuremath{q}}\fi}
\newunicodechar{ₙ}{\ifmmode_{n}\else\textsubscript{\ensuremath{n}}\fi}
\newunicodechar{ₐ}{\ifmmode_{a}\else\textsubscript{\ensuremath{a}}\fi}
\newunicodechar{ᵢ}{\ifmmode_{i}\else\textsubscript{\ensuremath{i}}\fi}
\newunicodechar{ᵣ}{\ifmmode_{r}\else\textsubscript{\ensuremath{r}}\fi}
\newunicodechar{ₖ}{\ifmmode_{k}\else\textsubscript{\ensuremath{k}}\fi}
\newunicodechar{ᵦ}{\ifmmode_{\beta}\else\textsubscript{\ensuremath{\beta}}\fi}
\newunicodechar{ₓ}{\ifmmode_{x}\else\textsubscript{\ensuremath{x}}\fi}
\newfontfamily\popdisplay{ArchivoBlack-Regular.ttf}[Path=../../../fonts/]
\newfontfamily\poplogo{SpaceGrotesk-Bold.ttf}[Path=../../../fonts/]
\newfontfamily\popsemi{PlusJakartaSans-SemiBold.ttf}[Path=../../../fonts/]
\newfontfamily\popxbold{PlusJakartaSans-ExtraBold.ttf}[Path=../../../fonts/]
\newfontfamily\popxboldls{PlusJakartaSans-ExtraBold.ttf}[Path=../../../fonts/,LetterSpace=3.0]

\setlist[enumerate]{leftmargin=1.7em,itemsep=5pt,topsep=4pt}
\setlist[itemize]{leftmargin=1.7em,itemsep=4pt,topsep=4pt,label={\color{poporange}\textbullet}}
\setlength{\parindent}{0pt}
\setlength{\parskip}{5pt}
\renewcommand{\arraystretch}{1.25}

% pgfplots : style Pop par défaut
\pgfplotsset{
  every axis/.append style={axis line style={popink,line width=1pt},tick style={popink},
    grid style={popline,line width=0.5pt},label style={font=\small},tick label style={font=\footnotesize}},
  popcurve/.style={poporange,line width=1.8pt,no markers,smooth},
}
% Même style disponible dans un \draw TikZ simple (hors pgfplots)
\tikzset{popcurve/.style={poporange,line width=1.8pt}}
\tikzset{popgrid/.style={popline,very thin}}
\colorlet{popcurve}{poporange} % le nom du style est parfois utilisé comme couleur

% ============================================================
% Pill / squiggle
% ============================================================
\newcommand{\poppill}[3]{%
  \tikz[baseline=-0.55ex]\node[draw=popink,line width=1.2pt,rounded corners=2.6pt,%
    fill=#2,inner xsep=6.5pt,inner ysep=3pt]{%
    \popxboldls\fontsize{7}{7}\selectfont\color{#3}\MakeUppercase{#1}};%
}
\newcommand{\popsquiggle}[1]{%
  \tikz[baseline=-0.4ex]{
    \draw[#1,line width=2.6pt,line cap=round]
      plot [smooth] coordinates {(0,0.09) (0.34,0.01) (0.68,0.21) (1.02,0.01) (1.36,0.10)};
  }%
}
% Mot-clé surligné (marqueur orange sous le texte)
\newcommand{\cle}[1]{{\popxbold\color{popdark}#1}}

% ============================================================
% En-tête / pied
% ============================================================
\pagestyle{fancy}
\fancyhf{}
\renewcommand{\headrulewidth}{0pt}
\renewcommand{\footrulewidth}{0pt}
\lhead{\raisebox{-0.15ex}{\poplogo\fontsize{13}{13}\selectfont\color{popink}OnBuch\color{poporange}+}}
\rhead{\poppill{\DOCMATIERE\ \textbullet\ \DOCNIVEAU\ \textbullet\ Leçon \DOCLECON}{white}{popink}}
\lfoot{\popxbold\fontsize{7.6}{7.6}\selectfont\color{popmuted}\MakeUppercase{Cours signé OnBuch+}\ \textbullet\ {\popsemi\DOCTITRE}}
\rfoot{\tikz[baseline=-0.6ex]\node[rounded corners=8.5pt,fill=popink,minimum height=17pt,minimum width=17pt,inner xsep=5pt,inner sep=0]{\popxbold\fontsize{8}{8}\selectfont\color{popcream}\thepage\,/\,\pageref*{LastPage}};}

% ============================================================
% Titres
% ============================================================
\newcommand{\chaptitre}[2]{%
  \begin{tcolorbox}[enhanced,colback=poporange,colframe=popink,boxrule=2.6pt,arc=11pt,
      drop shadow=popink,left=15pt,right=15pt,top=12pt,bottom=11pt,before skip=3pt,after skip=18pt]
    \poppill{#2}{popink}{popcream}\par\vspace{9pt}
    {\raggedright\hyphenpenalty=10000\exhyphenpenalty=10000\popdisplay\fontsize{22}{24}\selectfont\color{popcream}\MakeUppercase{#1}\par}
  \end{tcolorbox}
}
% Section numérotée : pastille orange avec le numéro + titre Archivo
\newcounter{popsec}
\newcommand{\coursec}[1]{%
  \stepcounter{popsec}\setcounter{popsub}{0}%
  \par\vspace{16pt}\noindent
  \begin{minipage}{\linewidth}
  \tikz[baseline=-0.7ex]\node[draw=popink,line width=1.6pt,fill=poporange,rounded corners=4pt,minimum size=22pt,inner sep=0]{\popdisplay\fontsize{12}{12}\selectfont\color{popcream}\thepopsec};%
  \hspace{8pt}{\popdisplay\fontsize{15}{17}\selectfont\color{popink}\MakeUppercase{#1}}\par
  \vspace{4pt}\noindent{\color{poporange}\rule{36pt}{3.6pt}}
  \end{minipage}\par\nopagebreak\vspace{8pt}
}
\newcounter{popsub}
\newcommand{\courssub}[1]{%
  \stepcounter{popsub}%
  \par\vspace{9pt}\noindent{\popxbold\fontsize{11.5}{13}\selectfont\color{popink}{\color{poporange}\thepopsec.\thepopsub}\hspace{6pt}#1}\par\nopagebreak\vspace{4pt}
}

% ============================================================
% Boîtes pédagogiques — même recette : fond blanc, bordure épaisse colorée,
% ombre dure encre, badge plein. Titre optionnel [#1].
% ============================================================
\tcbset{popbase/.style={breakable,enhanced,colback=white,boxrule=2.3pt,arc=9pt,
  shadow={3pt}{-3pt}{0pt}{popink},left=14pt,right=14pt,top=11pt,bottom=11pt,before skip=11pt,after skip=15pt,
  fontupper=\popsemi\fontsize{10.6}{14.5}\selectfont\color{popink},
  fontlower=\popsemi\fontsize{10.6}{14.5}\selectfont\color{popink}}}
% Coche dessinée (le glyphe ✓ n'existe pas dans les polices du document)
\newcommand{\popcheck}[1]{\tikz[baseline=-0.5ex]\draw[#1,line width=1.6pt,line cap=round,line join=round](-0.13,0.0)--(-0.04,-0.09)--(0.13,0.1);}
\providecommand{\coche}{\popcheck{popgreen}}
\providecommand{\resultat}[1]{\par\smallskip\noindent\fcolorbox{popgreen}{popgreenL}{\parbox{\dimexpr\linewidth-2\fboxsep-2\fboxrule\relax}{\textbf{Résultat.} #1}}\par\smallskip}
\newcommand{\popboxhead}[3]{\poppill{#1}{#2}{white}\ifx\relax#3\relax\else\hspace{6pt}{\popxbold\fontsize{10.6}{12}\selectfont\color{popink}#3}\fi\par\vspace{7pt}}

\newtcolorbox{aretenir}[1][]{popbase,colframe=popgreen,before upper={\popboxhead{À retenir}{popgreen}{#1}}}
\newtcolorbox{exemplebox}[1][]{popbase,colframe=poporange,before upper={\popboxhead{Exemple}{poporange}{#1}}}
\newtcolorbox{definition}[1][]{popbase,colframe=popblue,before upper={\popboxhead{Définition}{popblue}{#1}}}
\newtcolorbox{propriete}[1][]{popbase,colframe=poppurple,before upper={\popboxhead{Propriété}{poppurple}{#1}}}
\newtcolorbox{methode}[1][]{popbase,colframe=popgold,before upper={\popboxhead{Méthode}{popgold}{#1}}}
\newtcolorbox{attention}[1][]{popbase,colframe=poppink,before upper={\popboxhead{Attention}{poppink}{#1}}}
\newtcolorbox{experience}[1][]{popbase,colframe=popblue,colback=popblueL,before upper={\popboxhead{Expérience}{popblue}{#1}}}
% « Le savais-tu ? » : carte violette pleine (contexte, culture, Cameroun)
\newtcolorbox{savaistu}[1][]{popbase,colback=poppurple,colframe=popink,
  fontupper=\popsemi\fontsize{10.4}{14.2}\selectfont\color{white},
  before upper={\poppill{Le savais-tu\,?}{white}{poppurple}\ifx\relax#1\relax\else\hspace{6pt}{\popxbold\color{white}#1}\fi\par\vspace{7pt}}}
% Exercice résolu : énoncé en haut, \tcblower puis la solution sur fond crème
\newcounter{popexo}\newcounter{popexr}
\newtcolorbox{exoresolu}[1][]{popbase,colframe=poporange,colbacklower=popcream,
  segmentation style={popink,line width=1.2pt,dashed},
  before upper={\stepcounter{popexr}\popboxhead{Exercice résolu \thepopexr}{poporange}{#1}},
  before lower={\poppill{Solution}{popink}{popcream}\par\vspace{6pt}}}
% Exercice (non résolu en place) — la correction est dans la partie Corrigés
\newtcolorbox{exercice}[1][]{popbase,colframe=popink,
  before upper={\stepcounter{popexo}\popboxhead{Exercice \thepopexo}{popink}{#1}}}
% Corrigé
\newtcolorbox{corrige}[1][]{popbase,colframe=popgreen,colback=popgreenL,
  before upper={\popboxhead{Corrigé}{popgreen}{#1}}}
% Objectifs / prérequis / bilan
\newtcolorbox{objectifs}{popbase,colframe=popink,colback=poporangeL,
  before upper={\popboxhead{Objectifs de la leçon}{popink}{}}}
\newtcolorbox{prerequis}{popbase,colframe=popink,colback=popblueL,
  before upper={\popboxhead{Prérequis}{popblue}{}}}
\newtcolorbox{bilan}{popbase,colframe=popink,colback=white,boxrule=2.8pt,
  before upper={\popboxhead{Fiche bilan}{poporange}{L'essentiel à connaître pour l'examen}}}
% Figure encadrée avec légende
\newtcolorbox{popfigure}[1][]{enhanced,breakable=false,colback=white,colframe=popline,boxrule=1.4pt,arc=8pt,
  left=10pt,right=10pt,top=10pt,bottom=8pt,before skip=10pt,after skip=12pt,halign=center,
  lower separated=false,halign lower=center,
  fontlower=\popsemi\fontsize{9}{11.5}\selectfont\color{popmuted},
  #1}
\newcommand{\legende}[1]{\par\vspace{6pt}{\popsemi\fontsize{9}{11.5}\selectfont\color{popmuted}#1}}

% ============================================================
% Signature OnBuch+
% ============================================================
\newcommand{\poplogoinline}[1]{{\poplogo\fontsize{#1}{#1}\selectfont\color{popink}OnBuch\color{poporange}+}}
\newcommand{\signatureonbuch}{%
  \begin{tcolorbox}[enhanced,colback=popink,colframe=popink,boxrule=2.2pt,arc=10pt,
      drop shadow=poporange,left=14pt,right=14pt,top=10pt,bottom=10pt,before skip=0pt,after skip=0pt]
    \tikz[baseline=-0.7ex]\node[circle,fill=popgreen,draw=popcream,line width=1.3pt,minimum size=22pt,inner sep=0]{\popcheck{white}};%
    \hspace{9pt}\begin{minipage}[c]{0.62\linewidth}
      {\popxbold\fontsize{12}{13}\selectfont\color{popcream}Signé\ \textbullet\ {\poplogo OnBuch\color{poporange}+}}\par\vspace{2pt}
      {\raggedright\popsemi\fontsize{8}{10.5}\selectfont\color{popline}\MakeUppercase{Équipe pédagogique OnBuch+}\\\MakeUppercase{Conforme au programme officiel MINESEC}\par}
    \end{minipage}\hfill
    {\poplogo\fontsize{22}{22}\selectfont\color{popcream}OnBuch\color{poporange}+}
  \end{tcolorbox}%
}

% ============================================================
% Page de garde
% #1 titre · #2 matière · #3 niveau · #4 module · #5 n° leçon · #6 durée ·
% #7 description / objectifs (texte ou itemize)
% ============================================================
\newcommand{\pagegarde}[7]{%
  \thispagestyle{empty}
  \newgeometry{a4paper,margin=1.7cm,top=1.6cm,bottom=1.4cm}%
  \begin{tikzpicture}[remember picture,overlay]
    \fill[poporange!18] ([xshift=-3.2cm,yshift=-3.6cm]current page.north east) circle (4.6cm);
    \fill[poppurple!14] ([xshift=2.4cm,yshift=7.6cm]current page.south west) circle (3.8cm);
  \end{tikzpicture}%
  {\poplogo\fontsize{30}{30}\selectfont\color{popink}OnBuch\color{poporange}+}\hfill
  \raisebox{0.6ex}{\poppill{Cours signé \textbullet\ Premium}{popink}{popcream}}\par
  \vspace{1pt}\popsquiggle{poppurple}\par
  \vspace{8pt}
  \begin{tcolorbox}[enhanced,colback=poporange,colframe=popink,boxrule=2.8pt,arc=13pt,
      drop shadow=popink,left=18pt,right=18pt,top=12pt,bottom=13pt,after skip=10pt]
    \poppill{#2 \textbullet\ #3}{popink}{popcream}\hspace{5pt}\poppill{#4}{white}{popink}\par\vspace{8pt}
    {\popdisplay\fontsize{14}{14}\selectfont\color{popink}LEÇON #5}\par\vspace{4pt}
    {\raggedright\hyphenpenalty=10000\exhyphenpenalty=10000\popdisplay\fontsize{25}{27}\selectfont\color{popcream}\MakeUppercase{#1}\par}
    \vspace{8pt}
    {\popxbold\fontsize{9.5}{9.5}\selectfont\color{popink}\MakeUppercase{Durée conseillée : #6}}
  \end{tcolorbox}
  \begin{tcolorbox}[enhanced,colback=white,colframe=popink,boxrule=2.2pt,arc=10pt,
      drop shadow=popink,left=16pt,right=16pt,top=10pt,bottom=10pt,after skip=8pt]
    \poppill{Dans cette leçon}{poppurple}{white}\par\vspace{5pt}
    \popsemi\fontsize{9.3}{12.6}\selectfont\color{popink}#7
  \end{tcolorbox}
  \noindent\begin{minipage}[t]{0.487\linewidth}
    \begin{tcolorbox}[enhanced,colback=popgreen,colframe=popink,boxrule=2.2pt,arc=10pt,
        drop shadow=popink,left=11pt,right=11pt,top=10pt,bottom=10pt,height=3.1cm,valign=center]
      \tcbox[colback=white,colframe=popink,boxrule=1.3pt,arc=5pt,boxsep=0pt,nobeforeafter,
        left=3pt,right=3pt,top=3pt,bottom=3pt,valign=center]{\qrcode[height=1.9cm]{https://play.google.com/store/apps/details?id=com.luvvix.onbuch&hl=fr}}%
      \hspace{8pt}\begin{minipage}[c]{0.5\linewidth}
        {\popdisplay\fontsize{11}{12.5}\selectfont\color{white}\MakeUppercase{Télécharge OnBuch}}\par\vspace{4pt}
        {\raggedright\popsemi\fontsize{8.4}{11}\selectfont\color{white}Gratuit \textbullet\ Google Play\\Tuteur IA, annales, concours, résultats\par}
      \end{minipage}
    \end{tcolorbox}
  \end{minipage}\hfill
  \begin{minipage}[t]{0.487\linewidth}
    \begin{tcolorbox}[enhanced,colback=poppurple,colframe=popink,boxrule=2.2pt,arc=10pt,
        drop shadow=popink,left=11pt,right=11pt,top=10pt,bottom=10pt,height=3.1cm,valign=center]
      \tikz[baseline=-0.7ex]\node[circle,fill=white,draw=popink,line width=1.3pt,minimum size=1.6cm,inner sep=0]{\popdisplay\fontsize{24}{24}\selectfont\color{poppurple}+};%
      \hspace{8pt}\begin{minipage}[c]{0.6\linewidth}
        {\popdisplay\fontsize{11}{12.5}\selectfont\color{white}\MakeUppercase{Passe à OnBuch+}}\par\vspace{4pt}
        {\raggedright\popsemi\fontsize{8.4}{11}\selectfont\color{white}Tous les cours signés, Tuteur IA illimité, fiches PDF — onglet OnBuch+ de l'app\par}
      \end{minipage}
    \end{tcolorbox}
  \end{minipage}
  \vspace{8pt}
  \popcontenu
  \vfill
  \signatureonbuch
  \restoregeometry
  \newpage
}

% Rangée « Ce cours contient » (page de garde)
\newcommand{\poptile}[3]{% #1 couleur #2 gros texte #3 légende
  \begin{tcolorbox}[enhanced,colback=#1,colframe=popink,boxrule=2pt,arc=9pt,shadow={2.5pt}{-2.5pt}{0pt}{popink},
      left=6pt,right=6pt,top=9pt,bottom=9pt,halign=center,valign=center,height=1.75cm,width=\linewidth,nobeforeafter]
    {\popdisplay\fontsize{12.5}{13}\selectfont\color{white}\MakeUppercase{#2}}\par\vspace{3pt}
    {\centering\popsemi\fontsize{7.8}{9.5}\selectfont\color{white}#3\par}
  \end{tcolorbox}}
\newcommand{\popcontenu}{%
  \par\noindent{\popxboldls\fontsize{7.5}{7.5}\selectfont\color{popmuted}CE COURS SIGNÉ CONTIENT}\par\vspace{7pt}
  \noindent\begin{minipage}[t]{0.235\linewidth}\poptile{popblue}{Cours}{complet, illustré, pas à pas}\end{minipage}\hfill
  \begin{minipage}[t]{0.235\linewidth}\poptile{poporange}{Méthodes}{et exercices résolus}\end{minipage}\hfill
  \begin{minipage}[t]{0.235\linewidth}\poptile{poppink}{Exercices}{type examen + corrigés}\end{minipage}\hfill
  \begin{minipage}[t]{0.235\linewidth}\poptile{popgreen}{Fiche bilan}{l'essentiel à retenir}\end{minipage}\par}

% Sommaire
\newtcolorbox{sommaire}{enhanced,colback=white,colframe=popink,boxrule=2.2pt,arc=10pt,
  drop shadow=popink,left=18pt,right=18pt,top=13pt,bottom=13pt,before skip=6pt,after skip=20pt,
  before upper={\poppill{Sommaire}{poppurple}{white}\par\vspace{9pt}\popsemi\fontsize{10.6}{14.5}\selectfont\color{popink}}}

% Signature de fin de cours — poussée tout en bas de la dernière page
\newcommand{\findecours}{%
  \par\vfill\nopagebreak
  \begin{center}\popsquiggle{poporange}\end{center}\vspace{4pt}
  \signatureonbuch
}
\AtBeginDocument{\def\llbracket{\mathopen{[\mkern-3.5mu[}}\def\rrbracket{\mathclose{]\mkern-3.5mu]}}} % intervalles d'entiers (glyphe ⟦ absent de la police)
% Glyphes absents de Fira Math : redéfinis à partir de glyphes disponibles
\AtBeginDocument{%
  \def\setminus{\mathbin{\backslash}}\let\smallsetminus\setminus
  \def\vdots{\mathord{\vbox{\baselineskip3.2pt\lineskiplimit0pt\kern4pt\hbox{.}\hbox{.}\hbox{.}}}}%
  \def\ddots{\mathinner{\mkern1mu\raise6.5pt\hbox{.}\mkern2mu\raise3.6pt\hbox{.}\mkern2mu\raise0.7pt\hbox{.}\mkern1mu}}%
  \def\triangle{\mathord{\scalebox{0.9}{$\symup{\Delta}$}}}%
  \def\nabla{\mathord{\raisebox{\depth}{\scalebox{1}[-1]{$\symup{\Delta}$}}}}%
  \def\checkmark{\popcheck{popdark}}%
  \def\star{\mathbin{\ast}}\def\bigstar{\mathord{\ast}}%
  \def\diamond{\mathbin{\scalebox{0.6}{$\Diamond$}}}%
  \def\bigcup{\mathop{\mathchoice{\raisebox{-0.3em}{\scalebox{1.7}{$\cup$}}}{\scalebox{1.25}{$\cup$}}{\cup}{\cup}}\displaylimits}%
  \def\bigcap{\mathop{\mathchoice{\raisebox{-0.3em}{\scalebox{1.7}{$\cap$}}}{\scalebox{1.25}{$\cap$}}{\cap}{\cap}}\displaylimits}%
  \def\lesseqqgtr{\mathrel{\lessgtr}}\def\bigoplus{\mathop{\oplus}\displaylimits}%
}
\providecommand{\ssi}{si et seulement si} % abréviation fréquente
\AtBeginDocument{\providecommand{\wideparen}{\overparen}} % arc de cercle

%% FIGFIX-BEGIN v5 — correctifs globaux des figures (docs/onbuch-generation/tools/figfix)
\usepackage{adjustbox}\usepackage{etoolbox}\tcbuselibrary{raster}
% toute figure TikZ/pgfplots plus large que la ligne est réduite à la largeur disponible
\BeforeBeginEnvironment{tikzpicture}{\begin{adjustbox}{max width=\linewidth}}
\AfterEndEnvironment{tikzpicture}{\end{adjustbox}}
% légendes pgfplots lisibles par-dessus les courbes
\pgfplotsset{every axis/.append style={legend style={fill=white,fill opacity=0.92,text opacity=1,draw=black!15,inner sep=3pt}}}
% étiquettes d'arêtes (syntaxe edge["..."]) avec fond blanc
\tikzset{every edge quotes/.append style={fill=white,inner sep=1.2pt}}
%% FIGFIX-END
\begin{document}

\pagegarde{La période allemande}{Histoire}{1re Tech}{Histoire – classe de Première technique}{N05}{2 séances (fiche de progression harmonisée)}{Cette leçon étudie la naissance de l'entité territoriale « Cameroun » sous domination allemande (1884-1916) : la conquête, la délimitation des frontières, l'organisation administrative et économique de la colonie. Elle analyse ensuite les réactions des populations camerounaises face à cette administration étrangère, entre résistances armées et collaborations.
\vspace{4pt}
\begin{multicols}{2}\small
\begin{itemize}
\item À la fin de cette leçon, je sais situer dans le temps et dans l'espace la période allemande au Cameroun (1884-1916)
\item À la fin de cette leçon, je sais expliquer les circonstances de la signature du traité germano-douala du 12 juillet 1884
\item À la fin de cette leçon, je sais décrire les étapes de la conquête et de la délimitation territoriale du Kamerun
\item À la fin de cette leçon, je sais construire une carte ou un croquis montrant l'évolution territoriale du Cameroun allemand
\item À la fin de cette leçon, je sais présenter l'organisation administrative et l'exploitation économique du Kamerun
\item À la fin de cette leçon, je sais identifier et caractériser les formes de résistance des Camerounais face à l'administration allemande
\end{itemize}\end{multicols}}

\begin{sommaire}
\begin{multicols}{2}
\begin{enumerate}
\item Introduction : définitions et repères
\item La fondation du Kamerun (1884-1911)
\item L'organisation administrative et économique du Kamerun
\item Les Camerounais face à l'administration allemande : les résistances
\item Collaborations et fin de la période allemande
\item Travaux pratiques : construire la carte du Kamerun
\item Activité d'intégration
\item Méthodes et exercices résolus
\item Exercices
\item Corrigés des exercices
\item Fiche bilan
\end{enumerate}
\end{multicols}
\end{sommaire}

\begin{objectifs}
\begin{itemize}
\item À la fin de cette leçon, je sais situer dans le temps et dans l'espace la période allemande au Cameroun (1884-1916)
\item À la fin de cette leçon, je sais expliquer les circonstances de la signature du traité germano-douala du 12 juillet 1884
\item À la fin de cette leçon, je sais décrire les étapes de la conquête et de la délimitation territoriale du Kamerun
\item À la fin de cette leçon, je sais construire une carte ou un croquis montrant l'évolution territoriale du Cameroun allemand
\item À la fin de cette leçon, je sais présenter l'organisation administrative et l'exploitation économique du Kamerun
\item À la fin de cette leçon, je sais identifier et caractériser les formes de résistance des Camerounais face à l'administration allemande
\item À la fin de cette leçon, je sais distinguer résistance et collaboration à partir de documents historiques
\item À la fin de cette leçon, je sais rédiger et justifier une conclusion argumentée à partir de sources (cartes, photographies, données chiffrées)
\item À la fin de cette leçon, je sais appliquer les notions de colonisation et de résistance à des situations concrètes de la vie courante (débats sur la mémoire coloniale, restitution des biens culturels)
\end{itemize}
\end{objectifs}

\begin{prerequis}
\begin{itemize}
\item La traite négrière et les premiers contacts euro-africains (classe de 4e)
\item La conférence de Berlin (1884-1885) et le partage de l'Afrique (classe de 4e)
\item La localisation des grandes régions et des peuples du Cameroun actuel
\item La notion de colonisation et de protectorat
\item La lecture d'une carte historique (légende, échelle, orientation)
\end{itemize}
\end{prerequis}

\newpage

\coursec{Introduction : définitions et repères}

\courssub{Situation d'accroche et problématique}

À Douala, les élèves passent chaque jour devant le \cle{Palais des Bell} et la \cle{Pagode}, bâtiments construits sous l'administration allemande au début du XX\textsuperscript{e} siècle. Ces édifices, encore debout au cœur de la capitale économique du Cameroun, sont des témoins silencieux d'une époque où notre pays n'était pas un État indépendant, mais une possession de l'Empire allemand. Pourquoi des bâtiments allemands se dressent-ils encore au cœur de nos villes ? Comment le territoire que nous appelons aujourd'hui « Cameroun » a-t-il été dessiné par des Européens entre 1884 et 1916 ? Et comment nos ancêtres ont-ils réagi face à cette administration étrangère ?

Comprendre cette période, c'est comprendre l'origine de nos frontières actuelles et de notre mémoire collective. C'est aussi comprendre pourquoi le Cameroun, contrairement à beaucoup d'États africains, a connu successivement trois puissances coloniales : l'Allemagne, la France et le Royaume-Uni.

\begin{aretenir}[Problématique de la leçon]
Comment l'Allemagne a-t-elle constitué l'entité territoriale « Cameroun » entre 1884 et 1916, et comment les populations camerounaises ont-elles réagi face à cette administration étrangère ?
\end{aretenir}

Pour répondre à cette problématique, il faut d'abord définir précisément ce qu'est la période allemande, maîtriser le vocabulaire de la colonisation, situer le Kamerun dans l'espace et mémoriser les repères chronologiques essentiels.

\courssub{Qu'est-ce que la période allemande ?}

\textbf{Intuition.} Imagine qu'un étranger s'installe dans ta maison, en change les règles, y impose sa langue et décide de ce qui s'y produit, tout en affirmant qu'il est là pour ton bien. C'est, à l'échelle d'un territoire entier, ce qu'a vécu le Cameroun entre 1884 et 1916 avec l'Allemagne.

\begin{definition}[La période allemande]
La \cle{période allemande} désigne la présence et l'administration de l'Empire allemand sur le territoire du Cameroun, de \textbf{1884} (signature du traité germano-douala) à \textbf{1916} (départ des Allemands pendant la Première Guerre mondiale). Pendant ces 32 années, le Cameroun est une colonie allemande appelée \cle{Kamerun}.
\end{definition}

Pour bien comprendre cette définition, il faut distinguer deux notions fondamentales du vocabulaire colonial : la colonie et le protectorat.

\begin{definition}[Colonie]
Une \cle{colonie} est un territoire conquis et administré par une puissance étrangère qui y exerce sa souveraineté. La puissance coloniale y nomme ses propres administrateurs, y applique ses lois et exploite ses ressources à son profit.
\end{definition}

\begin{definition}[Protectorat]
Un \cle{protectorat} est un territoire dont le souverain local conserve une autorité nominale (en apparence), mais qui est placé sous la « protection » d'une puissance étrangère. En réalité, la puissance « protectrice » contrôle les décisions essentielles : politique extérieure, économie, armée.
\end{definition}

\begin{attention}[Ne pas confondre colonie et protectorat]
Au Probatoire et au Baccalauréat technique, une erreur fréquente consiste à employer les deux termes comme des synonymes. Retiens la différence essentielle : dans une \textbf{colonie}, la souveraineté locale disparaît totalement ; dans un \textbf{protectorat}, le chef local reste en place mais n'a plus de pouvoir réel. Le Kamerun allemand a été juridiquement constitué comme un \emph{Schutzgebiet} (territoire de protection) par le traité de 1884 : les chefs douala croyaient garder leur autorité, mais l'administration allemande a très vite imposé une souveraineté effective, transformant le protectorat de fait en administration coloniale directe.
\end{attention}

\courssub{D'où vient le nom « Kamerun » ?}

\textbf{Intuition.} Les noms de pays ne tombent pas du ciel : ils racontent souvent une histoire. Celui du Cameroun raconte... des crevettes !

Au XV\textsuperscript{e} siècle, des navigateurs portugais remontent la côte atlantique de l'Afrique. En atteignant l'estuaire du Wouri (l'actuelle baie de Douala), ils sont frappés par l'abondance des crevettes dans les eaux du fleuve. Ils baptisent alors ce lieu \emph{Rio dos Camarões}, c'est-à-dire « rivière aux crevettes » en portugais.

Ce nom a voyagé à travers les langues européennes :
\begin{itemize}
\item les \textbf{Portugais} disent \emph{Camarões} ;
\item les \textbf{Espagnols} transforment en \emph{Camarones} ;
\item les \textbf{Allemands}, en s'installant en 1884, germanisent le nom en \cle{Kamerun} ;
\item les \textbf{Anglais} diront \emph{Cameroons} et les \textbf{Français} \emph{Cameroun}.
\end{itemize}

\begin{savaistu}[Un pays nommé d'après des crevettes]
Le mot « Cameroun » vient du portugais \emph{camarões}, qui signifie « crevettes ». Ces crustacés étaient si abondants dans l'estuaire du Wouri au XV\textsuperscript{e} siècle que les navigateurs portugais en ont fait le nom du lieu. Le Cameroun est donc l'un des rares pays du monde dont le nom désigne un animal ! Cette étymologie rappelle aussi que les contacts entre l'Europe et nos côtes sont anciens : ils remontent à près de quatre siècles avant la colonisation allemande.
\end{savaistu}

\begin{aretenir}[À retenir]
\cle{Kamerun} est le nom donné par les Allemands à leur colonie. Il dérive du portugais \emph{Rio dos Camarões} (« rivière aux crevettes »), nom donné à l'estuaire du Wouri par les navigateurs portugais du XV\textsuperscript{e} siècle.
\end{aretenir}

\courssub{Les repères chronologiques essentiels (1884--1916)}

La période allemande est courte : 32 ans seulement. Mais elle est dense en événements. Trois grandes dates structurent toute la leçon et doivent être connues par cœur.

\begin{itemize}
\item \textbf{1884 : le traité germano-douala.} Le 12 juillet 1884, les chefs douala (notamment les rois Bell et Akwa) signent avec les représentants allemands un traité qui place leur territoire sous « protection » allemande. C'est l'acte de naissance officiel du Kamerun.
\item \textbf{1901 : transfert de la capitale à Buea.} Le gouverneur Jesko von Puttkamer déplace le siège de l'administration de Douala vers Buea, au pied du mont Cameroun, pour des raisons climatiques et sanitaires. Yaoundé ne deviendra capitale qu'ultérieurement (vers 1909-1910).
\item \textbf{1911 : l'accord de Rabat (Neukamerun).} À la suite de la crise d'Agadir, la France cède à l'Allemagne une vaste zone à l'est et au sud (le « Neukamerun »), portant la colonie à environ 760 000 km\textsuperscript{2}.
\item \textbf{1914--1916 : les campagnes militaires de la Première Guerre mondiale.} Quand la guerre éclate en Europe en août 1914, elle s'étend aux colonies. Les troupes françaises et britanniques attaquent le Kamerun ; les Allemands résistent, notamment autour de Yaoundé et de Mora.
\item \textbf{1916 : le départ des Allemands.} Vaincus, les Allemands abandonnent le Kamerun en février 1916. Le territoire est alors partagé entre la France et le Royaume-Uni : c'est la fin de la période allemande.
\end{itemize}

\begin{popfigure}
\begin{center}
\begin{tikzpicture}[scale=1]
% Axe principal
\draw[->,thick,popdark] (0,0) -- (13.5,0);
% Graduations et dates
\foreach \x/\annee in {0.5/1884, 2.5/1894, 4.5/1901, 6.5/1905, 8.5/1911, 10.5/1914, 12.5/1916}{
  \draw[popdark] (\x,-0.12) -- (\x,0.12);
}
% Événements au-dessus
\node[above,align=center,font=\small,text=popblue] at (0.5,0.25) {\textbf{1884}\\Traité germano-douala\\naissance du Kamerun};
\node[above,align=center,font=\small,text=popgreen] at (4.5,0.25) {\textbf{1901}\\Transfert de la\\capitale à Buea};
\node[above,align=center,font=\small,text=popgold] at (8.5,0.25) {\textbf{1911}\\Accord de Rabat :\\le Neukamerun};
% Événements en dessous
\node[below,align=center,font=\small,text=popmuted] at (2.5,0.25) {\textbf{Années 1890}\\Conquête de l'intérieur\\et premières résistances};
\node[below,align=center,font=\small,text=poporange] at (6.5,0.25) {\textbf{1909--1910}\\Yaoundé devient\\la nouvelle capitale};
\node[below,align=center,font=\small,text=poporange] at (10.5,0.25) {\textbf{1914--1916}\\Campagnes militaires\\de la 1\textsuperscript{re} Guerre mondiale};
\node[below,align=center,font=\small,text=poppink] at (12.5,0.25) {\textbf{1916}\\Départ des Allemands\\Partage franco-britannique};
\end{tikzpicture}
\end{center}
\legende{Frise chronologique de la période allemande (1884--1916). Les dates en couleur sont les repères essentiels à mémoriser : 1884 (fondation), 1901 (Buea), 1911 (Neukamerun), 1914--1916 (guerre), 1916 (fin).}
\end{popfigure}

\begin{methode}[Construire une frise chronologique]
Pour construire une frise chronologique claire et correcte, suis ces étapes :
\begin{enumerate}
\item \textbf{Choisir une échelle régulière} : décide d'une correspondance fixe entre le temps et l'espace (par exemple 1 cm = 5 ans) et respecte-la sur toute la frise. Ne place jamais les dates « au hasard ».
\item \textbf{Tracer un axe orienté} : une flèche indique le sens du temps, du passé vers le présent (généralement de gauche à droite).
\item \textbf{Placer les dates clés} : marque d'abord les événements les plus importants (début, tournants, fin) à leur position exacte sur l'échelle.
\item \textbf{Hiérarchiser les événements} : distingue visuellement les événements majeurs (en gras, au-dessus de l'axe) des événements secondaires (en dessous, en plus petit).
\item \textbf{Légender} : chaque date doit être accompagnée d'un intitulé court et précis (date + événement).
\end{enumerate}
\end{methode}

\courssub{Localiser le Kamerun : en Afrique centrale et dans l'empire colonial allemand}

\textbf{Où se trouve le Kamerun ?} Le Kamerun se situe en \cle{Afrique centrale}, au fond du golfe de Guinée, là où la côte atlantique de l'Afrique change brutalement de direction. Cette position est stratégique : le Kamerun est une « charnière » entre l'Afrique de l'Ouest et l'Afrique équatoriale, ce qui explique que les Allemands l'aient appelé « la clé de l'Afrique centrale ».

\textbf{Le Kamerun dans l'empire colonial allemand.} Arrivée tardivement dans la course coloniale (l'unification de l'Allemagne ne date que de 1871), l'Allemagne possède quatre colonies en Afrique :
\begin{itemize}
\item le \textbf{Togo} (Afrique de l'Ouest), petite colonie côtière ;
\item le \textbf{Kamerun} (Afrique centrale), la plus vaste possession allemande d'Afrique de l'Ouest et centrale ;
\item le \textbf{Sud-Ouest africain allemand} (actuelle Namibie) ;
\item l'\textbf{Afrique orientale allemande} (actuels Tanzanie continentale, Rwanda et Burundi).
\end{itemize}

\begin{popfigure}
\begin{center}
\begin{tikzpicture}[scale=0.9]
% Contour très simplifié de l'Afrique
\draw[thick,popdark,fill=popcream] 
  (0,6) -- (2.5,7) -- (5,7.2) -- (6.5,6.5) -- (7.5,5.5) -- (8,4) -- (7.2,2.5) 
  -- (6,1) -- (4.5,0) -- (3.5,0.3) -- (2.5,1.5) -- (1.5,2.8) -- (0.5,3.5) 
  -- (0,4.5) -- cycle;
% Togo (petite bande côte ouest)
\fill[popgreen!70] (1.1,3.9) -- (1.35,3.9) -- (1.35,4.6) -- (1.1,4.6) -- cycle;
% Kamerun
\fill[popblue!60] (1.5,2.9) -- (2.6,2.9) -- (2.9,3.8) -- (2.2,4.5) -- (1.6,4.2) -- cycle;
% Sud-Ouest africain
\fill[popgold!70] (2.6,0.5) -- (4,0.4) -- (4.2,1.6) -- (3,1.8) -- (2.6,1.2) -- cycle;
% Afrique orientale allemande
\fill[poporange!60] (5.6,2.4) -- (6.9,2.7) -- (6.7,4.2) -- (5.7,4.3) -- (5.4,3.3) -- cycle;
% Étiquettes
\node[font=\scriptsize,align=center] at (2.1,3.7) {\textbf{Kamerun}};
\node[font=\scriptsize,align=center,text=popdark] at (0.4,5.1) {\textbf{Togo}};
\draw[->,popdark] (0.6,4.95) -- (1.15,4.35);
\node[font=\scriptsize,align=center] at (3.4,1.1) {\textbf{Sud-Ouest}\\\textbf{africain}};
\node[font=\scriptsize,align=center] at (6.15,3.35) {\textbf{Afrique}\\\textbf{orientale}};
% Flèche du nord
\draw[->,thick,popdark] (7.8,6.2) -- (7.8,7);
\node[font=\small,popdark] at (7.8,7.2) {\textbf{N}};
% Océan
\node[font=\scriptsize\itshape,text=popmuted] at (0.6,2) {Océan Atlantique};
\end{tikzpicture}
\end{center}
\legende{Localisation schématique du Kamerun (en bleu) dans l'empire colonial allemand en Afrique : Togo (vert), Sud-Ouest africain (jaune, actuelle Namibie), Afrique orientale allemande (orange). Croquis simplifié, sans précision cartographique.}
\end{popfigure}

\begin{popfigure}
\begin{center}
\begin{tikzpicture}[scale=1.2]
% Contour Kamerun 1884 (côte + arrière-pays immédiat)
\draw[thick,popblue,fill=popblue!20] 
  (0,0) -- (3,0) -- (3.5,1) -- (3.2,2.5) -- (2,3.5) -- (0.5,3) -- (-0.2,2) -- (-0.2,0.5) -- cycle;
% Extension Neukamerun 1911 (Est et Sud)
\draw[thick,poporange,fill=poporange!20,opacity=0.6] 
  (3.2,2.5) -- (5,2.5) -- (5,0.5) -- (4.5,0) -- (3,0) -- (3.5,1) -- cycle;
\draw[thick,poporange,fill=poporange!20,opacity=0.6] 
  (0.5,3) -- (2,3.5) -- (2.5,5) -- (-0.5,5) -- (-0.2,2) -- cycle;
% Villes
\node[circle,fill=popdark,inner sep=1.5pt,label=below:\tiny Douala] at (0.5,0.8) {};
\node[circle,fill=popdark,inner sep=1.5pt,label=above:\tiny Buea] at (1.5,1.5) {};
\node[circle,fill=popdark,inner sep=1.5pt,label=right:\tiny Yaoundé] at (2.5,2.5) {};
\node[circle,fill=popdark,inner sep=1.5pt,label=above:\tiny Garoua] at (3.5,3.5) {};
\node[circle,fill=popdark,inner sep=1.5pt,label=above:\tiny Mora] at (4.5,1.5) {};
% Légende interne
\node[font=\tiny,align=left,text=popblue] at (-1.2,3.5) {\textbf{Kamerun 1884}\\(zone côtière initiale)};
\node[font=\tiny,align=left,text=poporange] at (5.2,1.5) {\textbf{Neukamerun 1911}\\(extension Est/Sud)};
% Fleche Nord
\draw[->,thick,popdark] (5.5,4.5) -- (5.5,5);
\node[font=\tiny,popdark] at (5.5,5.2) {N};
\end{tikzpicture}
\end{center}
\legende{Évolution des frontières du Kamerun : en bleu, le territoire initial de 1884 (zone côtière et arrière-pays proche) ; en orange, l'extension du \emph{Neukamerun} (1911) vers l'est (bassin de la Sangha) et le sud (bassin du Congo). Les points noirs indiquent les principales stations administratives. Croquis schématique.}
\end{popfigure}

\begin{attention}[Piège fréquent à l'examen]
Ne confonds pas le \textbf{Kamerun allemand} avec le \textbf{Cameroun actuel} ! Les limites ne sont pas identiques : le Kamerun de 1911 (avec le Neukamerun) était plus vaste à l'est et au sud que le Cameroun actuel. Après 1916, le territoire est partagé entre la France et le Royaume-Uni, et le Cameroun indépendant de 1960 ne correspond pas exactement au Kamerun de 1911. Par ailleurs, ne dis jamais que l'Allemagne a « découvert » le Cameroun : le territoire était habité et organisé par des sociétés (chefferies, royaumes, lamidats) bien avant 1884.
\end{attention}

\courssub{Application : vérifier ses repères}

\begin{exoresolu}[Replacer les événements dans le temps]
\textbf{Énoncé.} Un élève affirme : « Le Kamerun a été fondé en 1914 et les Allemands sont partis en 1884, après la signature du traité germano-douala. » Relever les erreurs de cette phrase et rédiger une version correcte et justifiée.
\tcblower
\textbf{Solution détaillée.} La phrase contient deux inversions chronologiques majeures.
\begin{enumerate}
\item Le Kamerun n'a pas été fondé en 1914 : il est né en \textbf{1884}, avec la signature du \textbf{traité germano-douala} (12 juillet 1884), par lequel les chefs douala placent leur territoire sous protection allemande.
\item Les Allemands ne sont pas partis en 1884 : ils sont partis en \textbf{1916}, à la fin des campagnes militaires de la Première Guerre mondiale (1914--1916), après leur défaite face aux troupes franco-britanniques.
\end{enumerate}
\textbf{Version correcte :} « Le Kamerun a été fondé en 1884 par le traité germano-douala ; les Allemands l'ont quitté en 1916, à la fin de la Première Guerre mondiale. » La période allemande a donc duré 32 ans ($1916 - 1884 = 32$).
\end{exoresolu}

\begin{exercice}[Pour s'entraîner]
\begin{enumerate}
\item Définis en deux phrases la « période allemande » au Cameroun, en citant les deux dates qui la bornent.
\item Explique l'origine du mot « Kamerun » en trois lignes maximum.
\item Construis, sur ta feuille, une frise chronologique de 1884 à 1916 à l'échelle 1 cm = 4 ans, en y plaçant les cinq dates clés de la leçon (1884, 1901, 1911, 1914, 1916).
\end{enumerate}
\end{exercice}

\begin{corrige}[Exercice : Pour s'entraîner]
\begin{enumerate}
\item La période allemande est la période pendant laquelle l'Empire allemand a administré le territoire du Cameroun, sous le nom de Kamerun. Elle s'étend de 1884 (traité germano-douala) à 1916 (départ des Allemands pendant la Première Guerre mondiale).
\item « Kamerun » est le nom allemand de la colonie. Il dérive du portugais \emph{Rio dos Camarões} (« rivière aux crevettes »), nom donné à l'estuaire du Wouri par les navigateurs portugais du XV\textsuperscript{e} siècle, frappés par l'abondance des crevettes.
\item À l'échelle 1 cm = 4 ans, les 32 années de la période occupent 8 cm. On place : 1884 à 0 cm (traité germano-douala), 1901 à 4,25 cm (transfert capitale à Buea), 1911 à 6,75 cm (Neukamerun), 1914 à 7,5 cm (début guerre), 1916 à 8 cm (départ Allemands). L'axe est orienté de gauche à droite, les événements majeurs sont notés au-dessus.
\end{enumerate}
\end{corrige}

\begin{aretenir}[L'essentiel de l'introduction]
\begin{itemize}
\item La \cle{période allemande} : présence et administration de l'Allemagne au Cameroun, de \textbf{1884 à 1916} (32 ans).
\item \cle{Kamerun} : nom allemand de la colonie, dérivé du portugais \emph{camarões} (crevettes).
\item Cinq dates clés : \textbf{1884} (traité germano-douala), \textbf{1901} (capitale à Buea), \textbf{1911} (Neukamerun), \textbf{1914--1916} (campagnes de la Première Guerre mondiale), \textbf{1916} (départ des Allemands et partage).
\item Le Kamerun se situe en Afrique centrale ; l'empire colonial allemand comprend aussi le Togo, le Sud-Ouest africain et l'Afrique orientale allemande.
\item Problématique : comment l'Allemagne a-t-elle constitué l'entité territoriale « Cameroun » et comment les populations ont-elles réagi ?
\end{itemize}
\end{aretenir}

\coursec{La fondation du Kamerun (1884-1911)}

Dans la section précédente, nous avons posé les définitions essentielles (colonie, protectorat, traité) et situé les grands repères chronologiques de la période allemande. Nous savons donc \emph{ce qu'est} une colonie. Reste à comprendre \emph{comment} l'entité territoriale appelée « Kamerun » est née, entre 1884 et 1911. C'est l'objet de cette section : partir d'un simple comptoir commercial sur la côte, et suivre, étape par étape, la construction d'une colonie qui atteindra près de \num{790000} km² en 1911.

\courssub{Le contexte : la rivalité des puissances européennes dans le golfe de Guinée}

\textbf{Une course entre trois puissances.} Dans la seconde moitié du XIX\textsuperscript{e} siècle, trois puissances européennes se disputent les côtes du golfe de Guinée : la \cle{Grande-Bretagne}, installée au Nigeria voisin ; la \cle{France}, qui avance depuis le Gabon et le Congo ; et l'\cle{Allemagne}, unie en 1871, qui cherche à son tour des débouchés commerciaux et des matières premières (huile de palme, ivoire, caoutchouc). L'estuaire du Wouri, où se trouvent les royaumes côtiers Douala, devient un point de convoitise : c'est une porte d'entrée vers l'intérieur du continent.

\textbf{La conférence de Berlin (novembre 1884 -- février 1885).} Réunie à l'initiative du chancelier allemand \cle{Bismarck}, cette conférence rassemble les puissances européennes pour fixer les règles du partage de l'Afrique. Elle pose deux principes essentiels :
\begin{itemize}
\item la \cle{liberté de navigation} et de commerce sur les grands fleuves (Congo, Niger) ;
\item le principe d'\cle{occupation effective} : pour revendiquer un territoire, une puissance doit l'occuper réellement (traités avec les chefs locaux, postes administratifs, troupes) et en informer les autres puissances.
\end{itemize}
Aucun Africain n'est invité à Berlin. La conférence ne partage pas l'Afrique à proprement parler : elle fixe les \emph{règles du jeu} de la conquête. Pour l'Allemagne, cela signifie qu'il faut agir vite, avant les Britanniques et les Français.

\begin{attention}[Ne pas confondre deux dates]
Erreur fréquente au Probatoire : confondre la \textbf{conférence de Berlin} (novembre 1884 -- février 1885), qui est une réunion \emph{européenne} fixant les règles du partage de l'Afrique, avec le \textbf{traité germano-douala} du \textbf{12 juillet 1884}, signé \emph{au Cameroun}, qui fonde le Kamerun. Le traité est \emph{antérieur} de quelques mois à la conférence : l'Allemagne a pris les devants, puis a fait reconnaître son fait accompli à Berlin.
\end{attention}

\courssub{Les commerçants allemands sur la côte : Woermann et Jantzen \& Thormählen}

Avant l'arrivée de l'État allemand, ce sont des \textbf{commerçants privés} qui préparent le terrain. Deux grandes maisons de commerce hambourgeoises sont installées à Douala :
\begin{itemize}
\item la maison \cle{Woermann}, la plus importante, spécialisée dans le commerce de l'huile de palme, des noix de palme et de l'ivoire ;
\item la maison \cle{Jantzen \& Thormählen}, dirigée localement par Johannes Thormählen.
\end{itemize}
Ces commerçants entretiennent des relations suivies avec les chefs Douala, qui leur vendent les produits de l'intérieur. Mais ils subissent la concurrence britannique et française, et craignent que les chefs locaux ne passent des accords avec leurs rivaux. Ils font donc pression sur le gouvernement de Berlin pour que l'Allemagne prenne officiellement le contrôle de la côte. Le chancelier Bismarck, d'abord réticent à l'aventure coloniale, finit par céder : il dépêche sur place un explorateur et diplomate, \cle{Gustav Nachtigal}, nommé consul général pour l'Afrique de l'Ouest.

\begin{savaistu}[Les rois Douala avaient d'abord choisi... les Britanniques]
Avant 1884, les rois Bell et Akwa, inquiets des rivalités commerciales et soucieux de protéger leur position, avaient adressé plusieurs demandes de \emph{protectorat} à la Grande-Bretagne (notamment en 1879 et 1881). Mais Londres tardait à répondre. Quand Nachtigal arrive devant Douala à bord de la canonnière \emph{Möwe}, en juillet 1884, la réponse britannique n'est toujours pas arrivée. Les Allemands signent le traité le 12 juillet ; le navire britannique n'apparaît que quelques jours plus tard, trop tard. C'est un cas célèbre où la lenteur d'une administration a changé le destin d'un pays : le Cameroun aurait pu devenir une colonie anglaise.
\end{savaistu}

\courssub{Le traité germano-douala du 12 juillet 1884}

\begin{definition}[Traité]
Un \cle{traité} est un accord écrit signé entre deux parties (États, ou ici un État européen et des chefs locaux), qui fixe des engagements réciproques : cessions de droits, promesses, obligations. En droit, il engage les signataires ; dans les faits, au XIX\textsuperscript{e} siècle, les traités coloniaux étaient souvent déséquilibrés et imposés.
\end{definition}

\textbf{Les signataires.} Le 12 juillet 1884, à Douala, Gustav \cle{Nachtigal}, au nom de l'empereur d'Allemagne, signe un traité avec les deux principaux rois Douala : le roi \cle{Bell} (Ndumb'a Lobe) et le roi \cle{Akwa} (Dika Mpondo). Le roi Deïdo et d'autres chefs s'y associent ensuite. Les commerçants allemands (Woermann, Jantzen \& Thormählen) jouent les intermédiaires et les témoins.

\begin{popfigure}
\begin{center}
\begin{tikzpicture}[scale=1, every node/.style={font=\small}]
% Partie allemande
\node[draw=popblue, fill=popblueL, rounded corners, thick, minimum width=4.2cm, minimum height=1.6cm, align=center] (allemagne) at (0,3) {\textbf{Empire allemand}\\ Gustav Nachtigal\\ (consul général)};
\node[draw=popgold, fill=popgoldL, rounded corners, thick, minimum width=4.2cm, minimum height=1.2cm, align=center] (commercants) at (0,0.6) {\textbf{Commerçants allemands}\\ Woermann -- Jantzen \& Thormählen};
% Partie douala
\node[draw=popgreen, fill=popgreenL, rounded corners, thick, minimum width=3.4cm, minimum height=1.2cm, align=center] (bell) at (7.2,3.9) {\textbf{Roi Bell}\\ (Ndumb'a Lobe)};
\node[draw=popgreen, fill=popgreenL, rounded corners, thick, minimum width=3.4cm, minimum height=1.2cm, align=center] (akwa) at (7.2,2.1) {\textbf{Roi Akwa}\\ (Dika Mpondo)};
% Traité au centre
\node[draw=popink, fill=popcream, rounded corners, thick, minimum width=3.6cm, minimum height=1.4cm, align=center] (traite) at (3.8,3) {\textbf{Traité germano-douala}\\ 12 juillet 1884};
\draw[-{Stealth}, very thick, popblue] (allemagne.east) -- (traite.west);
\draw[-{Stealth}, very thick, popgreen] (bell.west) -- (traite.east |- bell.west);
\draw[-{Stealth}, very thick, popgreen] (akwa.west) -- (traite.east |- akwa.west);
\draw[-{Stealth}, thick, popgold, dashed] (commercants.north) -- (allemagne.south) node[midway, right, font=\footnotesize, align=left]{pression sur\\ Berlin};
\draw[-{Stealth}, thick, popgold, dashed] (commercants.east) to[bend left=10] node[midway, below, font=\footnotesize]{intermédiaires} (traite.south);
\end{tikzpicture}
\end{center}
\legende{Schéma des acteurs du traité du 12 juillet 1884. Flèches pleines : signataires. Flèches pointillées : rôle d'intermédiaire et de pression des commerçants allemands.}
\end{popfigure}

\textbf{Les clauses du traité.} Le texte comporte des engagements des deux côtés, mais ils sont profondément inégaux :
\begin{itemize}
\item les rois Bell et Akwa \textbf{cèdent la souveraineté} de leurs territoires à l'Allemagne, ainsi que le droit exclusif de commercer et de légiférer ;
\item en échange, l'Allemagne promet la \textbf{protection} des rois et de leurs populations, le \textbf{respect des coutumes et des traditions} locales, et le maintien des droits des chefs sur leurs terres ;
\item le drapeau allemand est hissé à Douala le 14 juillet 1884 : le « \cle{Kamerun} » est proclamé protectorat allemand.
\end{itemize}

\begin{attention}[Des promesses non tenues]
Les clauses de protection et de respect des coutumes ne seront \textbf{pas respectées} : expropriations de terres, travail forcé, impôts, répression brutale des résistances marqueront l'administration allemande (voir section 4). Dans une copie, il faut donc distinguer le \emph{texte} du traité de sa \emph{mise en œuvre} réelle.
\end{attention}

\begin{methode}[Analyser un traité historique]
Pour étudier un traité (exercice fréquent de commentaire de document), suivre cinq étapes :
\begin{enumerate}
\item \textbf{Identifier les signataires} : qui signe, au nom de qui ?
\item \textbf{Dater et situer} : quand, où, dans quel contexte ?
\item \textbf{Dégager les clauses} : que cède chaque partie ? que promet-elle ?
\item \textbf{Comparer les intérêts} : que gagne l'Allemagne (territoire, commerce) ? que cherchent les rois Douala (protection, avantage commercial face à leurs rivaux) ?
\item \textbf{Confronter le texte aux faits} : les promesses ont-elles été tenues ? quelles conséquences ?
\end{enumerate}
\end{methode}

\begin{exoresolu}[Analyser la clause de cession de souveraineté]
Énoncé : « Les souverains soussignés cèdent [...] les droits de souveraineté, de législation et d'administration [...] à l'Empire allemand » (traité du 12 juillet 1884). Question : que signifie cette clause et pourquoi les rois Douala l'ont-ils acceptée ?
\tcblower
\textbf{Solution.} \emph{Signification} : les rois abandonnent à l'Allemagne le pouvoir de gouverner, de faire les lois et de contrôler le commerce sur leurs territoires ; ils ne restent que des chefs locaux subordonnés. \emph{Raisons de l'acceptation} : les rois Douala, commerçants intermédiaires, pensaient obtenir la protection d'une puissance européenne contre leurs rivaux (autres chefs côtiers, commerçants britanniques) tout en conservant leurs privilèges commerciaux, comme le traité le promettait. Ils ont sans doute sous-estimé la portée de la cession : la notion européenne de « souveraineté » transférée à jamais n'avait pas d'équivalent dans leur conception du pouvoir. \emph{Bilan} : clause déséquilibrée, dont les promesses compensatoires (protection, respect des coutumes) ne seront pas tenues.
\end{exoresolu}

\courssub{L'extension vers l'intérieur : traités, expéditions et conquête militaire}

Le traité de 1884 ne concerne que la \textbf{côte}. Or la conférence de Berlin exige l'\emph{occupation effective} : pour revendiquer l'intérieur, l'Allemagne doit y pénétrer réellement. Trois moyens sont combinés :
\begin{enumerate}
\item \textbf{les expéditions d'exploration} : des missions (souvent mixtes, scientifiques et militaires) remontent les fleuves (Wouri, Sanaga, Nyong) et cartographient le pays ;
\item \textbf{les traités avec les chefs de l'intérieur} : sur le modèle du traité germano-douala, les Allemands font signer des dizaines de traités de « protection » aux chefs des régions intérieures ;
\item \textbf{la conquête militaire} : quand les chefs refusent de signer ou se révoltent, l'armée coloniale (la \emph{Schutztruppe}) intervient. C'est dans ce cadre qu'un poste militaire est fondé en \textbf{1889} à \cle{Yaoundé}, sur le territoire des Ewondo, par le lieutenant Kund : ce poste deviendra la capitale du Kamerun.
\end{enumerate}
L'avancée est lente et violente : forêts denses, maladies, résistances armées des populations. Vers 1901, l'essentiel du territoire compris entre la côte, le Sanaga et le nord est sous contrôle allemand, au prix de nombreuses campagnes de « pacification ».

\courssub{La délimitation des frontières avec les voisins}

Une colonie n'existe pleinement que lorsque ses \cle{frontières} sont fixées par des accords internationaux. L'Allemagne négocie avec ses deux voisins :
\begin{itemize}
\item \textbf{Avec la Grande-Bretagne} (qui occupe le Nigeria) : accords de \textbf{1885} puis de \textbf{1893}, qui fixent la frontière occidentale du Kamerun, de la côte jusqu'au lac Tchad ;
\item \textbf{Avec la France} (qui occupe le Congo et le Gabon) : accords de \textbf{1885}, \textbf{1894} puis \textbf{1908}, qui tracent les frontières sud et est, le long des fleuves (Sanaga, puis Campo, Ntem, Sangha).
\end{itemize}
Ces frontières, tracées à Paris, Londres et Berlin sans consulter les populations, séparent parfois des groupes humains homogènes : c'est l'une des conséquences durables de la colonisation.

\courssub{L'agrandissement de 1911 : le « Neukamerun »}

En 1911, la crise d'Agadir (au Maroc) oppose la France et l'Allemagne. Pour régler le conflit, les deux pays signent le \textbf{traité de Fès} (4 novembre 1911) : l'Allemagne reconnaît le protectorat français sur le Maroc et reçoit en compensation d'importants territoires du \textbf{Congo français}, rattachés au Kamerun sous le nom de \cle{Neukamerun} (« Nouveau-Cameroun »). Ces territoires correspondent en grande partie aux actuels Tchad méridional, Centrafrique occidentale et Congo septentrional, et donnent au Kamerun un accès au fleuve Congo.

\begin{exemplebox}[Un exemple chiffré : le Kamerun à son extension maximale]
Avant 1911, le Kamerun couvre environ \num{495000} km². Le Neukamerun ajoute environ \num{295000} km². Le Kamerun atteint alors près de \num{790000} km², contre \num{475000} km² pour le Cameroun actuel. Vérification du calcul : $495 + 295 = 790$ (en milliers de km²). Le Neukamerun représente donc une part de $\dfrac{295}{790} \approx \num{0,37}$, soit environ 37 \% du territoire total de 1911.
\end{exemplebox}

\begin{popfigure}
\begin{center}
\begin{tikzpicture}[scale=0.9, every node/.style={font=\small}]
% Zone 1884 : la côte
\fill[popblue!45] (0,0) -- (2.2,0) -- (2.2,1.2) -- (1.4,1.6) -- (0.4,1.4) -- cycle;
% Zone 1901 : intérieur
\fill[popgreen!35] (0.4,1.4) -- (1.4,1.6) -- (2.2,1.2) -- (2.2,0) -- (4.6,0.2) -- (5.2,1.6) -- (4.4,3.4) -- (2.6,4.2) -- (1.0,3.4) -- cycle;
% Zone 1911 : Neukamerun
\fill[poporange!40] (2.6,4.2) -- (4.4,3.4) -- (5.2,1.6) -- (7.6,1.8) -- (8.6,3.2) -- (7.4,5.2) -- (4.6,5.6) -- cycle;
% Villes
\fill[popink] (1.1,0.7) circle (0.09) node[below, font=\footnotesize]{Douala (1884)};
\fill[popink] (3.1,2.3) circle (0.09) node[right, font=\footnotesize]{Yaoundé (1889)};
% Nord
\draw[-{Stealth}, thick, popink] (0.3,5.6) -- (0.3,6.3) node[above, font=\footnotesize]{N};
% Légende intégrée
\node[draw=popblue, fill=popblue!45, minimum width=0.5cm, minimum height=0.35cm, anchor=west] at (9.5,5.0) {};
\node[right, align=left, font=\footnotesize] at (10.1,5.0) {\textbf{1884} : la côte (traité germano-douala)};
\node[draw=popgreen, fill=popgreen!35, minimum width=0.5cm, minimum height=0.35cm, anchor=west] at (9.5,4.4) {};
\node[right, align=left, font=\footnotesize] at (10.1,4.4) {\textbf{vers 1901} : extension vers l'intérieur\\ (traités et conquête militaire)};
\node[draw=poporange, fill=poporange!40, minimum width=0.5cm, minimum height=0.35cm, anchor=west] at (9.5,3.5) {};
\node[right, align=left, font=\footnotesize] at (10.1,3.5) {\textbf{1911} : le Neukamerun\\ (traité de Fès)};
\node[align=left, font=\footnotesize, text width=4cm] at (11.5,2.0) {Extension maximale en 1911 :\\ environ \num{790000} km²\\ (Cameroun actuel : \num{475000} km²)};
\end{tikzpicture}
\end{center}
\legende{Croquis schématique de l'évolution territoriale du Kamerun : 1. côte (1884) ; 2. intérieur conquis (vers 1901) ; 3. Neukamerun (1911). Représentation simplifiée, sans valeur cartographique exacte.}
\end{popfigure}

\begin{aretenir}[L'essentiel de la section]
\begin{itemize}
\item La conférence de \textbf{Berlin (1884-1885)} fixe les règles du partage de l'Afrique, dont l'occupation effective.
\item Le \textbf{12 juillet 1884}, Nachtigal signe avec les rois \textbf{Bell} et \textbf{Akwa} le traité germano-douala : cession de la souveraineté contre des promesses de protection non tenues. Le Kamerun est né.
\item L'extension vers l'intérieur se fait par expéditions, traités et conquête militaire ; \textbf{Yaoundé} est fondé en \textbf{1889}.
\item Les frontières sont fixées avec la Grande-Bretagne (1885, 1893) et la France (1885, 1894, 1908).
\item En \textbf{1911}, le traité de Fès rattache le \textbf{Neukamerun} : le Kamerun atteint environ \num{790000} km², son extension maximale.
\end{itemize}
\end{aretenir}

\begin{exercice}[Construire une frise et argumenter]
1. Place sur une frise chronologique : conférence de Berlin, traité germano-douala, fondation de Yaoundé, accord frontalier franco-allemand de 1894, traité de Fès.
2. En dix lignes maximum, explique pourquoi on peut dire que le Kamerun est né « d'un traité et d'un retard ».
\end{exercice}

\begin{corrige}[Exercice : frise et argumentation]
1. Ordre chronologique : traité germano-douala (12 juillet 1884) ; conférence de Berlin (novembre 1884 -- février 1885) ; fondation de Yaoundé (1889) ; accord franco-allemand (1894) ; traité de Fès (1911).
2. Éléments de réponse : le Kamerun naît d'un \emph{traité}, celui du 12 juillet 1884, par lequel les rois Bell et Akwa cèdent la souveraineté à l'Allemagne. Mais il naît aussi d'un \emph{retard} : les rois Douala avaient d'abord demandé un protectorat à la Grande-Bretagne, qui n'a pas répondu à temps. L'arrivée rapide de Nachtigal, poussée par les commerçants de Hambourg, a devancé les Britanniques de quelques jours. La fondation du Kamerun résulte donc à la fois d'une stratégie allemande délibérée et du hasard du calendrier diplomatique britannique.
\end{corrige}

\coursec{L'organisation administrative et économique du Kamerun}

Après avoir étudié comment l'Allemagne a fondé le Kamerun entre 1884 et 1911, en signant des traités avec les chefs Duala puis en repoussant les frontières vers l'intérieur, il faut maintenant répondre à une question simple : une fois le territoire conquis, comment les Allemands l'ont-ils gouverné et exploité ? C'est l'objet de cette section : comprendre la \cle{machine administrative} mise en place par Berlin et le système économique qui a fait du Kamerun une colonie de rente.

\courssub{Le statut du Kamerun : du protectorat à la colonie impériale}

Lorsque le consul général allemand Gustav Nachtigal hisse le drapeau allemand à Douala le 14 juillet 1884, le Kamerun n'est pas encore une colonie au sens strict. Il s'agit d'abord d'un \cle{protectorat} : l'Allemagne « protège » officiellement les chefs locaux qui ont signé des traités, mais en réalité, ce sont surtout les \cle{compagnies commerciales} (comme les firmes hambourgeoises Woermann et Jantzen und Thormählen) qui administrent le territoire à leur profit. Pour l'État allemand, cette formule est avantageuse : elle coûte peu cher, car ce sont des entreprises privées qui paient les frais d'administration.

Très vite, ce système montre ses limites : les compagnies ne pensent qu'au profit, négligent l'administration et provoquent des conflits avec les populations. L'Allemagne transforme alors le Kamerun en \cle{colonie impériale directe} (en allemand \emph{Reichskolonie}), placée sous l'autorité directe de l'empereur Guillaume II et de l'administration coloniale de Berlin. Un fonctionnaire, le \cle{gouverneur impérial}, devient le représentant tout-puissant de l'Empire sur place.

\begin{aretenir}[Deux statuts successifs]
Le Kamerun connaît deux statuts : d'abord un \textbf{protectorat} administré par les compagnies commerciales (à partir de 1884), puis une \textbf{colonie impériale directe} (\emph{Reichskolonie}) administrée par l'État allemand à travers un gouverneur impérial.
\end{aretenir}

\courssub{L'appareil administratif : gouverneur, districts et chefs auxiliaires}

Comment un petit nombre d'Allemands peut-il contrôler un territoire qui atteint, après l'agrandissement de 1911, environ 760 000 km$^2$ (près de 500 000 km$^2$ avant 1911) ? La réponse tient en une pyramide administrative à trois étages.

Au sommet se trouve le \cle{gouverneur impérial}, nommé par Berlin. Il concentre tous les pouvoirs : il commande la troupe coloniale (la \emph{Schutztruppe}), lève les impôts, rend la justice et dirige les travaux publics. Sa résidence est installée à \cle{Buéa}, sur les flancs du mont Cameroun, où le climat est plus frais qu'à Douala.

Le territoire est découpé en \cle{districts} (en allemand \emph{Bezirke}), dirigés par des administrateurs allemands (les \emph{Bezirksamtmänner}). Chaque district contrôle une portion du territoire : Douala, Buéa, Yaoundé, Ebolowa, Garoua, etc.

À la base de la pyramide, l'Allemagne s'appuie sur les \cle{chefs auxiliaires} : des chefs traditionnels recrutés ou maintenus par l'administration pour collecter l'impôt, recruter les porteurs et les travailleurs, et transmettre les ordres. C'est une administration \emph{indirecte} au quotidien : le chef local devient le rouage essentiel, mais aussi le bouc émissaire, du système colonial.

\begin{popfigure}
\begin{center}
\begin{tikzpicture}[
  box/.style={draw=popblue, fill=popblueL, rounded corners=2pt, align=center, minimum width=4.6cm, minimum height=1cm, font=\small},
  box2/.style={draw=poporange, fill=poporangeL, rounded corners=2pt, align=center, minimum width=3.4cm, minimum height=0.9cm, font=\small},
  box3/.style={draw=popgreen, fill=popgreenL, rounded corners=2pt, align=center, minimum width=3.4cm, minimum height=0.9cm, font=\small},
  link/.style={-{Stealth[length=2.5mm]}, thick, popdark}
]
\node[box, align=center] (emp) at (0,4.2){\textbf{Empereur et administration coloniale}\\ (Berlin)};
\node[box, align=center] (gouv) at (0,2.6){\textbf{Gouverneur impérial}\\ résidence à Buéa};
\node[box2, align=center] (d1) at (-4.4,0.9){District de\\ Douala};
\node[box2, align=center] (d2) at (0,0.9){District de\\ Yaoundé};
\node[box2, align=center] (d3) at (4.4,0.9){District de\\ Garoua};
\node[box3, align=center] (c1) at (-4.4,-0.9){Chefs auxiliaires\\ (impôt, corvées)};
\node[box3, align=center] (c2) at (0,-0.9){Chefs auxiliaires\\ (impôt, corvées)};
\node[box3, align=center] (c3) at (4.4,-0.9){Chefs auxiliaires\\ (impôt, corvées)};
\draw[link] (emp) -- (gouv);
\draw[link] (gouv.south) -- (d1.north);
\draw[link] (gouv.south) -- (d2.north);
\draw[link] (gouv.south) -- (d3.north);
\draw[link] (d1) -- (c1);
\draw[link] (d2) -- (c2);
\draw[link] (d3) -- (c3);
\end{tikzpicture}
\end{center}
\legende{Organigramme de l'administration coloniale allemande au Kamerun : une pyramide à trois étages (gouverneur, districts, chefs auxiliaires) sous l'autorité de Berlin.}
\end{popfigure}

\courssub{Les capitales successives et les principaux gouverneurs}

Le Kamerun allemand a connu \textbf{trois capitales successives}, et chaque déménagement raconte une étape de la colonisation :

\begin{enumerate}
\item \cle{Douala} (1884-1901) : premier point d'ancrage, c'est le grand port commercial par où passent toutes les marchandises ;
\item \cle{Buéa} (1901-1909) : capitale administrative choisie pour son climat sain, en altitude, au pied du mont Cameroun ;
\item \cle{Yaoundé} (1909-1916) : position centrale dans le territoire, plus pratique pour contrôler l'intérieur des terres après les agrandissements de 1911.
\end{enumerate}

Parmi les gouverneurs qui se succèdent, quatre noms sont à retenir :

\begin{center}
\begin{tabularx}{\linewidth}{|l|X|}
\hline
\rowcolor{popblueL} \textbf{Gouverneur} & \textbf{Action principale} \\
\hline
Julius von Soden & Premier gouverneur (à partir de 1885) ; organise l'administration naissante et fonde des postes dans l'intérieur. \\
\hline
Eugen von Zimmerer & Développe la colonisation vers le sud et l'intérieur ; encourage les plantations. \\
\hline
Jesko von Puttkamer & Gouverneur longtemps en poste ; réprime durement les résistances et favorise l'exploitation forestière. \\
\hline
Theodor Seitz & Dernier gouverneur (1910-1914) ; c'est lui qui capitule face aux Alliés en 1916. \\
\hline
\end{tabularx}
\end{center}

\courssub{L'exploitation économique : concessions et cultures de rente}

L'objectif économique de l'Allemagne est clair : faire du Kamerun une colonie \emph{rentable}, c'est-à-dire qui rapporte plus qu'elle ne coûte. Deux outils sont utilisés : les concessions aux compagnies et les cultures de rente.

\begin{definition}[Concession]
Une \cle{concession} est un vaste territoire accordé par l'État colonial à une compagnie privée pour l'exploiter. La compagnie y prélève les ressources (caoutchouc, ivoire, bois) et y impose souvent sa loi aux populations.
\end{definition}

L'État allemand accorde ainsi d'immenses concessions, notamment à la \cle{Südkamerun Gesellschaft} (Compagnie du Sud-Cameroun, créée en 1898) et à la \cle{Gesellschaft Nordwest-Kamerun} (Compagnie du Nord-Ouest, créée en 1899). Ces compagnies exploitent des millions d'hectares, souvent par des méthodes brutales qui provoquent de violentes révoltes dans le Sud forestier.

Parallèlement, se développent les \cle{cultures de rente} : des cultures destinées non pas à nourrir les populations, mais à être exportées vers l'Europe. Les principales sont :

\begin{itemize}
\item le \cle{caoutchouc}, récolté dans la forêt, très rentable mais source de terribles exactions ;
\item le \cle{cacao}, cultivé surtout autour de Yaoundé et dans le Centre ;
\item l'\cle{huile de palme}, produite dans les zones côtières ;
\item les \cle{bananes}, cultivées dans les grandes plantations de la côte, autour de Victoria (l'actuelle Limbé), où la \emph{Westafrikanische Pflanzungsgesellschaft Victoria} crée de vastes domaines.
\end{itemize}

\begin{exemplebox}[Un succès d'exportation : le cacao]
À la veille de la Première Guerre mondiale, le Kamerun exporte environ \textbf{12 000 tonnes de cacao par an}, ce qui en fait l'un des tout premiers producteurs mondiaux de l'époque. Cette réussite repose sur le travail des paysans camerounais, encadrés et pressurisés par l'administration et les missions.
\end{exemplebox}

\begin{popfigure}
\begin{center}
\begin{tikzpicture}
\begin{axis}[
  ybar, bar width=16pt,
  width=0.85\linewidth, height=6cm,
  symbolic x coords={Cacao, Huile de palme, Caoutchouc, Bananes},
  xtick=data,
  ylabel={Exportations (milliers de tonnes, vers 1913)},
  ymin=0, ymax=14,
  ytick={0,2,4,6,8,10,12,14},
  nodes near coords,
  every node near coord/.append style={font=\small},
  axis line style={popdark},
  tick label style={font=\small},
  label style={font=\small}
]
\addplot[fill=popblue, draw=popdark] coordinates {(Cacao,12) (Huile de palme,9) (Caoutchouc,4) (Bananes,2)};
\end{axis}
\end{tikzpicture}
\end{center}
\legende{Principaux produits d'exportation du Kamerun vers 1913 (ordres de grandeur, en milliers de tonnes). Le cacao domine nettement les exportations.}
\end{popfigure}

\courssub{Les infrastructures : chemins de fer, routes, télégraphe, port}

Pour évacuer les produits vers la côte, l'Allemagne construit les \textbf{premières infrastructures modernes} du Cameroun :

\begin{itemize}
\item le \cle{chemin de fer} : la \emph{ligne du Nord} (ou ligne de Bonabéri-Nkongsamba), construite entre 1906 et 1911, relie Douala-Bonabéri à Nkongsamba (environ 160 km) pour évacuer les produits de l'Ouest ; la \emph{ligne du Centre} (ou \emph{Mittellandbahn}), de Douala vers Eseka et l'intérieur, est en construction au début de la guerre de 1914 ;
\item les \cle{routes} et pistes, souvent ouvertes grâce aux corvées des populations ;
\item le \cle{télégraphe}, qui relie les principaux postes administratifs entre eux et à l'Europe par câble sous-marin ;
\item le \cle{port de Douala}, modernisé, qui reste la grande porte d'entrée et de sortie de la colonie.
\end{itemize}

\begin{popfigure}
\begin{center}
\begin{tikzpicture}[scale=0.9, every node/.style={font=\small}]
% Contour très simplifié du Kamerun allemand
\draw[thick, popdark, fill=popcream] (1.2,0) -- (3.4,0.3) -- (4.6,1.2) -- (5.2,2.6) -- (4.8,4.0) -- (3.6,5.2) -- (2.4,5.6) -- (1.6,4.6) -- (0.8,3.2) -- (0.4,1.8) -- cycle;
% Villes
\fill[popink] (0.9,1.6) circle (2pt) node[left=1pt]{Douala};
\fill[popink] (0.7,2.6) circle (2pt) node[left=1pt]{Buéa};
\fill[popink] (0.55,2.2) circle (1.5pt) node[left=1pt]{Victoria};
\fill[popink] (2.6,2.4) circle (2pt) node[above=1pt]{Yaoundé};
\fill[popink] (1.9,2.6) circle (1.5pt) node[above=1pt]{Nkongsamba};
\fill[popink] (3.4,4.4) circle (1.5pt) node[above=1pt]{Garoua};
% Chemin de fer ligne du Nord
\draw[line width=2pt, poporange] (0.9,1.6) -- (1.9,2.6);
% Ligne du Centre (en construction)
\draw[line width=2pt, poporange, dashed] (0.9,1.6) -- (2.6,2.4);
% Plantations côtières
\fill[popgreen!60] (0.4,2.0) rectangle (0.75,2.9);
\node[popgreen!50!black] at (1.05,3.35) {Plantations};
% Nord
\draw[-{Stealth}, thick, popdark] (5.6,4.6) -- (5.6,5.4) node[above]{N};
% Légende interne
\draw[line width=2pt, poporange] (3.0,0.5) -- (3.6,0.5) node[right, popdark]{rail (ligne du Nord)};
\draw[line width=2pt, poporange, dashed] (3.0,0.15) -- (3.6,0.15) node[right, popdark]{rail en chantier (Centre)};
\end{tikzpicture}
\end{center}
\legende{Croquis des infrastructures du Kamerun allemand : ligne du Nord Douala-Nkongsamba, ligne du Centre en construction, plantations de la côte (zone verte) et principales villes. Échelle indicative, tracé approximatif.}
\end{popfigure}

\begin{attention}[Erreur fréquente]
N'attribue pas toutes les infrastructures coloniales aux Français ou aux Britanniques ! Les \textbf{premières lignes de chemin de fer} du Cameroun (ligne du Nord vers Nkongsamba, début de la ligne du Centre), le télégraphe et la modernisation du port de Douala sont des réalisations \textbf{allemandes}, antérieures à 1914.
\end{attention}

\courssub{Le poids de l'exploitation : travail forcé et impôts}

Ces réalisations ont un prix, payé par les populations camerounaises. L'administration instaure l'\cle{impôt} (impôt de capitation, payable en argent), qui oblige les paysans à travailler pour les colons ou à vendre leurs récoltes afin de se procurer de la monnaie. Surtout, elle impose le travail forcé.

\begin{definition}[Travail forcé]
Le \cle{travail forcé} est un travail imposé gratuitement ou mal payé aux populations par l'administration coloniale : portage des charges, construction des routes et des voies ferrées, travail dans les plantations.
\end{definition}

Le portage et les chantiers du chemin de fer font de nombreuses victimes : épuisement, maladies, accidents, répression des récalcitrants. C'est ce poids de l'exploitation qui explique les résistances étudiées dans la section suivante.

\courssub{Les réalisations sanitaires et scolaires}

La colonisation allemande n'est pas seulement destructrice : elle laisse aussi des réalisations sanitaires et scolaires, même si celles-ci servent d'abord les intérêts coloniaux (main-d'œuvre en bonne santé, auxiliaires formés).

Sur le plan \cle{sanitaire}, des \cle{dispensaires} et des hôpitaux sont ouverts (notamment à Douala et Yaoundé), et des campagnes de lutte contre les grandes endémies comme la maladie du sommeil sont menées.

Sur le plan \cle{scolaire}, l'enseignement est surtout l'œuvre des \cle{missions} : les missionnaires \cle{protestants de Bâle} (installés dès 1886) et les missionnaires \cle{catholiques Pallottins} (arrivés en 1890) ouvrent des écoles où l'on enseigne la lecture, l'écriture, le catéchisme et des métiers manuels. Ces écoles forment les premiers employés, instituteurs et catéchistes camerounais.

\begin{savaistu}[La Pagode de Douala]
La célèbre \textbf{Pagode de Douala}, ancien palais du roi Auguste Manga Ndumbe Bell, a été construite en \textbf{1905} par des entrepreneurs allemands. Son architecture originale, qui évoque une pagode asiatique, en fait aujourd'hui l'un des monuments les plus connus de la ville et un témoin vivant de la période allemande.
\end{savaistu}

\begin{exoresolu}[Analyser un bilan économique]
Énoncé. Un élève affirme : « L'Allemagne n'a rien construit au Cameroun, elle a seulement exploité. » À partir du cours, rédige un paragraphe argumenté (5 à 8 lignes) qui nuance cette affirmation.
\tcblower
Solution détaillée. L'affirmation est excessive. Certes, l'objectif allemand est l'exploitation : concessions aux compagnies (Südkamerun Gesellschaft), cultures de rente (cacao, caoutchouc, huile de palme), impôts et travail forcé pèsent lourdement sur les populations. Mais l'Allemagne laisse aussi des réalisations concrètes : les premières lignes de chemin de fer (ligne du Nord Douala-Nkongsamba, début de la ligne du Centre), le télégraphe, la modernisation du port de Douala, des dispensaires et, grâce aux missions de Bâle et des Pallottins, les premières écoles. Le bilan est donc double : une exploitation brutale, mais aussi les fondations des infrastructures modernes du Cameroun.
\end{exoresolu}

\begin{aretenir}[L'essentiel de la section]
Le Kamerun passe du statut de protectorat administré par les compagnies à celui de colonie impériale directe dirigée par un gouverneur (Soden, Zimmerer, Puttkamer, Seitz), avec trois capitales successives : Douala, Buéa puis Yaoundé. L'économie repose sur les concessions, les cultures de rente (cacao : environ 12 000 tonnes exportées vers 1913) et le travail forcé. Les Allemands construisent les premières infrastructures (chemin de fer, télégraphe, port) et les missions ouvrent les premières écoles.
\end{aretenir}

\coursec{Les Camerounais face à l'administration allemande : les résistances}

Dans la section précédente, nous avons vu comment l'administration allemande a organisé le Kamerun : un gouverneur tout-puissant à Douala puis à Yaoundé, des compagnies concessionnaires qui exploitent le caoutchouc et l'huile de palme, des plantations qui réclament une main-d'œuvre abondante et bon marché. Mais cette organisation, imposée de l'extérieur, n'a pas été acceptée passivement par les populations. Partout, du littoral à l'intérieur, les Camerounais ont réagi : par les armes quand ils le pouvaient, par le refus et la fuite quand les armes étaient impossibles. C'est cette longue histoire des résistances que nous allons maintenant étudier.

\courssub{Comprendre la résistance : définition et causes}

\begin{definition}[La résistance]
La \cle{résistance} est l'ensemble des actions, armées ou passives, par lesquelles les populations s'opposent à la domination coloniale. Elle peut prendre la forme de révoltes ouvertes (guerres, soulèvements), mais aussi de comportements de refus : fuite, refus du travail, destruction de récoltes, pétitions, migrations.
\end{definition}

Pour comprendre pourquoi les Camerounais se sont soulevés, il faut d'abord comprendre ce qu'ils ont subi. Les causes des résistances sont directement liées aux méthodes de la colonisation allemande :

\begin{itemize}
\item \textbf{La conquête militaire} : le traité de 1884 n'avait été signé que par quelques chefs du littoral. Pour imposer son autorité sur tout le territoire, l'Allemagne a mené de nombreuses expéditions militaires (\emph{Strafexpeditionen}, « expéditions punitives ») contre les peuples de l'intérieur qui refusaient de se soumettre.
\item \textbf{Le travail forcé} : les plantations et les chantiers (routes, chemin de fer) exigeaient des milliers de porteurs et d'ouvriers. Les villages devaient fournir des hommes, souvent arrachés à leurs champs pendant des mois, dans des conditions terribles.
\item \textbf{Les corvées et les châtiments} : les retards ou refus étaient punis de coups de fouet (la \cle{chicotte}), d'emprisonnement, de prises d'otages. Ces brutalités ont nourri une haine profonde de l'administration.
\item \textbf{Les expropriations foncières} : les terres des villages étaient confisquées pour les plantations et les concessions. À Douala, le projet d'expropriation des Duala de Bonanjo en 1912-1914 en est l'exemple le plus célèbre.
\item \textbf{Le mépris des coutumes} : les Allemands imposaient leurs chefs, ignoraient les autorités traditionnelles légitimes et bousculaient les hiérarchies locales.
\end{itemize}

\begin{attention}[Erreur fréquente]
Ne réduis jamais la résistance aux seules révoltes armées. La \cle{résistance passive} — refus du travail, fuites vers les forêts ou les territoires voisins, migrations entières de villages, destruction des récoltes pour priver l'occupant de ressources — est tout aussi importante. Elle montre que même les peuples incapables de combattre militairement ont refusé la domination coloniale.
\end{attention}

\courssub{La résistance des Duala sur le littoral}

Les Duala, premiers en contact avec les Allemands, sont aussi parmi les premiers à résister.

\textbf{La révolte de Lock Priso (1884).} Dès décembre 1884, quelques mois seulement après la signature du traité germano-douala, le chef \cle{Lock Priso} (de Hickory Town, l'actuel Bonabéri) se soulève contre la présence allemande. Soutenu par des intérêts commerciaux britanniques qui craignent de perdre leurs positions, il attaque les partisans du roi Bell. La révolte est écrasée par la marine allemande, qui bombarde les villages rebelles. Cet épisode montre que la conquête allemande n'a jamais été une simple « signature de traité » : elle a été violente dès le premier jour.

\textbf{L'expropriation des Duala de Bonanjo (1912-1914).} Trente ans plus tard, le conflit rebondit autour de la terre. L'administration allemande décide de déplacer les habitants de Bonanjo, quartier historique des Duala à Douala, pour créer une ville européenne séparée de la ville indigène. Les Duala, expropriés de leurs terres ancestrales, réagissent par une arme nouvelle : la \cle{pétition}. Sous la conduite du chef \cle{Rudolf Duala Manga Bell}, ils adressent des protestations au gouverneur, puis au \cle{Reichstag} (le parlement allemand) à Berlin. Cette résistance juridique et diplomatique, menée par des chefs lettrés, inquiète profondément l'administration coloniale.

\courssub{La résistance des peuples de l'intérieur}

L'occupation de l'intérieur du territoire s'est faite contre des peuples déterminés à défendre leur indépendance.

\begin{itemize}
\item \textbf{Les Bafut} (Grassfields, Ouest) : le fon (roi) des Bafut, Abumbi Ier, affronte les Allemands dans une guerre acharnée (1901-1907). En 1907, le palais royal est détruit et le fon doit s'enfuir ; l'administration allemande impose ensuite son autorité.
\item \textbf{Les Bamoun} : le sultan \cle{Ibrahim Njoya}, après une première phase de collaboration active, entre en conflit avec l'administration allemande qui le soupçonne de double jeu à l'approche de la Première Guerre mondiale. Convoqué à Yaoundé en 1916, il y est retenu par les Allemands jusqu'à leur départ.
\item \textbf{Les Bakweri} (pentes du mont Cameroun) : ils mènent une guérilla contre les plantations qui accaparent leurs terres et réclament du travail forcé.
\item \textbf{Les Bulu} (Sud) : peuple de guerriers redoutés, ils opposent une résistance armée farouche à la pénétration allemande dans les années 1890-1900.
\item \textbf{Les Bassa} : ils résistent notamment aux recrutements forcés pour la construction du chemin de fer (ligne Nord et ligne principale).
\end{itemize}

\begin{popfigure}
\begin{tikzpicture}[scale=1.05, font=\small]
% Contour très simplifié du Kamerun
\draw[thick, popdark, fill=popcream] (2.2,0) -- (3.6,0.4) -- (4.6,1.4) -- (4.4,2.6) -- (5.0,3.6) -- (4.6,4.8) -- (3.4,5.6) -- (2.0,6.4) -- (1.2,6.0) -- (0.8,4.8) -- (0.2,3.4) -- (0.0,2.0) -- (0.6,0.8) -- cycle;
% Côte
\draw[popblue, thick] (0.0,2.0) -- (0.6,0.8);
\node[popblue, rotate=55] at (-0.35,1.3) {\scriptsize Atlantique};
% Foyers de résistance
\fill[poppink] (0.7,1.1) circle (2.2pt); \node[right, align=left] at (0.85,1.1) {\scriptsize \textbf{1. Douala} : Lock Priso 1884 ;\\ \scriptsize pétition Duala 1912-1914};
\fill[poppink] (1.5,0.7) circle (2.2pt); \node[right] at (1.65,0.7) {\scriptsize \textbf{2. Ebolowa} : M.-P. Samba (Bulu)};
\fill[poppink] (0.5,2.6) circle (2.2pt); \node[right] at (0.65,2.6) {\scriptsize \textbf{3. Bakweri} (mont Cameroun)};
\fill[poppink] (1.3,3.6) circle (2.2pt); \node[right] at (1.45,3.6) {\scriptsize \textbf{4. Bafut} : guerre 1901-1907};
\fill[poppink] (2.4,4.3) circle (2.2pt); \node[right] at (2.55,4.3) {\scriptsize \textbf{5. Bamoun} : tensions avec Njoya};
\fill[poppink] (1.9,1.9) circle (2.2pt); \node[right] at (2.05,1.9) {\scriptsize \textbf{6. Bassa} : refus chemin de fer};
% Nord
\draw[->, thick, popdark] (4.9,5.6) -- (4.9,6.3); \node[above] at (4.9,6.3) {\scriptsize N};
\end{tikzpicture}
\legende{Croquis des principaux foyers de résistance au Kamerun (1884-1916). Les points roses indiquent les zones de révoltes armées ou de résistance organisée. Contour schématique, sans précision cartographique.}
\end{popfigure}

\courssub{Les grandes figures de la résistance}

\textbf{Martin-Paul Samba.} Né chez les Bulu vers 1875, \cle{Martin-Paul Samba} (de son nom de naissance Mebenga m'Ebono) a été envoyé se former en Allemagne, où il a servi dans l'armée allemande. Rentré au pays, il est nommé chef de canton par l'administration. Mais la connaissance qu'il a des Allemands nourrit sa révolte : il devient l'un des plus farouches opposants à la colonisation et prépare un soulèvement général à partir d'Ebolowa.

\begin{exemplebox}[Un résistant formé par l'adversaire]
Martin-Paul Samba illustre un paradoxe fréquent : les colonisateurs ont formé eux-mêmes certains de leurs adversaires les plus déterminés. Scolarisé et militairement entraîné en Allemagne, Samba connaît les forces et les faiblesses de l'occupant. Il utilise cette connaissance pour organiser la résistance bulu et correspondre avec d'autres chefs, dont Duala Manga Bell.
\end{exemplebox}

\textbf{Rudolf Duala Manga Bell.} Chef des Duala de Bonanjo, \cle{Rudolf Duala Manga Bell} avait été scolarisé en Allemagne et parlait couramment allemand. Face à l'expropriation de son peuple, il mène la campagne de pétitions et cherche des alliances avec d'autres chefs (dont Samba) pour faire échec au projet allemand.

\textbf{Madola.} Chef \cle{bakoko} de la région d'Edéa, \cle{Madola} refuse de fournir des travailleurs forcés aux plantations et aux chantiers. Poursuivi, il se réfugie dans une grotte-forteresse (Ngog Lituba) et tient tête aux troupes coloniales avant d'être capturé et exécuté vers 1910. Sa résistance est restée célèbre dans la mémoire locale.

\courssub{Le complot de 1914 et les exécutions du 8 août 1914}

En 1914, l'administration allemande découvre la correspondance entre Rudolf Duala Manga Bell et Martin-Paul Samba, ainsi que leurs contacts avec d'autres chefs et, selon l'accusation, avec les puissances ennemies de l'Allemagne (France, Grande-Bretagne). Les deux hommes sont arrêtés et accusés de \cle{haute trahison}. Un procès expéditif, sans garanties judiciaires réelles, les condamne à mort. Le \textbf{8 août 1914}, Rudolf Duala Manga Bell et Martin-Paul Samba sont \textbf{pendus à Douala}, au moment même où éclate la Première Guerre mondiale en Europe. Ngoso Din, secrétaire et conseiller de Manga Bell, subit le même sort.

\begin{savaistu}[Un héros national]
Rudolf Duala Manga Bell, qui parlait couramment allemand et connaissait l'Europe, est aujourd'hui considéré comme un \cle{héros national} du Cameroun. Des rues, des écoles et un lycée de Douala portent son nom. Sa pendaison le 8 août 1914 est commémorée comme l'un des moments fondateurs de la conscience nationale camerounaise.
\end{savaistu}

\begin{popfigure}
\begin{tikzpicture}[font=\small]
% Cadre photo schématique
\draw[thick, popdark, fill=popcream] (0,0) rectangle (9,5.2);
% Personnage 1 : Manga Bell
\draw[fill=popblueL, draw=popdark] (2.3,2.2) circle (0.55);
\draw[fill=popblue!40, draw=popdark] (1.3,0.4) rectangle (3.3,1.9);
\draw[popdark] (2.3,1.9) -- (2.3,1.65);
\node[align=center] at (2.3,-0.35) {\scriptsize \textbf{Rudolf Duala Manga Bell}\\ \scriptsize costume européen, allure digne};
% Personnage 2 : Samba
\draw[fill=poporangeL, draw=popdark] (6.7,2.2) circle (0.55);
\draw[fill=poporange!40, draw=popdark] (5.7,0.4) rectangle (7.7,1.9);
\draw[popdark] (6.7,1.9) -- (6.7,1.65);
\node[align=center] at (6.7,-0.35) {\scriptsize \textbf{Martin-Paul Samba}\\ \scriptsize ancien soldat, regard déterminé};
% Annotations
\draw[->, poppink, thick] (4.5,4.6) -- (2.6,2.9);
\node[poppink] at (4.5,4.85) {\scriptsize tenue soignée : chefs lettrés, « évolués »};
\draw[->, poppink, thick] (4.5,0.9) -- (6.4,1.5);
\node[poppink, align=center] at (4.5,0.55) {\scriptsize posture fière :\\ \scriptsize dignité face à l'occupant};
\end{tikzpicture}
\legende{Schéma reproduisant une photographie historique type de Rudolf Duala Manga Bell et de Martin-Paul Samba. À observer : les vêtements européens (ces chefs maîtrisent les codes de l'occupant), la posture droite et fière (dignité, refus de la soumission), le contexte (photographie de propagande ou de procès, vers 1914).}
\end{popfigure}

\begin{methode}[Analyser une photographie historique]
\begin{enumerate}
\item \textbf{Décrire} : que voit-on ? Personnages, vêtements, attitudes, décor, sans interpréter.
\item \textbf{Identifier} : qui sont les personnages ? Quelle est la source et la date du document ?
\item \textbf{Situer le contexte} : à quel moment de la période coloniale la photo renvoie-t-elle ? Ici, la lutte contre l'expropriation des Duala et la répression de 1914.
\item \textbf{Dégager le message} : que veut montrer l'auteur de la photo ? Que révèle-t-elle malgré lui ? Ici, la dignité des chefs résistants contredit l'image de « sauvages » que la propagande coloniale voulait imposer.
\end{enumerate}
\end{methode}

\courssub{Les formes de résistance passive et le bilan des résistances}

Face à la \cle{Schutztruppe}, l'armée coloniale allemande bien armée et organisée, beaucoup de populations ne pouvaient pas combattre ouvertement. Elles ont alors pratiqué la résistance passive :

\begin{itemize}
\item les \textbf{fuites} : les villageois abandonnent leurs villages à l'arrivée des colonnes allemandes et se réfugient dans la forêt ;
\item le \textbf{refus du travail} : ralentissements, désertions des chantiers et des plantations ;
\item les \textbf{migrations} : des groupes entiers se déplacent vers les territoires français (Congo français, Oubangui-Chari) ou espagnol (Rio Muni) voisins pour échapper à l'administration allemande ;
\item la \textbf{destruction des récoltes} : pour priver l'occupant de vivres et décourager son avancée.
\end{itemize}

\begin{exoresolu}[Pourquoi les résistances ont-elles échoué ?]
Énoncé. À partir du cours, dégage deux raisons expliquant l'échec des résistances camerounaises face à l'Allemagne entre 1884 et 1914.
\tcblower
\textbf{Solution.} Première raison : les résistances sont \textbf{nombreuses mais dispersées}. Chaque peuple (Duala, Bulu, Bafut, Bakweri, Bassa, Bamoun) combat séparément, à des dates différentes, sans coordination générale ; les tentatives d'alliance (correspondance Manga Bell–Samba) sont découvertes et brisées. Deuxième raison : la \textbf{supériorité militaire allemande}. La Schutztruppe dispose d'armes à feu modernes (fusils à répétition, mitrailleuses, canons) face à des guerriers armés de lances, d'arcs et de quelques fusils anciens. Les expéditions punitives écrasent donc les révoltes les unes après les autres.
\end{exoresolu}

\begin{aretenir}[Bilan de la section]
Les Camerounais ont résisté à la colonisation allemande par les armes (Lock Priso en 1884, Bafut, Bulu, Bakweri, Madola) et par des moyens pacifiques ou passifs (pétition des Duala, fuites, refus du travail, migrations vers les colonies voisines). Les figures majeures sont Rudolf Duala Manga Bell et Martin-Paul Samba, pendus le 8 août 1914 à Douala pour « haute trahison ». Dispersées et confrontées à la supériorité militaire de la Schutztruppe, ces résistances ont été écrasées, mais elles ont légué au Cameroun une mémoire héroïque fondatrice.
\end{aretenir}

\coursec{Collaborations et fin de la période allemande}

\courssub{Introduction : des résistances aux stratégies d'adaptation}

Après avoir analysé les formes multiples de la résistance camerounaise face à la pénétration et à l'installation de l'administration allemande (section précédente), il est indispensable d'étudier l'autre volet de la relation coloniale : la \cle{collaboration}. L'histoire coloniale n'est pas un bloc monolithique de refus ; elle est traversée de stratégies complexes où des acteurs africains, pour des motifs variés (survie, intérêt dynastique, rivalités locales, conviction), choisissent de coopérer avec l'occupant. Cette section examine ces dynamiques, puis raconte comment le premier conflit mondial met brutalement fin à la présence allemande et redessine la carte du territoire.

\courssub{La collaboration : définition, motivations et acteurs}

\begin{definition}[Collaboration coloniale]
Attitude consistant, pour un individu ou un groupe africain, à coopérer activement avec l'administration coloniale, que ce soit par conviction idéologique, par recherche d'intérêts matériels ou politiques (protection, armes, prestige), ou sous la contrainte (menace de déposition, répression). Elle se distingue de la simple \emph{résignation} (subir sans agir) par l'engagement actif aux côtés de l'occupant.
\end{definition}

\begin{aretenir}[Les trois moteurs principaux de la collaboration]
\begin{enumerate}
    \item \textbf{La protection et la survie politique :} Face à la puissance de feu allemande (mitrailleuses, artillerie), certains chefs préfèrent traiter pour éviter la destruction de leur peuple et de leur capitale (ex. : Sultan Njoya).
    \item \textbf{Les intérêts économiques et le renforcement du pouvoir :} L'alliance donne accès aux armes à feu, aux marchandises d'importation et à l'appui militaire allemand contre des rivaux traditionnels.
    \item \textbf{Les rivalités inter-ethniques ou dynastiques :} L'ennemi de mon ennemi est mon ami. Des chefs collaborent pour faire tomber un rival récalcitrant ou pour s'affranchir d'une tutelle traditionnelle (ex. : certains chefs Bassa ou Bamiléké face à des rois traditionnels hostiles aux Allemands).
\end{enumerate}
\end{aretenir}

\begin{exemplebox}[Le rôle des auxiliaires de l'administration]
L'administration allemande, sous-administrée (peu d'officiers pour un vaste territoire), s'appuie massivement sur des \cle{chefs auxiliaires} (souvent créés \emph{ex nihilo} ou choisis parmi des notables dociles). Leurs missions :
\begin{itemize}
    \item Perception de l'\cle{impôt de capitation} (institué en 1909, payable en espèces ou en travail).
    \item Recrutement de la \cle{main-d'œuvre forcée} pour les chantiers (routes, chemin de fer, plantations) et comme porteurs pour les colonnes militaires.
    \item Maintien de l'ordre (police indigène) et renseignement.
\end{itemize}
En contrepartie, ils reçoivent un pourcentage sur l'impôt, des médailles, des uniformes, et la confirmation de leur autorité par le \emph{Bezirksamtmann} (chef de circonscription).
\end{exemplebox}

\courssub{Le cas emblématique : le Sultan Njoya et le royaume Bamoun}

\begin{savaistu}[Le Sultan Njoya (vers 1867-1933) : un modernisateur visionnaire]
Ibrahim Njoya, 17e sultan des Bamoun (Foumban), comprend très tôt (dès 1902-1903) que la résistance armée est vouée à l'échec face à la \emph{Schutztruppe}. Il choisit une \cle{alliance stratégique} avec les Allemands :
\begin{itemize}
    \item Il fournit des porteurs, des vivres et des renseignements ; il envoie ses fils à l'école allemande de Buea.
    \item En retour, les Allemands lui laissent une large autonomie interne, lui fournissent des armes et reconnaissent son autorité sur les chefferies voisines.
\end{itemize}
\underline{Héritage exceptionnel :} Njoya n'est pas qu'un collaborateur ; il est un réformateur génial. Il invente l'\textbf{alphabet Shümom} (vers 1896-1910, passé de 500 pictogrammes à 80 syllabes puis 35 lettres) pour écrire le bamoun, fait bâtir le \textbf{Palais royal de Foumban} (inspiré de l'architecture allemande mais aux symboles bamoun), développe l'agriculture (cacao, café), crée un musée, une pharmacopée et une administration écrite. Déposé par les Français en 1931, il meurt en exil à Yaoundé. Son palais est aujourd'hui un musée classé au patrimoine mondial de l'UNESCO (2023).
\end{savaistu}

\begin{attention}[Piège historiographique : l'anachronisme moral]
\emph{Erreur fréquente :} Juger les collaborateurs (Njoya, chefs auxiliaires, soldats) avec les valeurs nationalistes d'aujourd'hui (\og traitres \fg{}).
\emph{Rigueur historique :} Il faut \textbf{contextualiser}. Au tournant du XXe siècle, l'État-nation camerounais n'existe pas. Les loyautés sont dynastiques, lignagères ou religieuses. Njoya sauve son royaume et le modernise ; un chef Bassa collabore pour échapper à la tutelle Bamoun ou Tikar. L'historien analyse les \textbf{logiques d'acteurs} dans leurs contraintes structurelles (rapport de forces militaire, économie de la violence coloniale), sans les exonérer de leur responsabilité dans la répression de leurs propres frères.
\end{attention}

\courssub{Les Camerounais dans l'appareil militaire et administratif allemand}

Au-delà des chefs, des milliers de Camerounais intègrent les rouages de l'État colonial :

\begin{itemize}
    \item \textbf{La \cle{Schutztruppe} (troupe de protection) :} Composée d'officiers allemands et de \cle{soldats africains} (recrutés au Nord, au Cameroun même, ou venus du Togo/Dahomey). Ils assurent la conquête, la \emph{pacification} et la garde des postes. Leur fidélité tient à la solde régulière, au prestige de l'uniforme et à l'esprit de corps.
    \item \textbf{Les \cle{interprètes} :} Rouage indispensable. Maîtrisant l'allemand, le \emph{Pidgin English} (côte) ou le \emph{Fulfulde} (nord), ils sont les vrais médiateurs du pouvoir. Certains (comme \textbf{Charles Atangana} chez les Ewondo, bien que son rôle majeur soit plus tard sous mandat français) accumulent un pouvoir considérable.
    \item \textbf{Le personnel subalterne :} Infirmiers, commis, enseignants d'écoles gouvernementales, employés du chemin de fer (Nordbahn, Mittelbahn).
\end{itemize}

\begin{exoresolu}[Analyse d'un document : le rôle de l'interprète]
\textbf{Document :} Extrait du rapport du Gouverneur Seitz (1913) : \og Les interprètes sont les yeux et les oreilles du district. Sans eux, l'administration est aveugle et sourde. Mais leur pouvoir personnel est grand ; il faut les surveiller. \fg{}

\textbf{Question :} En quoi ce document illustre-t-il la position ambivalente des interprètes ?

\textbf{Solution détaillée :}
\begin{enumerate}
    \item \textbf{Identification :} Source administrative allemande (Gouverneur), 1913, contexte de consolidation de l'administration civile.
    \item \textbf{Analyse :} La métaphore \og yeux et oreilles \fg{} montre la \textbf{dépendance structurelle} de l'État colonial vis-à-vis de ses auxiliaires africains (barrière linguistique, ignorance des coutumes).
    \item \textbf{Nuance :} L'aveu \og leur pouvoir personnel est grand \fg{} révèle une \textbf{délégation de fait de la souveraineté}. L'interprète filtre l'information, peut arranger les procès, percevoir des pots-de-vin.
    \item \textbf{Conclusion :} L'interprète est un \textbf{acteur de pouvoir} à part entière, ni simple outil, ni simple victime. Il incarne la \cle{co-production de l'ordre colonial}.
\end{enumerate}
\end{exoresolu}

\courssub{La Première Guerre mondiale au Kamerun (1914-1916) : la fin brutale}

L'entrée en guerre de l'Allemagne (août 1914) transforme le Kamerun en théâtre d'opérations. La \cle{Schutztruppe} (env. 200 officiers allemands + 1 500 à 2 000 soldats africains) fait face à une coalition franco-britannique très supérieure en nombre (env. 18 000 hommes au total) et en logistique.

\begin{popfigure}
\begin{tikzpicture}[scale=0.75, every node/.style={font=\small}]
    % Contour très simplifié du Kamerun (polygone approximatif)
    \draw[popdark, thick, fill=popcream] 
        (0,0) -- (10,0) -- (11,3) -- (10.5,6) -- (9,8) -- (7,9) -- (4,8.5) -- (1,7) -- (0.5,4) -- (0,1) -- cycle;
    
    % Villes clés
    \node[fill=popblue, circle, inner sep=2pt, text=white] (garoua) at (2.5, 7.2) {};
    \node[above=2pt of garoua, popblue, font=\bfseries] {Garoua};
    
    \node[fill=poporange, circle, inner sep=2pt, text=white] (mora) at (4.5, 8.0) {};
    \node[above=2pt of mora, poporange, font=\bfseries] {Mora};
    
    \node[fill=popgreen, circle, inner sep=2pt, text=white] (yaounde) at (6.5, 4.5) {};
    \node[right=2pt of yaounde, popgreen, font=\bfseries] {Yaoundé (Jaunde)};
    
    \node[fill=poppurple, circle, inner sep=2pt, text=white] (douala) at (8.5, 1.5) {};
    \node[below right=2pt of douala, poppurple, font=\bfseries] {Douala};
    
    \node[fill=poppink, circle, inner sep=2pt, text=white] (buea) at (9.5, 0.5) {};
    \node[below right=2pt of buea, poppink, font=\bfseries] {Buéa};
    
    % Flèches d'attaque (Français depuis l'Est/Sud, Britanniques depuis l'Ouest/Nord)
    % Attaque Française depuis le Congo (Est)
    \draw[->, thick, popblue, decorate, decoration={snake, amplitude=1mm, segment length=4mm}] (3, 5) -- (5.5, 4.8);
    \node[popblue, rotate=20, anchor=south] at (4, 5.5) {Colonne française (Est)};
    
    % Attaque Britannique depuis le Nigeria (Ouest/Nord)
    \draw[->, thick, poporange, decorate, decoration={snake, amplitude=1mm, segment length=4mm}] (1, 7.5) -- (3, 7.3);
    \node[poporange, rotate=-10, anchor=south] at (1.5, 7.8) {Colonne britannique (Nord)};
    
    \draw[->, thick, poporange, decorate, decoration={snake, amplitude=1mm, segment length=4mm}] (9, 0) -- (7, 2);
    \node[poporange, rotate=30, anchor=south] at (8.5, 0.5) {Colonne britannique (Ouest)};
    
    % Retraite vers Rio Muni (Guinée espagnole)
    \draw[->, very thick, popdark, dashed] (6.5, 4.5) -- (3.5, 3.5) -- (1.5, 2) -- (0, 1.5);
    \node[popdark, rotate=-30, font=\bfseries\itshape] at (3, 3) {Retraite allemande $\rightarrow$ Rio Muni (Fév 1916)};
    
    % Flèche Nord
    \node at (11.5, 8.5) {\Huge \textbf{N}};
    \draw[->, thick] (11.5, 8.0) -- (11.5, 8.4);
\end{tikzpicture}
\legende{Carte schématique des opérations militaires au Kamerun (1914-1916). Les flèches en zigzag symbolisent la progression difficile dans la forêt/brousse. La retraite finale vers la Guinée espagnole (Rio Muni) marque la fin de la présence allemande. Villes clés : Garoua, Mora (Nord), Yaoundé (centre), Douala, Buéa (côte).}
\end{popfigure}

\begin{exemplebox}[Chronologie clé des opérations]
\begin{itemize}
    \item \textbf{Août-Sept 1914 :} Attaques alliées sur Douala (côte) et Garoua (Nord). Les Allemands évacuent la côte, se replient à l'intérieur (Yaoundé, Ngaundéré, Mora).
    \item \textbf{Sept 1914 - Fév 1916 :} \textbf{Siège de Mora} (Monts Mandara, Nord). Le capitaine von Raben résiste 18 mois avec 200 hommes. Reddition honorable (honneurs militaires) en février 1916.
    \item \textbf{Janvier 1916 :} Chute de Yaoundé (capitale). Gouvernement et troupes battent en retraite vers le Sud-Est.
    \item \textbf{Février 1916 :} Derniers combattants allemands (env. 6 000 à 10 000 personnes dont soldats, porteurs, familles) franchissent la frontière vers la \textbf{Guinée espagnole (Rio Muni)}. Ils y sont internés. C'est la fin officielle de la domination allemande.
\end{itemize}
\end{exemplebox}

\begin{exemplebox}[Chiffre clé : Le siège de Mora, une anomalie tactique]
\textbf{Durée :} \textbf{18 mois} (août 1914 -- février 1916).
\textbf{Forces en présence :} ~200 soldats allemands/africains (von Raben) vs ~1 500 à 2 000 alliés (Français/Britanniques).
\textbf{Signification :} Ce siège disproportionné illustre la \textbf{logistique défaillante des Alliés} en zone de montagne/brousse, la compétence tactique allemande, mais aussi la \textbf{fidélité des soldats africains} (Baya, Mafa, Peuls) restés loyaux à von Raben jusqu'au bout, partageant la faim et les privations. La reddition négociée (armes rendues, honneurs militaires) est unique dans la campagne.
\end{exemplebox}

\courssub{Le partage du Kamerun : de la conquête au mandat (1916-1922)}

La défaite allemande pose la question du sort du territoire. Les vainqueurs appliquent le principe du \cle{partage du butin colonial} (accords secrets Sykes-Picot / accords de Londres 1915-1916), contredisant les vœux du Président Wilson (droit des peuples).

\begin{center}
\begin{tabularx}{\linewidth}{|l|X|X|}
\hline
\rowcolor{popblueL} \textbf{Critère} & \textbf{Zone Française (Cameroun français)} & \textbf{Zone Britannique (British Cameroons)} \\
\hline
\textbf{Superficie} & \textbf{~4/5 du territoire} (~495 000 km²) & \textbf{~1/5 du territoire} (~88 000 km²) \\
\hline
\textbf{Population} & Majoritaire (forêt, savane, nord) & Minoritaire (zones frontalières Nigeria) \\
\hline
\textbf{Capitales} & Yaoundé (administrative), Douala (économique) & Buéa (administrative), Victoria (port) \\
\hline
\textbf{Administration} & \textbf{Assimilation / Administration directe} : turnover des cadres français, code de l'indigénat, travail forcé (prestations), mise en valeur économique intensive (plantations, chemins de fer). & \textbf{Indirect Rule (Gouvernement indirect)} : gouverner \emph{à travers} les chefferies traditionnelles (Emirats au Nord, chefferies au Sud), faible investissement infrastructurel, économie tournée vers le Nigeria (main-d'œuvre migrante). \\
\hline
\textbf{Frontière} & \textbf{Ligne Picot (1916) / Ligne Milner-Simon (1919)} : Traitée au ruler, coupant des peuples (Bamoun, Bamiléké, Mafa, Toupouri, Fulbé). & Même frontière. \\
\hline
\textbf{Statut juridique (1922)} & \textbf{Mandat de la SDN (Classe B)} : \og Le bien-être et le développement de ces peuples forment une mission sacrée de civilisation \fg{} (Article 22 Pacte SDN). En pratique : colonisation déguisée. & \textbf{Mandat de la SDN (Classe B)} : Administré \emph{comme partie intégrante du Nigeria}. \\
\hline
\end{tabularx}
\end{center}

\courssub{Bilan de la période allemande (1884-1916) : une entité née dans la contrainte}

\begin{aretenir}[Héritages contradictoires de 32 ans de colonisation allemande]
\begin{itemize}
    \item \textbf{Naissance de l'entité territoriale \og Cameroun \fg{} :} Frontières internationales reconnues (conférences de Berlin 1884, accords 1885, 1894, 1911), unification administrative embryonnaire (chef-lieu unique à Buéa puis Yaoundé), début d'intégration Nord/Sud par le rail et l'administration.
    \item \textbf{Héritages matériels :} Réseau ferroviaire (Nordbahn, Mittelbahn - ~600 km), routes carrossables, bâtiments administratifs (Gouvernorat Buéa/Yaoundé, postes de district), ports (Douala), plantations (cacao, café, palmiers, hévéas), introduction de cultures d'exportation.
    \item \textbf{Héritages mémoriels et symboliques :} Figures de la \textbf{résistance} (Rudolf Douala Manga Bell, Martin Paul Samba, Madola, chef Koupé, sultan Njoya \emph{ambigu}), mémoire des \textbf{violences} (travail forcé, répressions, déportations), mais aussi mémoire de la \textbf{modernisation} (écoles, écriture Shümom, médecine, architecture).
    \item \textbf{Héritage linguistique :} L'allemand langue de l'administration et de l'enseignement supérieur (jusqu'en 1916) ; survivance de l'allemand dans l'onomastique (noms de lieux, prénoms) et dans la mémoire collective du Nord (anciens \emph{Askari}).
\end{itemize}
\end{aretenir}

\begin{methode}[Rédiger une conclusion argumentée (type Bac/Probatoire)]
\textbf{Sujet type :} \og La période allemande (1884-1916) a-t-elle fondé le Cameroun moderne ? \fg{}

\begin{enumerate}
    \item \textbf{Rappeler la problématique :} Reformuler le sujet en question ouverte : \og En quoi 32 ans de colonisation allemande constituent-ils l'acte de naissance territorial et structurel du Cameroun, malgré leur brièveté et leur violence ? \fg{}
    \item \textbf{Synthétiser les acquis (plan dialectique) :}
    \begin{itemize}
        \item \textbf{Thèse (Oui, fondation) :} Création des frontières, unification administrative, infrastructures lourdes (rail, ports), introduction économie de plantation, naissance d'une élite lettrée (allemand), figures héroïques (résistance).
        \item \textbf{Antithèse (Non, simple occupation) :} Durée courte (32 ans), contrôle effectif limité (peu d'Allemands), violence structurelle (travail forcé, code de l'indigénat), pas d'unité nationale (peuples divisés, collaboration/résistance), partage franco-britannique qui brise l'unité allemande.
    \end{itemize}
    \item \textbf{Synthèse nuancée :} L'entité \textbf{territoriale et administrative} est allemande ; la \textbf{nation camerounaise} (conscience d'appartenance commune) se forge \emph{après}, dans la résistance au partage (années 1920-1950) et la réunification (1961). L'héritage est \textbf{structurel} (cadres), pas encore \textbf{identitaire} (peuple).
    \item \textbf{Ouverture :} \og Ce partage franco-britannique de 1916, en brisant l'unité allemande, va-t-il permettre l'émergence d'un nationalisme camerounais \emph{unifié} ou au contraire installer une fracture anglophone/francophone durable ? \fg{} (Transition vers la leçon sur les mandats).
\end{enumerate}
\end{methode}

\courssub{Travaux pratiques : Construire la carte du Kamerun allemand}

\begin{experience}[Réalisation d'une carte schématique du Kamerun (1911-1916)]
\textbf{Objectif :} Produire une carte légendée à la main (format A4) situant les éléments structurants de la colonie allemande à son apogée (après 1911).

\textbf{Matériel :} Feuille blanche A4, crayons de couleur (bleu, rouge, vert, marron, orange), règle, gomme, atlas ou carte de référence (manuel).

\textbf{Protocole (étapes) :}
\begin{enumerate}
    \item \textbf{Contour et repères :} Tracer le contour général du Kamerun (forme grossière en "pentagone irrégulier" élargi au Nord). Tracer l'équateur (traversant le Sud) et le méridien de Greenwich (à l'Ouest de Douala). Placer la flèche Nord.
    \item \textbf{Hydrographie structurante :} Dessiner en \textbf{bleu} : le fleuve Sanaga (axe Nord-Sud central), la Bénoué (affluent du Niger, vers le Nord-Ouest), le Logone/Chari (vers le Lac Tchad, Nord-Est), le Nyong, le Wouri (Douala).
    \item \textbf{Villes et postes administratifs :} Positionner et nommer en \textbf{noir} (souligné) : \textbf{Buéa} (1er chef-lieu), \textbf{Yaoundé/Jaunde} (chef-lieu 1909), \textbf{Douala} (port principal), \textbf{Garoua}, \textbf{Ngaoundéré}, \textbf{Mora}, \textbf{Foumban}, \textbf{Dschang}, \textbf{Bertoua}, \textbf{Kribi}, \textbf{Victoria/Limbe}.
    \item \textbf{Infrastructures de transport :} Tracer en \textbf{rouge} les deux lignes de chemin de fer : \emph{Nordbahn} (Douala -- Nkongsamba -- Kumba) et \emph{Mittelbahn} (Douala -- Yaoundé -- Ngaoundéré \emph{projeté/embranchement}). Tracer en \textbf{marron} les axes routiers principaux (pistes carrossables).
    \item \textbf{Économie :} Symboles en \textbf{vert} : zones de plantations (cacao/café : Sud-Ouest, littoral ; palmiers : littoral ; hévéas : Sud-Est). Mines (étain : Nord-Ouest, symbolisé par un picto marteau).
    \item \textbf{Limite 1911 (Neukamerun) :} Tracer en \textbf{orange pointillé} l'extension vers l'Est (Sangha, Lobaye) et le Nord (jusqu'au Lac Tchad) acquise en 1911 (Traité de Fès).
    \item \textbf{Légende organisée :} En bas à droite, encadrer une légende structurée : \textit{Relief/Hydro}, \textit{Villes/Admin}, \textit{Transports}, \textit{Économie}, \textit{Frontières (1884 / 1911)}.
    \item \textbf{Titre :} \og Le Kamerun allemand à son apogée (1911-1916) \fg.
\end{enumerate}

\textbf{Critères d'évaluation :} Respect de l'orientation, proportion grossière du contour, exhaustivité des items (villes, rails, fleuves), lisibilité de la légende, propreté.
\end{experience}

\courssub{Exercice d'application : Tableau comparatif Résistants / Collaborateurs}

\begin{exercice}[Mise en perspective des stratégies africaines]
Complétez le tableau ci-dessous en vous appuyant sur le cours et vos connaissances. Utilisez des noms précis (au moins 2 par colonne).
\end{exercice}

\begin{center}
\begin{tabularx}{\linewidth}{|l|X|X|}
\hline
\rowcolor{popblueL} \textbf{Critère de comparaison} & \textbf{Résistants (exemples : Rudolf Douala Manga Bell, Martin Paul Samba, Madola, Sultan Njoya \emph{phase initiale})} & \textbf{Collaborateurs / Auxiliaires (exemples : Chefs auxiliaires Bassa/Bamiléké, Interprètes, Soldats Schutztruppe, Sultan Njoya \emph{phase alliance})} \\
\hline
\textbf{Motivations principales} & & \\
\hline
\textbf{Moyens d'action} & & \\
\hline
\textbf{Relations avec l'administration allemande} & & \\
\hline
\textbf{Destin final (souvent)} & & \\
\hline
\textbf{Héritage mémoriel aujourd'hui} & & \\
\hline
\end{tabularx}
\end{center}

\begin{corrige}[Exercice : Tableau comparatif Résistants / Collaborateurs]
\begin{center}
\begin{tabularx}{\linewidth}{|l|X|X|}
\hline
\rowcolor{popblueL} \textbf{Critère} & \textbf{Résistants} & \textbf{Collaborateurs / Auxiliaires} \\
\hline
\textbf{Motivations} & Défense de la souveraineté, refus de l'impôt/travail forcé, fidélité aux ancêtres/terre, rejet de l'arbitraire (ex. expropriation Duala). & Protection du peuple/royaume, intérêts dynastiques (armes, territoire), rivalités locales, carrière/salaire (soldats, interprètes), conviction modernisatrice (Njoya). \\
\hline
\textbf{Moyens} & Guérilla, embuscades, refus de l'impôt, pétitions juridiques (Manga Bell au Reichstag), diplomatie, destruction symbolique (bornes). & Perception impôt, recrutement porteurs/soldats, renseignement, police, médiation culturelle, modernisation interne (écoles, écriture, palais). \\
\hline
\textbf{Relations avec Allemands} & Conflit ouvert, répression, trahison perçue par Allemands. & Partenariat inégal, dépendance (révocabilité du chef), surveillance mutuelle (cf. rapport Seitz). \\
\hline
\textbf{Destin final} & Exécution (Manga Bell 1914, Samba 1914), mort au combat, exil, déportation. & Maintien au pouvoir (provisoire), certains basculent vers Français (Atangana), d'autres internés à Rio Muni (soldats), Njoya déposé par Français (1931). \\
\hline
\textbf{Héritage mémoriel} & \textbf{Héros nationaux} (rues, statues, hymnes), symboles de l'indépendance. & \textbf{Figures ambivalentes} : Njoya réhabilité (génie culturel) ; chefs auxiliaires souvent oubliés ou critiqués ; soldats \emph{Askari} respectés au Nord pour bravoure. \\
\hline
\end{tabularx}
\end{center}
\end{corrige}

\coursec{Travaux pratiques : construire la carte du Kamerun}

Après avoir étudié les collaborations et la fin de la période allemande en 1916, nous savons désormais que le Kamerun n'est pas un territoire « donné » : il a été \cle{construit} par les traités, les conquêtes et les accords de frontières, et \cle{contesté} par les résistances des populations. Mais comment rendre visible cette histoire ? Comment montrer, d'un seul coup d'œil, qu'un territoire a grandi, s'est organisé et a été combattu ? La réponse des historiens est la \cle{carte historique}. Cette séance de travaux pratiques va vous transformer en cartographes : vous allez construire vous-mêmes la carte du Kamerun allemand, de la légende jusqu'au croquis de synthèse.

\courssub{Pourquoi cartographier le Kamerun ?}

Une carte historique n'est pas un dessin décoratif : c'est un \emph{raisonnement spatialisé}. Tracer trois frontières successives (1884, 1901, 1911), c'est démontrer que le territoire camerounais est le résultat d'un processus. Placer les foyers de résistance, c'est montrer que la conquête n'a jamais été acceptée passivement. La carte devient alors un \cle{document d'histoire} à part entière, au même titre qu'un texte ou une photographie.

\begin{methode}[Construire une carte historique : les six règles d'or]
Une carte historique conforme aux règles de la cartographie comporte obligatoirement six éléments :
\begin{enumerate}
\item \textbf{Un fond de carte} : le contour du territoire, tracé proprement, sans détail inutile.
\item \textbf{Un titre} : précis, daté, situé en haut (exemple : « L'extension du Kamerun, 1884-1911 »).
\item \textbf{Une légende} : chaque couleur et chaque symbole utilisé doit être expliqué dans un encadré.
\item \textbf{Une échelle} : graphique (segment gradué) ou chiffrée (1/\num{5000000}), pour mesurer les distances.
\item \textbf{L'orientation} : une flèche indiquant le nord.
\item \textbf{Les sources} : d'où viennent les informations (manuel, atlas, archives).
\end{enumerate}
Sans ces six éléments, votre production n'est pas une carte : c'est un coloriage.
\end{methode}

\begin{attention}[L'erreur qui coûte des points]
L'erreur la plus fréquente au Probatoire et au Baccalauréat technique est d'\textbf{oublier la légende ou l'échelle}. Une carte sans légende est \emph{illisible} : le correcteur ne peut pas deviner ce que signifient vos couleurs. Une carte sans titre ni date est \emph{inutilisable} : on ne sait pas quelle période elle représente. Vérifiez toujours, avant de rendre votre copie : titre, légende, échelle, nord, sources.
\end{attention}

\courssub{TP 1 : tracer les frontières du Kamerun (1884, 1901, 1911)}

\textbf{Consigne.} À partir du fond de carte muet fourni (reproduit ci-dessous), tracez :
\begin{itemize}
\item en \textbf{bleu} : le territoire du protectorat de 1884 (la côte, autour de Douala, après le traité germano-douala du 12 juillet 1884) ;
\item en \textbf{vert} : l'extension de 1901 (la conquête de l'intérieur : Yaoundé, l'Adamaoua, le Nord) ;
\item en \textbf{orange} : le \cle{Neukamerun} de 1911 (les territoires cédés par la France après la crise d'Agadir, accord du 4 novembre 1911).
\end{itemize}

\begin{popfigure}
\begin{center}
\begin{tikzpicture}[scale=1.05, every node/.style={font=\small}]
% Fond de carte muet : contour simplifié du Cameroun
\draw[thick, popdark, fill=popcream] (0,0) -- (0.8,-0.6) -- (1.6,-0.9) -- (2.6,-0.7) -- (3.4,-0.2) -- (4.2,0.4) -- (4.8,1.2) -- (5.0,2.2) -- (4.6,3.2) -- (4.9,4.2) -- (4.4,5.4) -- (3.6,6.4) -- (3.2,7.4) -- (2.4,8.2) -- (1.8,8.0) -- (1.4,7.0) -- (0.9,6.2) -- (0.2,5.6) -- (-0.4,4.6) -- (-0.6,3.6) -- (-0.9,2.6) -- (-0.8,1.4) -- (-0.3,0.5) -- cycle;
% Côte atlantique
\draw[popblue, very thick] (-0.3,0.5) -- (-0.8,1.4) -- (-0.9,2.6);
\node[popblue, rotate=75] at (-1.35,1.6) {Océan Atlantique};
% Flèche du nord
\draw[-{Stealth}, thick, popdark] (5.8,7.2) -- (5.8,8.2);
\node[above, popdark] at (5.8,8.2) {N};
% Échelle
\draw[thick, popdark] (5.4,0.2) -- (6.4,0.2);
\draw[thick, popdark] (5.4,0.1) -- (5.4,0.3);
\draw[thick, popdark] (6.4,0.1) -- (6.4,0.3);
\node[below, popdark] at (5.9,0.05) {200 km};
% Quelques repères muets
\fill[popmuted] (-0.55,1.9) circle (1.5pt);
\fill[popmuted] (1.3,2.6) circle (1.5pt);
\fill[popmuted] (2.2,3.4) circle (1.5pt);
\end{tikzpicture}
\end{center}
\legende{Modèle de fond de carte muet à reproduire : contour simplifié du Cameroun, côte atlantique, flèche du nord et échelle graphique. Les points gris sont des repères (villes) à identifier. Croquis schématique, sans précision cartographique.}
\end{popfigure}

\textbf{Travail à réaliser.} Reproduisez ce fond de carte sur une feuille, puis tracez les trois frontières. Observez le raisonnement géographique : le Kamerun de 1884 n'est qu'une bande côtière ; celui de 1901 s'enfonce vers le nord ; le Neukamerun de 1911 s'étire vers l'est, jusqu'au Congo et au Tchad. La carte prouve ainsi ce que le texte affirme : le territoire est une \emph{construction} progressive.

\courssub{TP 2 : capitales, chemins de fer et foyers de résistance}

\textbf{Consigne.} Sur la même carte, placez :
\begin{itemize}
\item les capitales successives : \cle{Douala} (1884), \cle{Buéa} (1901), \cle{Yaoundé} (1909), symbolisées par une étoile ;
\item les lignes de chemin de fer : la ligne du Nord (Douala--Nkongsamba) et la ligne du Centre (Douala--Yaoundé), symbolisées par un trait noir interrompu de croix ;
\item les principaux foyers de résistance : les Douala (résistance diplomatique puis armée), les Bulu du Sud (révolte de 1895), les Bamiléké et les Bali à l'ouest, les peuples du Nord (guerres de conquête), symbolisés par un triangle rouge.
\end{itemize}

\begin{exemplebox}[Exemple de légende normalisée]
Voici comment organiser votre légende dans un encadré, en bas à droite de la carte :
\begin{itemize}
\item \textcolor{popblue}{\rule{0.8cm}{2.5pt}} : frontière du protectorat (1884)
\item \textcolor{popgreen}{\rule{0.8cm}{2.5pt}} : extension de 1901
\item \textcolor{poporange}{\rule{0.8cm}{2.5pt}} : Neukamerun (1911)
\item $\bigstar$ : capitale \quad $\blacktriangle$ : foyer de résistance \quad \rule[1pt]{0.25cm}{1pt}\,\rule[1pt]{0.1cm}{1pt}\,\rule[1pt]{0.25cm}{1pt} : chemin de fer
\end{itemize}
Chaque symbole de la carte doit apparaître dans la légende, et rien dans la légende ne doit être absent de la carte.
\end{exemplebox}

\courssub{TP 3 : titre, légende et conformité cartographique}

Rédigez maintenant le titre définitif et la légende complète de votre carte. Un bon titre de carte historique répond à trois questions : \emph{quoi ?} (le sujet), \emph{où ?} (l'espace), \emph{quand ?} (la période).

\begin{exoresolu}[Rédiger un titre conforme]
Énoncé : un élève propose le titre « Carte du Cameroun ». Corrigez ce titre pour le rendre conforme aux règles de la cartographie historique.
\tcblower
Solution : le titre « Carte du Cameroun » est inacceptable pour trois raisons : il ne précise ni le sujet (frontières ? résistances ? économie ?), ni la période, et il utilise le nom actuel « Cameroun » pour une entité qui s'appelait alors « Kamerun ». Un titre conforme serait : \textbf{« L'extension territoriale du Kamerun et les résistances à la conquête allemande, 1884-1911 »}. Ce titre annonce le sujet (extension + résistances), l'espace (le Kamerun) et la période (1884-1911).
\end{exoresolu}

\courssub{TP 4 : diagramme en barres des exportations du Kamerun}

La carte montre l'espace ; le graphique montre les quantités. L'économie du Kamerun reposait sur l'\cle{exportation} de matières premières vers l'Allemagne : cacao, caoutchouc, huile de palme. Voici des données simplifiées (en milliers de tonnes, ordres de grandeur) :

\begin{center}
\begin{tabularx}{\linewidth}{|l|X|X|X|}
\hline
\rowcolor{popblueL} \textbf{Produit} & \textbf{1895} & \textbf{1905} & \textbf{1913} \\
\hline
Cacao & 0,2 & 1,5 & 5,0 \\
\hline
Caoutchouc & 0,8 & 1,2 & 0,6 \\
\hline
Huile de palme & 1,0 & 2,0 & 3,5 \\
\hline
\end{tabularx}
\end{center}

\begin{methode}[Construire un diagramme en barres]
\begin{enumerate}
\item Tracez deux axes : l'axe horizontal porte les catégories (les années), l'axe vertical porte les valeurs.
\item \textbf{Graduez} l'axe vertical régulièrement (0, 1, 2, 3, 4, 5) et indiquez l'\textbf{unité} (milliers de tonnes).
\item Donnez un \textbf{titre} au graphique.
\item Dessinez des barres de même largeur, dont la \textbf{hauteur est proportionnelle} à la valeur : une valeur double donne une barre deux fois plus haute.
\item Ajoutez une légende si plusieurs séries sont représentées (une couleur par produit).
\end{enumerate}
\end{methode}

\begin{popfigure}
\begin{center}
\begin{tikzpicture}
\begin{axis}[
ybar, bar width=10pt,
width=0.95\linewidth, height=6.5cm,
ylabel={Milliers de tonnes},
xlabel={Années},
symbolic x coords={1895,1905,1913},
xtick=data,
ymin=0, ymax=6,
legend style={at={(0.02,0.97)}, anchor=north west, font=\small},
nodes near coords, every node near coord/.append style={font=\scriptsize},
axis line style={popdark},
]
\addplot[fill=popgold, draw=popdark] coordinates {(1895,0.2) (1905,1.5) (1913,5.0)};
\addplot[fill=popgreen, draw=popdark] coordinates {(1895,0.8) (1905,1.2) (1913,0.6)};
\addplot[fill=poporange, draw=popdark] coordinates {(1895,1.0) (1905,2.0) (1913,3.5)};
\legend{Cacao, Caoutchouc, Huile de palme}
\end{axis}
\end{tikzpicture}
\end{center}
\legende{Exportations du Kamerun vers l'Allemagne (ordres de grandeur simplifiés, en milliers de tonnes). On observe l'essor spectaculaire du cacao après 1905 et le déclin du caoutchouc, concurrencé par les plantations d'Asie.}
\end{popfigure}

\textbf{Lecture du graphique.} Le cacao est multiplié par 25 entre 1895 et 1913 : à la veille de la Première Guerre mondiale, le Kamerun est devenu l'un des tout premiers exportateurs mondiaux de cacao. Le caoutchouc recule après 1905 : la concurrence des plantations d'Asie ruine la collecte du caoutchouc sauvage. L'huile de palme progresse régulièrement. Conclusion : l'économie coloniale allemande est une économie de \emph{plantation et d'extraction}, tournée vers l'extérieur, qui profite d'abord à la métropole.

\courssub{TP 5 : croquis de synthèse « Le Kamerun allemand : un territoire construit et contesté »}

Le croquis de synthèse est l'exercice le plus difficile : il faut \emph{superposer} sur une seule carte les trois dimensions de la leçon — la construction territoriale, l'organisation économique, les résistances — sans surcharger le dessin.

\begin{methode}[Réussir un croquis de synthèse]
\begin{enumerate}
\item Choisissez \textbf{trois informations maximum} par thème : trois frontières, trois villes, deux lignes de chemin de fer, trois foyers de résistance.
\item Hiérarchisez : les éléments les plus importants (frontières) en traits forts, les éléments secondaires en symboles discrets.
\item Rédigez une légende \textbf{organisée par thème} : I. Territoire ; II. Économie ; III. Résistances.
\item Terminez par titre, nord, échelle et sources.
\end{enumerate}
\end{methode}

\begin{popfigure}
\begin{center}
\begin{tikzpicture}[scale=0.95, every node/.style={font=\scriptsize}]
% Contour
\draw[thick, popdark, fill=popcream] (0,0) -- (0.8,-0.6) -- (1.6,-0.9) -- (2.6,-0.7) -- (3.4,-0.2) -- (4.2,0.4) -- (4.8,1.2) -- (5.0,2.2) -- (4.6,3.2) -- (4.9,4.2) -- (4.4,5.4) -- (3.6,6.4) -- (3.2,7.4) -- (2.4,8.2) -- (1.8,8.0) -- (1.4,7.0) -- (0.9,6.2) -- (0.2,5.6) -- (-0.4,4.6) -- (-0.6,3.6) -- (-0.9,2.6) -- (-0.8,1.4) -- (-0.3,0.5) -- cycle;
% Frontière 1884 (zone côtière)
\draw[popblue, very thick] (-0.75,2.4) -- (0.2,2.2) -- (0.9,2.0) -- (1.4,1.6);
\node[popblue, right] at (1.5,1.7) {1884};
% Frontière 1901
\draw[popgreen, very thick] (-0.5,4.4) -- (0.6,4.0) -- (1.8,3.9) -- (3.0,4.1) -- (4.0,4.4);
\node[popgreen, right] at (3.4,4.7) {1901};
% Neukamerun 1911 (extension est)
\draw[poporange, very thick, dashed] (2.6,-0.7) -- (3.4,-0.2) -- (4.2,0.4) -- (4.8,1.2) -- (5.0,2.2) -- (4.6,3.2);
\node[poporange, right] at (4.7,0.6) {Neukamerun};
\node[poporange, right] at (4.7,0.2) {1911};
% Villes
\node[popdark] at (-0.55,1.9) {$\bigstar$};
\node[left, popdark] at (-0.65,1.9) {Douala};
\node[popdark] at (-0.2,2.6) {$\bigstar$};
\node[left, popdark] at (-0.3,2.75) {Buéa};
\node[popdark] at (2.2,3.0) {$\bigstar$};
\node[above, popdark] at (2.2,3.1) {Yaoundé};
% Chemin de fer Douala - Yaoundé
\draw[popdark, thick, dash pattern=on 3pt off 2pt on 1pt off 2pt] (-0.55,1.9) -- (0.8,2.3) -- (2.2,3.0);
% Foyers de résistance
\node[poppink] at (0.3,1.5) {$\blacktriangle$};
\node[poppink] at (1.6,2.2) {$\blacktriangle$};
\node[below, poppink] at (1.6,2.05) {Bulu};
\node[poppink] at (0.9,3.6) {$\blacktriangle$};
\node[above, poppink] at (0.9,3.7) {Ouest};
\node[poppink] at (2.6,6.6) {$\blacktriangle$};
\node[above, poppink] at (2.6,6.75) {Nord};
% Nord et échelle
\draw[-{Stealth}, thick, popdark] (5.8,7.2) -- (5.8,8.2);
\node[above, popdark] at (5.8,8.2) {N};
\draw[thick, popdark] (5.4,0.2) -- (6.4,0.2);
\node[below, popdark] at (5.9,0.05) {200 km};
\end{tikzpicture}
\end{center}
\legende{Exemple de croquis de synthèse : « Le Kamerun allemand, un territoire construit et contesté (1884-1916) ». 1 : frontières de 1884 (bleu), 1901 (vert), Neukamerun 1911 (orange) ; 2 : capitales ($\bigstar$) ; 3 : chemin de fer Douala--Yaoundé ; 4 : foyers de résistance ($\blacktriangle$). Croquis schématique.}
\end{popfigure}

\begin{attention}[Pièges du croquis de synthèse]
\begin{itemize}
\item Ne surchargez pas : dix symboles illisibles valent moins que cinq symboles clairs.
\item N'inventez pas de précision : un croquis est \emph{schématique}, il montre des positions relatives, pas des coordonnées exactes.
\item N'oubliez pas que le croquis doit \emph{prouver la thèse du titre} : ici, que le territoire a été à la fois construit (frontières, villes, voies ferrées) et contesté (triangles de résistance répartis sur tout l'espace).
\end{itemize}
\end{attention}

\courssub{Restitution : présenter et justifier sa carte}

Chaque groupe dispose de \textbf{3 minutes} pour présenter sa carte à la classe. La présentation suit un plan en trois temps :
\begin{enumerate}
\item \textbf{Décrire} : « Notre carte montre... » (titre, période, éléments représentés).
\item \textbf{Expliquer} : « Nous avons choisi ces couleurs et ces symboles parce que... » (justifier chaque choix de représentation).
\item \textbf{Conclure} : « Cette carte démontre que... » (formuler la leçon d'histoire que la carte rend visible).
\end{enumerate}

\begin{exercice}[Pour s'entraîner]
\begin{enumerate}
\item Pourquoi le nom « Kamerun » doit-il être utilisé sur une carte de la période 1884-1916 plutôt que « Cameroun » ?
\item Sur le diagramme en barres, calculez le taux de variation des exportations de cacao entre 1895 et 1913.
\item Citez les six éléments obligatoires d'une carte historique et vérifiez leur présence sur votre croquis de synthèse.
\end{enumerate}
\end{exercice}

\begin{corrige}[Exercice]
\begin{enumerate}
\item Parce que « Kamerun » est le nom officiel de l'entité territoriale créée par l'Allemagne en 1884 ; « Cameroun » est le nom postérieur. Une carte historique doit respecter le vocabulaire de l'époque qu'elle représente.
\item Taux de variation $= \dfrac{5{,}0 - 0{,}2}{0{,}2} \times 100 = \num{2400}\,\%$. Les exportations de cacao ont été multipliées par 25.
\item Fond de carte, titre, légende, échelle, orientation (nord), sources.
\end{enumerate}
\end{corrige}

\begin{savaistu}[Les cartes allemandes du Kamerun, des archives vivantes]
Les cartes coloniales dressées par les officiers et les géographes allemands entre 1884 et 1916 sont aujourd'hui conservées dans des archives en Allemagne, notamment aux Archives fédérales de Berlin. Les historiens camerounais les consultent pour reconstituer les frontières, les routes et les villages de l'époque. Ces cartes, autrefois instruments de la conquête, sont devenues des outils de recherche au service de la mémoire du Cameroun : preuve qu'un document change de sens selon celui qui l'utilise.
\end{savaistu}

\begin{aretenir}[L'essentiel de la séance]
Une carte historique est un raisonnement, pas un dessin : elle exige fond de carte, titre daté, légende, échelle, nord et sources. Le Kamerun s'est construit en trois étapes (1884 : la côte ; 1901 : l'intérieur ; 1911 : le Neukamerun) et cette construction s'est accompagnée de résistances sur tout le territoire. Le diagramme en barres complète la carte en quantifiant l'économie d'exportation (essor du cacao, déclin du caoutchouc). Savoir construire et présenter ces documents, c'est savoir faire de l'histoire.
\end{aretenir}

\newpage

\coursec{Activité d'intégration}

\courssub{Objectifs de l'activité}
Mobiliser les connaissances sur l'administration coloniale allemande (organisation territoriale, politique foncière, travail forcé) et les réactions camerounaises (collaboration, résistance armée, résistance juridique et politique) pour construire une argumentation historique rigoureuse. L'activité vise à faire distinguer les registres de l'action (juridique, diplomatique, militaire), à croiser des sources de nature hétérogène (traité, carte, archives judiciaires, données économiques) et à rédiger une conclusion argumentée justifiant l'entrée de Rudolf Duala Manga Bell dans le panthéon des héros nationaux camerounais.

\courssub{Situation}
\emph{Contexte historico-juridique.} Nous sommes en août 1914 à Douala. La guerre européenne vient d'éclater, mais le \emph{Kamerun} est encore sous administration allemande. Le tribunal militaire allemand siège pour juger Rudolf Duala Manga Bell, chef supérieur des Bell, et Martin-Paul Samba (originaire du pays Ewondo-Bulu, éduqué et formé militairement en Allemagne), accusés de \emph{Hochverrat} (haute trahison) pour avoir tenté d'entrer en contact avec les puissances alliées (France, Grande-Bretagne) afin de dénoncer les expropriations massives et de demander la protection de leurs peuples.

\emph{Le dossier documentaire fourni à chaque groupe (extraits).}
\begin{enumerate}
    \item \textbf{Traité germano-douala du 12 juillet 1884 :} signé par les rois et chefs Duala (Bell, Akwa, Deido) avec les représentants des firmes allemandes et le consul impérial. Il place le pays sous protection allemande tout en stipulant que les droits de propriété foncière, les coutumes et le monopole commercial des indigènes seront respectés.
    \item \textbf{Ordonnance foncière de 1913 (extrait) :} déclaration d'utilité publique pour la zone de Bonanjo (terres Bell) et Bali (terres Akwa) en vue de la création d'un quartier européen strictement ségrégué. Les populations doivent être déplacées vers \emph{Neu-Bell} (New Bell), zone marécageuse et insalubre, sans indemnisation équitable.
    \item \textbf{Carte schématique des expropriations de Bonanjo (1913-1914) :} superficie concernée : environ 350 hectares. Population déplacée : environ 3 000 personnes. Distance Bonanjo $\to$ New Bell : 4 km. Forte surmortalité observée à New Bell (1913-1914), liée au paludisme et à la dysenterie.
    \item \textbf{Correspondance de Rudolf Duala Manga Bell (1913-1914) :} pétitions au Reichstag (parlement allemand) et au \emph{Reichskolonialamt} (ministère des Colonies) : « Nous ne demandons pas l'indépendance, nous demandons le respect du traité de 1884. Si l'Allemagne ne nous écoute pas, nous nous tournerons vers d'autres puissances pour obtenir justice. »
    \item \textbf{Témoignage d'un administrateur colonial (1914) :} « Manga Bell abuse de son influence. Son éducation allemande le rend dangereux. Il faut briser le pouvoir des chefs traditionnels pour moderniser la colonie. »
    \item \textbf{Données économiques des plantations (1913) :} les grandes compagnies (Woermann, Jantzen \& Thormählen, sociétés de plantation du Sud-Kamerun) contrôlent des dizaines de milliers d'hectares de cacao, palmiers et bananes, exploités grâce au recrutement forcé de la main-d'œuvre (corvées, portage), source de profits considérables.
    \item \textbf{Acte d'accusation du tribunal militaire (août 1914) :} « Conspiration avec l'ennemi en temps de guerre. Peine requise : mort par pendaison. »
\end{enumerate}

\begin{attention}[Vigilance documentaire]
Les documents 5, 6 et 7 sont des sources \emph{allemandes} : elles expriment le point de vue de l'administration coloniale et doivent être croisées avec les documents 1 et 4 (sources camerounaises ou contractuelles). En histoire, on ne prend jamais un document pour la vérité : on l'interroge (qui parle ? à qui ? dans quel but ?).
\end{attention}

\emph{Mise en scène.} La classe est divisée en 4 collèges (groupes de 5-6 élèves) :
\begin{itemize}
    \item \textbf{Collège A -- Procureurs impériaux :} doivent prouver la trahison en s'appuyant sur la correspondance avec l'ennemi (doc. 4) et la menace pour l'ordre colonial (doc. 5, 7).
    \item \textbf{Collège B -- Avocats de la défense :} doivent démontrer la légalité de la démarche (respect du traité de 1884, doc. 1), le caractère abusif de l'expropriation (doc. 2, 3) et l'absence d'intention séditieuse avant 1914.
    \item \textbf{Collège C -- Témoins duala et bulu :} représentent les populations déplacées (doc. 3), les travailleurs forcés des plantations (doc. 6) et les chefs traditionnels déchus. Ils apportent la preuve du préjudice subi.
    \item \textbf{Collège D -- Historiens (jury) :} doivent synthétiser les arguments, qualifier historiquement les faits (résistance juridique ou trahison) et rendre un verdict motivé.
\end{itemize}

\courssub{Consignes}
\begin{enumerate}
    \item \textbf{Analyse documentaire (20 min) :} chaque collège identifie dans le dossier les 3 pièces les plus favorables à sa thèse et note les contradictions ou silences des pièces adverses. Compléter un tableau d'analyse : \emph{Nature du doc / Auteur / Date / Information clé / Fiabilité / Utilisation pour ma thèse}.
    \item \textbf{Rédaction des plaidoiries (30 min) :}
    \begin{itemize}
        \item Procureurs : rédiger le réquisitoire (max 300 mots) structuré en 3 chefs d'accusation.
        \item Avocats : rédiger la plaidoirie (max 300 mots) structurée en 3 moyens de défense.
        \item Témoins : rédiger 3 dépositions courtes (50 mots chacune) incarnées (je, nous) appuyées sur les informations du dossier.
        \item Historiens : préparer une grille d'évaluation (légalité, proportionnalité, intentionnalité, légitimité).
    \end{itemize}
    \item \textbf{Simulation du procès (30 min) :} ordre : Procureurs (5 min) $\to$ Avocats (5 min) $\to$ Témoins (3 $\times$ 2 min) $\to$ Questions des Historiens (10 min) $\to$ Délibération publique des Historiens (5 min) $\to$ Verdict.
    \item \textbf{Production individuelle écrite (20 min) :} chaque élève (tous collèges confondus) rédige un paragraphe argumenté (150-200 mots) répondant à la question : \og Pourquoi Rudolf Duala Manga Bell est-il considéré aujourd'hui comme un héros national camerounais ? \fg La réponse doit mobiliser \textbf{au moins 3 documents du dossier} et utiliser les concepts : \cle{résistance juridique}, \cle{spoliation foncière}, \cle{traité de 1884}, \cle{mémoire collective}.
\end{enumerate}

\courssub{Ressources à mobiliser}
\begin{itemize}
    \item \textbf{Notions juridiques :} traité international, \emph{protectorat} (le traité de 1884 instaure un protectorat, non une annexion : la souveraineté interne des chefs est en principe préservée), propriété foncière coutumière face à l'appropriation écrite coloniale, haute trahison (droit pénal allemand de 1914).
    \item \textbf{Concepts historiques :} administration coloniale allemande s'appuyant sur les chefs traditionnels (intermédiaires rémunérés et contrôlés) ; rupture de 1913 (l'expropriation de Bonanjo marque la fin de l'alliance passée en 1884). \cle{Résistance} : armée (soulèvements du Sud et de l'Ouest, campagnes de « pacification »), politique et juridique (Duala Manga Bell, Martin-Paul Samba), culturelle (églises, écoles).
    \item \textbf{Repères chronologiques :} 1884 (traité germano-douala), 1885-1907 (conquête et « pacification » militaire de l'intérieur), 1911 (élargissement du Kamerun par l'accord de Marrakech-Fès : le \emph{Neukamerun}), 1913 (ordonnance foncière), 8 août 1914 (exécution de Manga Bell et Samba), 1914-1916 (campagne du Kamerun, fin de la domination allemande).
    \item \textbf{Méthode de commentaire de document :} Identifier (nature, auteur, date, contexte), Analyser (vocabulaire, arguments, chiffres), Confronter (croiser les sources), Conclure (portée historique).
\end{itemize}

\begin{aretenir}[Repères essentiels]
Le traité germano-douala du 12 juillet 1884 fonde le protectorat allemand du \emph{Kamerun}. L'administration allemande organise la colonie autour de Douala, développe les plantations et les voies de communication (chemin de fer du Nord et du Centre) grâce au travail forcé. Face aux exactions (expropriations, corvées, impôts), les Camerounais réagissent par la résistance armée, mais aussi par la résistance juridique et politique, dont Rudolf Duala Manga Bell est la figure emblématique : pendu le 8 août 1914 avec Martin-Paul Samba, il est devenu un héros national.
\end{aretenir}

\courssub{Démarche et corrigé commenté}
\begin{exoresolu}[Corrigé type et analyse des arguments attendus]
\textbf{1. Analyse documentaire type (extrait pour le Collège B -- Défense).}
\begin{center}
\begin{tabularx}{\linewidth}{|l|X|X|X|}\hline
\textbf{Doc} & \textbf{Information clé} & \textbf{Fiabilité / Limite} & \textbf{Utilisation défense} \\ \hline
1. Traité 1884 & Respect de la propriété foncière et des coutumes. & Source officielle, signée par les rois Bell, Akwa, Deido. & Preuve du contrat bilatéral violé unilatéralement par l'Allemagne. \\ \hline
2. Ordonnance 1913 & Expropriation pour créer un quartier européen ségrégué. & Acte administratif unilatéral, postérieur au traité. & Preuve de la violation : but ségrégationniste, pas utilité publique générale. \\ \hline
3. Carte Bonanjo & 350 ha, 3 000 déplacés, surmortalité à New Bell. & Données administratives allemandes (sous-estimation probable). & Preuve du préjudice matériel et humain disproportionné. \\ \hline
4. Lettres Manga Bell & Pétitions au Reichstag, puis menace d'appel à d'autres puissances. & Source primaire, démarche d'abord pacifique et légale. & Preuve de l'épuisement des voies de droit internes avant la rupture. \\ \hline
6. Données plantations & Vastes concessions, main-d'œuvre forcée, gros profits. & Statistiques des compagnies, à croiser. & Contexte : l'expropriation sert l'économie de plantation, pas l'intérêt général. \\ \hline
\end{tabularx}
\end{center}

\textbf{2. Réquisitoire type (Collège A -- Procureurs).}
\og Messieurs les juges, l'accusé Rudolf Duala Manga Bell a commis le crime de haute trahison (\emph{Hochverrat}) en trois points :
1. \textbf{Correspondance avec l'ennemi :} le document 4 établit qu'en 1914 l'accusé a cherché à entrer en contact avec les puissances ennemies (France, Grande-Bretagne) pour solliciter leur intervention. En temps de guerre, cela constitue une intelligence avec l'ennemi passible de mort (doc. 7).
2. \textbf{Incitation à la sédition :} en tant que chef supérieur reconnu et rémunéré par l'administration, il a utilisé son autorité pour mobiliser les populations contre l'ordonnance foncière de 1913 (doc. 2), menaçant l'ordre public et l'économie de la colonie (doc. 6).
3. \textbf{Complicité avec Martin-Paul Samba :} ce dernier, formé militairement en Allemagne, a préparé un soulèvement. La convergence des démarches prouve une conspiration ourdie de longue date contre l'œuvre allemande. \fg

\textbf{3. Plaidoirie type (Collège B -- Avocats).}
\og Messieurs les juges, mon client n'est pas un traître, c'est un justiciable qui a épuisé toutes les voies de droit.
1. \textbf{Violation du traité fondateur :} le traité de 1884 (doc. 1) garantit les droits de propriété des indigènes. L'ordonnance de 1913 (doc. 2) viole ce traité. Nul ne peut être condamné pour avoir dénoncé l'illégalité de l'État qui le juge.
2. \textbf{Absence d'intention séditieuse :} le document 4 montre une démarche \emph{hiérarchique et légale} : pétitions au Reichstag et au ministère des Colonies pendant plus d'un an. La mention des autres puissances n'est qu'un \emph{ultimatum juridique}, non un acte de guerre.
3. \textbf{Disproportion de la répression :} les documents 3 et 6 prouvent que l'expropriation sert des intérêts privés (les compagnies de plantation) au prix de vies humaines (surmortalité à New Bell). Résister à une spoliation n'est pas de la trahison, c'est l'exercice d'un droit. \fg

\textbf{4. Dépositions types (Collège C -- Témoins).}
\begin{itemize}
    \item \textit{Témoin 1 (Chef de quartier de Bonanjo) :} « Nous habitions nos terres depuis des générations. En 1913, on nous a ordonné de partir vers New Bell, un marais. Beaucoup des nôtres y sont morts du paludisme (doc. 3). On nous a chassés pour faire place aux Allemands. »
    \item \textit{Témoin 2 (Ancien travailleur forcé des plantations) :} « Nous étions des milliers chaque année à porter les charges et à défricher sous le fouet (doc. 6). Le chef Manga Bell a refusé de nous livrer comme du bétail. Il a dit : "Le traité dit que nous sommes libres." Pour cela, on le pend. »
    \item \textit{Témoin 3 (Notable allié de Martin-Paul Samba) :} « Samba est revenu d'Allemagne avec des idées de liberté. Il a vu que le Blanc n'est pas un dieu. Il nous a dit : "Le traité de 1884 est notre arme." Aujourd'hui, on juge Manga Bell, mais c'est tout le peuple camerounais qu'on juge. »
\end{itemize}

\textbf{5. Délibération et verdict des Historiens (Collège D).}
\og \textbf{Verdict historique :} le tribunal militaire allemand a condamné à mort Rudolf Duala Manga Bell et Martin-Paul Samba, exécutés par pendaison le 8 août 1914 à Douala. \textbf{Au regard de l'histoire, ce procès est une mascarade judiciaire.}
\textbf{Motifs :}
\begin{enumerate}
    \item \textbf{Illégalité de la cause :} la spoliation foncière de 1913 viole le traité de 1884 (doc. 1, 2). L'administration allemande rompt le pacte fondateur du protectorat.
    \item \textbf{Légitimité de la résistance juridique :} Manga Bell use des armes du droit (pétitions, parlement, presse) pendant plus d'un an (1913-1914). C'est une \cle{résistance juridique} exemplaire et non violente.
    \item \textbf{Instrumentalisation de la guerre :} l'accusation de \emph{Hochverrat} utilise le contexte de guerre européenne (août 1914) pour liquider un opposant gênant pour les intérêts des compagnies (doc. 6).
    \item \textbf{Héroïsation posthume :} le sacrifice consenti et la convergence des luttes fondent la \cle{mémoire collective} d'un héros national, symbole de la dignité africaine face à l'arbitraire colonial.
\end{enumerate}
\fg

\textbf{6. Paragraphe argumenté type (production individuelle attendue).}
\og Rudolf Duala Manga Bell est considéré comme un héros national camerounais parce qu'il a incarné, en 1913-1914, une \textbf{résistance juridique} lucide et courageuse face à la \textbf{spoliation foncière} organisée par l'administration allemande. Le traité germano-douala de 1884 (doc. 1), base juridique du protectorat, garantissait le respect des droits de propriété et des coutumes. Or, l'ordonnance de 1913 (doc. 2) et la carte des expropriations de Bonanjo (doc. 3) prouvent que l'État colonial a violé unilatéralement ce contrat pour créer un quartier européen ségrégué, déportant environ 3 000 personnes vers une zone insalubre, au profit des intérêts économiques des compagnies de plantation (doc. 6). Manga Bell a d'abord épuisé les voies légales (pétitions au Reichstag, doc. 4) avant d'évoquer un appel à d'autres puissances, non par trahison, mais par fidélité au droit. Son exécution le 8 août 1914 (doc. 7) a transformé ce notable éduqué en Allemagne en martyr de la dignité camerounaise. La \textbf{mémoire collective} retient qu'il a refusé la soumission à l'arbitraire, fondant ainsi une référence identitaire partagée par la nation camerounaise. \fg
\tcblower
\textbf{Commentaire pédagogique :} ce corrigé montre la progression attendue : de l'extraction d'informations (tableau) à la construction rhétorique (plaidoiries), puis à la synthèse historienne (verdict) et enfin à l'écriture argumentée individuelle. L'évaluation portera sur : (1) l'exactitude des références documentaires (citation explicite doc. 1, 3, 6...), (2) la pertinence des concepts (\cle{résistance juridique}, \cle{spoliation}, \cle{traité}), (3) la structure logique (thèse -- arguments -- exemples -- conclusion), (4) la nuance (distinction collaboration/résistance, contexte de guerre).
\end{exoresolu}

\courssub{Bilan de l'activité}
Cette activité d'intégration a permis de mettre en évidence plusieurs acquis majeurs du programme de Première technique :
\begin{itemize}
    \item \textbf{Maîtrise chronologique et causale :} les élèves situent la rupture de 1913 (expropriation de Bonanjo, fin de l'alliance de 1884) comme cause directe du procès de 1914.
    \item \textbf{Distinction des formes de résistance :} le procès met en lumière la spécificité de la \cle{résistance juridique et politique} (élites lettrées utilisant le droit et la pétition) par rapport à la \cle{résistance armée} étudiée en amont.
    \item \textbf{Esprit critique et traitement de l'information :} le croisement des sources (traité, ordonnance, carte, statistiques économiques, archives judiciaires) oblige à hiérarchiser la fiabilité : les faits matériels (doc. 3, 6) contredisent le discours officiel allemand (doc. 5, 7).
    \item \textbf{Argumentation écrite et orale :} la rédaction du paragraphe final évalue la capacité à \og rédiger et justifier une démarche \fg{} (savoir-faire officiel) en mobilisant le vocabulaire conceptuel (\cle{spoliation}, \cle{traité}, \cle{mémoire}) et non le simple récit.
    \item \textbf{Construction de la mémoire nationale :} l'activité ancre la figure de Manga Bell dans le temps long de l'histoire camerounaise : il n'est pas une victime passive, mais un acteur lucide dont le sacrifice fonde une référence identitaire partagée (rues, lycées et monuments portant son nom au Cameroun).
\end{itemize}
L'évaluation sommative portera sur la production écrite individuelle (notée sur 10 : 4 pts connaissances et mobilisation des documents, 3 pts structure et argumentation, 3 pts expression et concepts) et sur la participation orale lors du procès (grille d'observation : pertinence des interventions, écoute, vocabulaire historique).

\newpage

\coursec{Méthodes et exercices résolus}

\courssub{Méthodes et exercices résolus}

\begin{methode}[Analyser un traité de protectorat (Traité de juillet 1884)]
\textbf{Objectif :} Comprendre la mise en place de la souveraineté allemande par la voie juridique.
\begin{enumerate}
    \item \textbf{Identifier le document :} Nature (traité), date (12 juillet 1884), signataires (chefs douala : Bell, Akwa, Deido / Nachtigal pour l'Allemagne), lieu (Douala).
    \item \textbf{Repérer les clauses clés :} Protection allemande (art. 1), maintien des droits coutumiers des chefs (art. 2), monopole commercial allemand (art. 3), juridiction allemande sur les Européens (art. 4).
    \item \textbf{Analyser la portée juridique :} Distinction \emph{protectorat} (protection extérieure, autonomie interne théorique) vs \emph{colonie} (annexion pure). Noter l'ambiguïté : les chefs croient signer un traité d'amitié/commerce, l'Allemagne y voit une cession de souveraineté.
    \item \textbf{Contextualiser :} Conférence de Berlin (1884-1885), principe de l'effectivité, concurrence franco-allemande (traité de 1883 avec la France non ratifié).
    \item \textbf{Conclure :} Acte de naissance juridique du \cle{Kamerun}, base légale de l'administration coloniale.
\end{enumerate}
\end{methode}

\begin{exoresolu}[Analyse du traité germano-douala de juillet 1884]
\textbf{Énoncé :} À l'aide du texte du traité signé le 12 juillet 1884 entre le Dr Nachtigal et les rois Bell, Akwa et Deido, montrez que ce document fonde l'entité territoriale « Kamerun » tout en créant un malentendu sur la nature de l'autorité allemande.
\tcblower
\textbf{Solution détaillée :}
\begin{enumerate}
    \item \textbf{Identification :} Ce document est un \textbf{traité de protectorat} signé le \textbf{12 juillet 1884} à Douala. Les parties sont : pour l'Allemagne, le \textbf{Dr Gustav Nachtigal} (commissaire impérial) ; pour les populations locales, les \textbf{rois Bell (Ndumbé Lobé), Akwa (Dika Mpondo) et Deido (Jim Ekwalla)}. C'est un document diplomatique officiel, source primaire.
    \item \textbf{Analyse des clauses fondatrices :}
    \begin{itemize}
        \item \textbf{Article 1 :} « Les rois... placent leur pays sous la protection de l'Empire allemand. » $\rightarrow$ Acte juridique fondateur : naissance du \textbf{Kamerun} comme entité politique reconnue par le droit international (Conférence de Berlin).
        \item \textbf{Article 2 :} « Les droits des rois et chefs sur leurs terres et leurs sujets sont maintenus. » $\rightarrow$ Garantie formelle de l'autonomie interne (coutumes, justice locale, propriété foncière).
        \item \textbf{Article 3 :} « Le commerce est libre... mais l'Allemagne a le monopole des relations extérieures. » $\rightarrow$ Contrôle économique et diplomatique exclusif.
        \item \textbf{Article 4 :} « Les différends entre Allemands et indigènes sont jugés par les tribunaux allemands. » $\rightarrow$ Instauration de la \textbf{justice d'exception} (supériorité du droit allemand).
    \end{itemize}
    \item \textbf{Mise en évidence du malentendu (point crucial) :}
    \begin{itemize}
        \item \textbf{Côté camerounais :} Les chefs perçoivent un \emph{traité d'alliance} entre égaux : protection militaire contre ennemis intérieurs (Bakoko, Bassa) et extérieurs (Anglais/Français), maintien de leur autorité coutumière, libre commerce.
        \item \textbf{Côté allemand :} Nachtigal et le chancelier Bismarck y voient une \emph{cession de souveraineté} (Hoheitsrechte). L'expression « Schutzgebiet » (territoire de protection) masque une annexion déguisée. L'article 4 instaure l'extraterritorialité judiciaire, signe de soumission.
    \end{itemize}
    \item \textbf{Contexte et conséquences :} Ce traité s'inscrit dans la course aux colonies post-\textbf{Conférence de Berlin (nov. 1884 - fév. 1885)}. L'Allemagne obtient la \textbf{reconnaissance internationale} de ses droits sur le Kamerun (frontières fixées plus tard : 1885, 1894, 1911). Les résistances précoces (Douala 1884, Bassa 1890) naissent de la violation perçue de l'article 2 par l'administration coloniale.
\end{enumerate}
\textbf{Conclusion :} Le traité de 1884 est l'acte de naissance juridique du Kamerun. Il instaure un \textbf{protectorat de fait transformé en colonie de droit} par la pratique administrative, source d'un contentieux permanent entre l'autorité traditionnelle et l'administration allemande.
\end{exoresolu}

\begin{methode}[Expliquer l'organisation administrative : Règle directe vs Règle indirecte]
\textbf{Objectif :} Maîtriser le dualisme administratif allemand et ses acteurs.
\begin{enumerate}
    \item \textbf{Distinguuer les deux niveaux :}
    \begin{itemize}
        \item \textbf{Niveau central (Règle directe) :} Gouverneur (représente l'Empereur) $\rightarrow$ Siège : Buea (1901-1919) puis Douala (fin de guerre). Services : Secrétariat général, Justice, Finances, Troupes de protection (\cle{Schutztruppe}). Administration par des \textbf{fonctionnaires allemands}.
        \item \textbf{Niveau local (Règle indirecte) :} Subdivision en \textbf{Bezirke} (circonscriptions) dirigées par des \textbf{Bezirksamtmänner} (administrateurs de circonscription, Allemands). Au village : \textbf{Chefs coutumiers} maintenus comme \textbf{auxiliaires} (perception d'impôts, corvées, recrutement porteurs/soldats).
    \end{itemize}
    \item \textbf{Identifier les instruments de contrainte :} \cle{Impôt de capitation} (1891, payable en espèces $\rightarrow$ monétarisation forcée), \cle{Travaux forcés} (corvées routes/chemins de fer), \cle{Code pénal indigène} (1896, travail forcé légalisé, justice sommaire).
    \item \textbf{Nuancer selon les zones :} Zone côtière (règle plus directe, chefferies affaiblies/révoquées ex: Douala) vs Zone des Hauts-Plateaux/Adamawa (règle indirecte forte via chefferies structurées : Bamoun, Bafut, Lamidats).
    \item \textbf{Évaluer le bilan :} Modernisation infrastructurelle (chemins de fer, ports, routes) au prix d'une exploitation brutale (main-d'œuvre forcée, expropriations).
\end{enumerate}
\end{methode}

\begin{exoresolu}[L'administration du Kamerun : dualisme et contrainte]
\textbf{Énoncé :} « L'administration allemande au Kamerun a combiné une autorité centrale forte et une gestion locale déléguée aux chefs traditionnels. » À l'aide d'exemples précis, expliquez cette organisation et montrez comment elle s'est traduite par une contrainte économique et sociale sur les populations.
\tcblower
\textbf{Solution détaillée :}
\begin{enumerate}
    \item \textbf{L'architecture administrative (Le dualisme) :}
    \begin{itemize}
        \item \textbf{Sommet : Le Gouverneur.} Représentant direct de l'Empereur (Kaiser). Siège transféré de Douala à \textbf{Buea (1901)} pour raisons climatiques/sanitaires, puis retour à Douala (1914-1916). Ex : \textbf{Jesko von Puttkamer} (1895-1907), artisan de la « pacification » brutale ; \textbf{Theodor Seitz} (1907-1910), réformateur ; \textbf{Karl Ebermaier} (1910-1916), dernier gouverneur.
        \item \textbf{Relais : Les Bezirksamtmänner.} Administrateurs de circonscription (ex: \textbf{Bezirk Dschang}, \textbf{Bezirk Ngaoundéré}, \textbf{Bezirk Yaoundé}). Seuls Blancs sur place, pouvoir discrétionnaire immense (police, justice, impôt).
        \item \textbf{Base : Les chefs coutumiers (auxiliaires).} Système de la \textbf{Règle Indirecte} (\emph{Indirekte Herrschaft}). Les chefs perçoivent l'impôt, fournissent les porteurs/recrues, font appliquer les ordres. \textbf{Exemple type :} Le \textbf{Sultan Njoya} (Bamoun) ou les \textbf{Lamibé} de l'Adamaoua conservent leur cour mais deviennent des rouages de l'État colonial.
    \end{itemize}
    \item \textbf{Les instruments de la contrainte économique et sociale :}
    \begin{itemize}
        \item \textbf{L'impôt de capitation (1891) :} Impôt par tête (3 à 5 Marks). \textbf{Conséquence :} Obligation de vendre sa force de travail (plantations, chantiers) pour obtenir des espèces $\rightarrow$ \textbf{prolétarisation forcée} et monétarisation de l'économie de subsistance.
        \item \textbf{Le travail forcé (Zwangsarbeit) :} Légalisé par le \textbf{Code pénal indigène (1896)}. Utilisé pour : chemins de fer (Nord : 1901-1911 ; Centre : 1906-1911 ; Sud : 1910-1914), routes, port de Douala. Conditions : taux de mortalité élevé, pas de salaire, réquisition de porteurs (portage à dos d'homme).
        \item \textbf{L'expropriation foncière :} Décrets de 1896, 1899, 1901, 1913 (Douala). Déclaration des « terres vacantes et sans maître » (Konklusen) $\rightarrow$ attribution aux compagnies (Woermann, Jantzen \& Thormählen, Plantations sociétés) $\rightarrow$ déplacements de populations (ex: Bassa, Bakoko, Duala).
    \end{itemize}
    \item \textbf{Nuance spatiale :}
    \begin{itemize}
        \item \textbf{Côte (Douala, Victoria) :} Règle quasi-directe. Chefs déposés ou marginalisés (ex: \textbf{Rudolf Duala Manga Bell} 1914, chefs Bassa/Bakoko), tribunaux allemands compétents pour tous.
        \item \textbf{Intérieur (Grassfields, Adamaoua) :} Règle indirecte effective. Chefs puissants (Njoya, Galega de Bafut) collaborent pour conserver prérogatives internes.
    \end{itemize}
\end{enumerate}
\textbf{Conclusion :} L'organisation administrative allemande repose sur un \textbf{État fort au sommet (fonctionnaires, armée)} et une \textbf{délégation contrainte à la base (chefs auxiliaires)}. L'efficacité de cette machine vise l'extraction de la richesse (caoutchouc, cacao, main-d'œuvre) par la contrainte légale (impôt, code pénal, expropriation), provoquant résistances et collaborations.
\end{exoresolu}

\begin{methode}[Caractériser une résistance camerounaise : Typologie et analyse causale]
\textbf{Objectif :} Structurer l'étude d'un mouvement de résistance (causes, formes, acteurs, issue, bilan).
\begin{enumerate}
    \item \textbf{Classer le type de résistance :}
    \begin{itemize}
        \item \textbf{Armées / Politiques :} Conflit ouvert (Bafut, Bamoun, Douala 1914, Mora, Sanaga Maritime).
        \item \textbf{Passives / Quotidiennes :} Fuite (exode vers forêt/Nigeria britannique), refus d'impôt, sabotage cultures, ralentissement travail.
        \item \textbf{Religieuses / Culturelles :} Mouvements prophétiques (ex: Maria Ngoss, plutôt 1920s, ou cultes de résistance), refus école missionnaire, invention écriture (Shümom).
    \end{itemize}
    \item \textbf{Analyser les causes (grille CAS) :}
    \begin{itemize}
        \item \textbf{Conjoncturelles :} Exaction ponctuelle (viols, vols bétail, brutalité agent), recrutement forcé porteurs.
        \item \textbf{Structurelles :} Perte souveraineté, spoliation foncière, impôt de capitation, travail forcé, atteinte chefferie.
        \item \textbf{Identitaires :} Défense coutumes, religion traditionnelle, honneur du chef.
    \end{itemize}
    \item \textbf{Décrire les formes d'action :} Guérilla (embuscades, connaissance terrain), siège (Bafut), diplomatie (lettres au Kaiser, pétitions Manga Bell), alliance inter-ethnique (rare, ex: Bassa-Bakoko 1890s, Douala-Bassa-Intérieur 1913).
    \item \textbf{Expliquer l'échec militaire (facteurs) :} Supériorité technique allemande (Mitrailleuses \cle{Maxim}, artillerie, \cle{Schutztruppe} entraînées, médecine tropicale quinine), division africaine (rivalités inter-chefferies, collaborateurs), isolement diplomatique.
    \item \textbf{Évaluer les conséquences :} Répression (exécutions, déportations, villages brûlés), déportation chefs (Manga Bell, Ngoso Din, Abumbi I), mais aussi adaptations allemandes (réformes 1907-1911 : suppression travail forcé privé, tribunaux indigènes).
\end{enumerate}
\end{methode}

\begin{exoresolu}[La résistance du peuple Bafut (1901-1910) : analyse structurée]
\textbf{Énoncé :} En vous appuyant sur l'exemple de la résistance de Bafut (Grassfields), montrez que les résistances camerounaises face à l'administration allemande relevaient d'une logique de défense de la souveraineté politique et économique, mais ont échoué face à la supériorité technique et à la division.
\tcblower
\textbf{Solution détaillée :}
\begin{enumerate}
    \item \textbf{Contexte et causes (Grille CAS) :}
    \begin{itemize}
        \item \textbf{Structurelles :} Refus de la \textbf{perte de souveraineté} du Fon (roi) \textbf{Abumbi I}. Refus de l'\textbf{impôt de capitation} et du \textbf{travail forcé} (portage pour chemin de fer). Défense de l'intégrité territoriale face à l'expansion des plantations et postes militaires.
        \item \textbf{Conjoncturelles :} Arrogance du \textbf{Bezirksamtmann von Pavel} (poste de Bafut créé 1901). Exactions des soldats africains (Askaris) : pillages, viols. Humiliation du Fon (obligation de porter les bagages de l'administrateur).
        \item \textbf{Identitaires :} Défense de la \textbf{religion traditionnelle} et du rôle sacré du Fon (intermédiaire ancêtres) menacé par l'administration et les missions (Baptistes/Basel Mission).
    \end{itemize}
    \item \textbf{Acteurs et formes de la résistance :}
    \begin{itemize}
        \item \textbf{Acteur principal :} \textbf{Fon Abumbi I} (jeune, énergique, stratège). Soutenu par ses notables et prêtres traditionnels.
        \item \textbf{Forme :} \textbf{Guérilla de montagne} (1901-1910). Utilisation du relief (grottes, collines), embuscades sur les pistes, attaques nocturnes du poste allemand. \textbf{Diplomatie :} Tentatives d'alliance avec voisins (Bamenda, Kom, Bali) $\rightarrow$ \textbf{échec} (rivalités traditionnelles, certains chefs collaborent avec Allemands pour affaiblir Bafut).
    \end{itemize}
    \item \textbf{La riposte allemande et facteurs de l'échec :}
    \begin{itemize}
        \item \textbf{Moyens :} Expéditions punitives (1901, 1904, 1907, 1910). \textbf{Schutztruppe} + Askaris + porteurs réquisitionnés. Armement : \textbf{Fusils Mauser 98, Mitrailleuses Maxim, Canon de montagne 3,7 cm}.
        \item \textbf{Tactique :} « Terre brûlée » (destruction greniers, cases), prise d'otages (familles notables), blocus alimentaire.
        \item \textbf{Facteurs décisifs de l'échec Bafut :}
        \begin{enumerate}
            \item Supériorité de feu écrasante (Mitrailleuse vs Fusils à pierre/arc).
            \item Isolement diplomatique (pas d'alliance régionale solide).
            \item Usure démographique et économique (blocus, pertes humaines).
            \item Trahison / Renseignement fournis par chefferies rivales (Bali, Bamenda).
        \end{enumerate}
    \end{itemize}
    \item \textbf{Issues et bilan :}
    \begin{itemize}
        \item \textbf{Militaire :} Défaite finale 1910. Fon Abumbi I déposé, exilé à Douala (1910-1918). Village incendié.
        \item \textbf{Politique :} Installation d'un Fon collaborateur (Fon Achirimbi II, fils, élevé à l'école missionnaire).
        \item \textbf{Héritage :} Symbole de résistance héroïque. A poussé l'administration (Gouverneur \textbf{Theodor Seitz}, Secrétaire d'État \textbf{Wilhelm Solf}) à réformer : \textbf{Arrêté de 1911} limitant travail forcé, création \textbf{Tribunaux indigènes} (1911) pour apaiser.
    \end{itemize}
\end{enumerate}
\textbf{Conclusion :} La résistance de Bafut illustre le \textbf{paradoxe camerounais} : une volonté farouche de défendre l'autonomie politique (souveraineté du Fon) et économique (refus impôt/travail forcé), brisée par l'asymétrie technique (Maxim vs arcs) et l'incapacité à fédérer au-delà du cadre étatique précolonial. Elle a toutefois forcé l'Allemagne à « humaniser » son administration in extremis.
\end{exoresolu}

\begin{methode}[Comparer Collaboration et Résistance : Grille d'analyse comparative]
\textbf{Objectif :} Dépasser le manichéisme « collabos / résistants » pour analyser des stratégies d'acteurs rationnels.
\begin{enumerate}
    \item \textbf{Définir les concepts :}
    \begin{itemize}
        \item \textbf{Collaboration :} Acceptation active de l'autorité coloniale pour en tirer profit (maintien pouvoir, enrichissement, protection contre rivaux). $\neq$ Soumission passive.
        \item \textbf{Résistance :} Refus actif (armes, droit, fuite) de la domination. $\neq$ Révolte spontanée non politique.
        \item \textbf{Zones grises :} \textbf{Accommodation} (subir en attendant), \textbf{Instrumentalisation} (utiliser l'école/écriture pour défendre ses droits : ex. Rudolf Duala Manga Bell).
    \end{itemize}
    \item \textbf{Construire un tableau comparatif (Critères : Acteurs, Motivations, Moyens, Issue, Mémoire).}
    \item \textbf{Appliquer à deux figures emblématiques :}
    \begin{itemize}
        \item \textbf{Sultan Ibrahim Njoya (Bamoun) :} Collaboration \textbf{stratégique}. Accepte protectorat (1906) $\rightarrow$ Garde trône, territoire, développe écriture (Shümom), agriculture, artisanat. Utilise Allemands contre rivaux (Nso, Bafut). \textbf{Limite :} Déposé par Français (1924).
        \item \textbf{Rudolf Duala Manga Bell (Douala) :} Résistance \textbf{juridique et politique}. Éduqué en Allemagne, avocat. Utilise \textbf{droit allemand} (pétitions au Reichstag 1910, 1913) contre expropriation (décret 1913). Alliance multi-ethnique (Bassa, Bakoko, Ewondo). \textbf{Issue :} Pendu pour haute trahison (8 août 1914).
    \end{itemize}
    \item \textbf{Synthétiser :} La collaboration préserve l'institution chefferie à court terme mais l'affaiblit moralement. La résistance juridique échoue immédiatement mais fonde le nationalisme futur.
\end{enumerate}
\end{methode}

\begin{exoresolu}[Njoya et Manga Bell : deux stratégies face à la domination]
\textbf{Énoncé :} Comparez les stratégies du Sultan Njoya (Bamoun) et du Roi Rudolf Duala Manga Bell (Douala) face à l'administration allemande. En quoi ces deux figures illustrent-elles la complexité des rapports de pouvoir coloniaux ?
\tcblower
\textbf{Solution détaillée :}
\begin{center}
\begin{tabularx}{\linewidth}{|l|X|X|}
\hline
\textbf{Critères} & \textbf{Sultan Ibrahim Njoya (Bamoun - Hauts-Plateaux)} & \textbf{Roi Rudolf Duala Manga Bell (Douala - Côte)} \\
\hline
\textbf{Profil / Formation} & Héritier tradition, éducation cour, autodidacte génie inventif (écriture Shümom, moulin, horloge). & Éducation allemande (école mission, lycée en Allemagne), culture biculturelle, avocat. \\
\hline
\textbf{Date contact / Statut} & Protectorat accepté \textbf{1906} (traité). Reconnaissance mutuelle. & Traité 1884 (père). Héritier 1908. Statut \textbf{sujet allemand} (pas citoyen). \\
\hline
\textbf{Stratégie principale} & \textbf{Collaboration active / Instrumentalisation.} & \textbf{Résistance légale / Politique moderne.} \\
\hline
\textbf{Motivations profondes} & 1. Sauver le \textbf{royaume Bamoun} (menacé Nso, Allemands). 2. Moderniser (écriture, religion syncrétique, agriculture). 3. Garder autonomie interne. & 1. Défense \textbf{droits fonciers} (décret 1913 ségrégation spatiale). 2. Défense \textbf{dignité/statut} (égalité droits). 3. Unité \textbf{nationale naissante} (alliance côtiers/intérieur). \\
\hline
\textbf{Moyens d'action} & - Allié militaire fiable (fournit porteurs, soldats).\newline - Développement interne (écriture, écoles, agriculture, palais).\newline - Diplomatie : cadeaux à l'Empereur, lettres. & - \textbf{Pétitions au Reichstag} (1910, 1913) : argumentaire juridique (violation traité 1884 art. 2).\newline - Presse allemande (lettres journaux).\newline - \textbf{Alliance inter-ethnique} (Bassa, Bakoko, Ewondo, Bafia).\newline - Conseil d'avocats allemands. \\
\hline
\textbf{Réaction allemande} & \textbf{Confiance / Soutien.} Gouverneur Seitz le décore (Ordre de la Couronne, Aigle Rouge). Modèle « chef modèle ». & \textbf{Méfiance / Répression.} Perçu comme « dangereux agitateur ». Surveillance, interdiction réunions. \\
\hline
\textbf{Issue sous Allemands} & \textbf{Apogée :} Royaume puissant, modernisé, respecté (1914). & \textbf{Tragédie :} Arrêté (mai 1914), procès expéditif, \textbf{pendu 8 août 1914} (haute trahison, guerre déclarée). \\
\hline
\textbf{Héritage / Mémoire} & Héros culturel (écriture, palais, musée). Symbole \textbf{modernité endogène}. Limite : collaboration perçue comme soumission par certains nationalistes. & \textbf{Martyr nationaliste.} Père de l'indépendance (stade, boulevard, hymne). Symbole \textbf{résistance par le droit} et unité nationale. \\
\hline
\end{tabularx}
\end{center}
\textbf{Analyse synthétique :}
\begin{enumerate}
    \item \textbf{Contexte structurel diffère :} Njoya (intérieur, État centralisé fort, menace extérieure Nso) $\rightarrow$ Négociation de force. Manga Bell (côte, chefferie affaiblie, pression foncière intense, élite lettrée) $\rightarrow$ Confrontation juridique.
    \item \textbf{Complexité des rapports :} Njoya \textbf{utilise} le colonisateur pour renforcer son État (instrumentalisation). Manga Bell \textbf{utilise les outils du colonisateur} (droit, presse, parlement) pour le combattre (détournement d'arme).
    \item \textbf{L'illusion de la collaboration :} Njoya sauve son royaume \emph{provisoirement} (perdu sous Français 1924). La collaboration ne garantit pas la survie à long terme.
    \item \textbf{La modernité de la résistance :} Manga Bell invente le \textbf{nationalisme camerounais} : dépasse le cadre ethnique (alliance Douala/Bassa/Intérieur), utilise le droit international/allemand, vise l'égalité civile.
\end{enumerate}
\textbf{Conclusion :} Ces deux figures prouvent que la ligne « collaboration/résistance » est poreuse. Njoya résiste culturellement (écriture Shümom) en collaborant politiquement. Manga Bell collabore socialement (élite allemande) pour résister politiquement. Tous deux cherchent la \textbf{survie de leur peuple et de leur pouvoir} face à une machine coloniale qui, in fine, les broie ou les marginalise.
\end{exoresolu}

\begin{methode}[Réaliser un croquis historique : Le Kamerun (1884-1914/1916)]
\textbf{Objectif :} Produire une carte schématique rigoureuse (exigence Bac : croquis légendé, organisé, esthétique).
\begin{enumerate}
    \item \textbf{Préparer la base (au crayon) :} Contour simplifié du Kamerun 1911 (forme « triangle » + « jambe » vers lac Tchad). Repères : Méridien 10°-16° Est, Parallèles 2°-13° Nord. \textbf{Flèche Nord} (haut de page). \textbf{Échelle} (ex: 1 cm = 100 km).
    \item \textbf{Structurer la légende (3 à 4 items max) :} Organiser par \textbf{thèmes} (couleurs distinctes) et non par liste.
    \begin{itemize}
        \item \textbf{Item 1 : Administration / Villes.} Points proportionnels : Douala (capitale éco, port), Buea (capitale adm 1901-1919), Yaoundé (poste militaire 1888, futur admin), Ngaoundéré, Garoua, Dschang. Symboles : $\square$ Poste admin, $\triangle$ Poste militaire.
        \item \textbf{Item 2 : Infrastructures / Mise en valeur.} Traits : \textbf{Chemins de fer} (Nord 1911, Centre 1911, Sud 1914 - couleurs distinctes ou traits pointillés/pleins). \textbf{Routes principales} (traits fins). \textbf{Ports} (ancre).
        \item \textbf{Item 3 : Économie de plantation / Ressources.} Zones colorées (hachures) : \textbf{Cacao} (Sud-Ouest, Sanaga), \textbf{Caoutchouc} (Sud, Est), \textbf{Palmiers/Huile} (Côte), \textbf{Coton/Bétail} (Nord/Adamaoua). \textbf{Mines} (symboles $\times$) : Étain (Bertoua), Or (Bétaré-Oya).
        \item \textbf{Item 4 : Résistances / Limites.} Flèches rouges : Directions expéditions allemandes. Étoiles rouges : Foyers résistance majeurs (Bafut, Bamoun, Douala 1914, Ngaoundéré, Mora). Trait discontinu épais : \textbf{Frontière 1884} vs \textbf{Frontière 1911} (Neukamerun - hachuré).
    \end{itemize}
    \item \textbf{Encrer et colorier :} Feutre fin noir (noms, traits). Crayons de couleur (zones). Pas de surcharge. Titre centré en haut : \textbf{« Le Kamerun allemand : organisation, économie et résistances (1884-1914) »}.
    \item \textbf{Vérifier :} Légende complète ? Hiérarchie visuelle respectée ? Noms orthographes exactes ? Nord/Échelle présents ?
\end{enumerate}
\end{methode}

\begin{exoresolu}[Réalisation commentée : Croquis du Kamerun allemand]
\textbf{Énoncé :} Réalisez un croquis légendé et organisé du Kamerun sous administration allemande (1884-1914) mettant en évidence : l'organisation administrative, les axes de communication, les zones de production et les principaux foyers de résistance.
\tcblower
\textbf{Solution détaillée (Description du croquis attendu - Guide de correction) :}
\textbf{1. Titre :} \emph{Le Kamerun allemand : administration, économie, résistances (1884-1914)}.
\textbf{2. Base cartographique (TikZ simplifié pour l'examen) :}
\begin{itemize}
    \item Contour : Triangle côtier (Douala-Victoria-Kribi) + prolongement Nord (Bénoué) + « jambe » Extrême-Nord (lac Tchad) + excroissance Est (Neukamerun 1911 : Sangha, Lobaye).
    \item Repères : Fleuves (Wouri, Sanaga, Nyong, Bénoué, Logone, Chari). Lacs (Tchad, Nyos, Monoun).
\end{itemize}
\textbf{3. Légende structurée (4 items) :}
\begin{center}
\begin{tabular}{|l|l|l|}
\hline
\textbf{Item} & \textbf{Thème} & \textbf{Sémiologie (Symboles/Couleurs)} \\
\hline
1 & \textbf{Administration \& Villes} & \textcolor{popblue}{$\blacksquare$} Douala (Capitale éco, 1er port) ; \textcolor{popblue}{$\blacktriangle$} Buea (Capitale adm 1901) ; \textcolor{popblue}{$\bullet$} Yaoundé, Ngaoundéré, Garoua, Dschang, Bertoua (Chefs-lieux de Bezirk/Residenz). Noms en \textbf{noir}. \\
\hline
2 & \textbf{Infrastructures \& Transport} & \textcolor{popdark}{\textbf{---}} Chemin de fer Nord (Douala-Ngaoundéré, ach. 1911) ; \textcolor{popdark}{\textbf{-- --}} Chemin de fer Centre (Douala-Yaoundé, ach. 1911) ; \textcolor{popdark}{\textbf{$\cdots$}} Chemin de fer Sud (Douala-Kribi/Lolodorf, en cours 1914). \textcolor{popgray}{--} Routes principales. \text{Port} Ports (Victoria, Douala, Kribi). \\
\hline
3 & \textbf{Mise en valeur économique} & \textcolor{popgreen!50!popdark}{\textbf{Hachures $\diagup$}} Zone Cacao (Sud-Ouest, Mungo, Sanaga). \textcolor{poporange!70!popdark}{\textbf{Hachures $\diagdown$}} Zone Caoutchouc sauvage (Sud, Est, Nyong). \textcolor{popgold!80!popdark}{\textbf{Points}} Palmiers à huile (Côte, Wouri). \textcolor{poppurple}{\textbf{Losanges}} Coton / Bétail (Adamaoua, Nord). $\textcolor{popred}{\times}$ Mines (Étain Bertoua, Or Bétaré-Oya). \\
\hline
4 & \textbf{Résistances \& Frontières} & \textcolor{popred}{$\star$} Foyers résistance : Bafut, Bamoun (Njoya ambivalent), Douala (Manga Bell 1914), Bassa/Bakoko (1890s), Ngaoundéré (Mora), Adamaoua (Lamibé). \textcolor{popred}{\textbf{Flèches}} Directions colonnes allemandes (pacification). \textbf{Trait discontinu épais \textcolor{popblue}{- - -}} Frontière 1884. \textbf{Trait continu épais \textcolor{popdark}{—}} Frontière 1911 (Neukamerun hachuré \textcolor{popgray}{\%}). \\
\hline
\end{tabular}
\end{center}
\textbf{4. Qualité graphique (Barème Bac) :}
\begin{itemize}
    \item Propreté / Soin (règle, équerre, crayons affûtés) : 2 pts.
    \item Respect sémiologie (ponctuel/linéaire/surfacique) : 3 pts.
    \item Légende hiérarchisée, complète, exploitable : 3 pts.
    \item Exactitude géographique (position relative villes/fleuves/frontières) : 2 pts.
    \item Titre, Nord, Échelle, Source : 2 pts. \textbf{Total : 12 pts.}
\end{itemize}
\textbf{Conclusion :} Ce croquis synthétise la \textbf{logique extractive} (chemins de fer vers zones de production, ports d'exportation), la \textbf{pénétration administrative} (quadrillage poste militaire/admin) et la \textbf{réaction africaine} (foyers résistance corrélés aux zones de forte contrainte ou chefferies fortes).
\end{exoresolu}

\begin{methode}[Rédiger un paragraphe argumenté type Bac : Bilan de la période allemande]
\textbf{Objectif :} Répondre à une question de synthèse (ex: « Le bilan de la période allemande est-il positif ? ») avec une structure thèse/antithèse/synthèse.
\begin{enumerate}
    \item \textbf{Analyser le sujet :} Mots-clés : « Bilan », « Période allemande (1884-1916) », « Positif/Négatif » $\rightarrow$ Nuancer : \textbf{Contraste} entre modernisation infrastructurelle et violence structurelle.
    \item \textbf{Construire le plan dialectique :}
    \begin{itemize}
        \item \textbf{Thèse (Apports / Modernisation) :} État centralisé, frontières fixées, infrastructures (rails, routes, ports), sanitaire (hôpitaux, lutte sommeil/variole), éducation (écoles missions subventionnées, alphabetisation), écriture (Njoya), introduction cultures rentes (cacao, café, palmier).
        \item \textbf{Antithèse (Coûts / Violences) :} Perte souveraineté, spoliation foncière, travail forcé (mortalité), impôt injuste, répression sanglante résistances, justice inéquitable (Code pénal indigène), ségrégation raciale (décret 1913 Douala), démographie stagnante/déclin local.
        \item \textbf{Synthèse (Héritage contradictoire) :} « Modernisation sans développement humain ». Héritage institutionnel (État, frontières, administration, justice, éducation) récupéré par Français/Anglais. Traumatisme mémoriel (Manga Bell, Bafut). Base du nationalisme (revendication droits, unité).
    \end{itemize}
    \item \textbf{Rédiger :} Introduction (accroche, définition, problématique, plan). Développement (2 paragraphes équilibrés, arguments \textbf{chiffrés/datés} : km rails, tonnes cacao, nombre exécutés). Conclusion (réponse nuancée, ouverture sur mandat français/britannique).
    \item \textbf{Relire :} Vérifier : Dates exactes ? Noms propres ? Vocabulaire précis (protectorat, mandat, chefferie, Zwangsarbeit) ? Orthographe ?
\end{enumerate}
\end{methode}

\begin{exoresolu}[Bilan contradictoire de la période allemande (1884-1916)]
\textbf{Énoncé :} « La colonisation allemande a modernisé le Cameroun mais au prix d'une violence inouïe. » Discutez cette affirmation en structurant votre réponse en thèse, antithèse et synthèse.
\tcblower
\textbf{Solution détaillée (Plan de dissertation rédigé) :}
\textbf{Introduction}
\begin{itemize}
    \item \textbf{Accroche :} Le 12 juillet 1884, le traité Douala-Allemagne scelle le sort du « Kamerun » pour 32 ans.
    \item \textbf{Définition :} Période allemande = Protectorat (1884) $\rightarrow$ Colonie de peuplement/exploitation (1884-1916). Fin : conquête franco-britannique (1916), mandat SDN (1919/1922).
    \item \textbf{Problématique :} Comment évaluer un héritage colonial qui a posé les bases de l'État moderne camerounais par la contrainte, la spoliation et la répression ?
    \item \textbf{Plan :} I. Une modernisation infrastructurelle et institutionnelle indéniable (Thèse). II. Une entreprise de domination violente et prédatrice (Antithèse). III. Un héritage ambivalent : les fondations d'un État sans nation (Synthèse).
\end{itemize}
\textbf{Développement}
\textbf{I. Une modernisation infrastructurelle et institutionnelle indéniable (Thèse)}
\begin{itemize}
    \item \textbf{État et Frontières :} Création de l'entité « Kamerun » (475 000 km² en 1911). Frontières internationales fixées (traité 1885, 1894, 1911). Administration centralisée (Gouverneur, Bezirke), justice codifiée, monnaie (Mark), poste/télégraphe.
    \item \textbf{Infrastructures de mise en valeur :} \textbf{~1 000 km de chemins de fer} (Nord 1911, Centre 1911, Sud 1914) reliant l'intérieur au port de Douala (port en eau profonde, quai, grues). \textbf{Routes} (3 000 km pistes carrossables). Télégraphe (liaison Europe).
    \item \textbf{Santé \& Éducation :} Lutte contre maladies (variole, trypanosomiase - camps de concentration sanitaires controversés mais efficaces). Hôpitaux (Douala, Buea, Yaoundé). Écoles missions (Baptistes, Basel, Pallotins, Catholiques) subventionnées : \textbf{~200 000 élèves} en 1914. Alphabétisation en langues locales (Duala, Bassa, Bamoun, Fulfulde). \textbf{Invention écriture Shümom (Njoya, 1896)}.
    \item \textbf{Économie de rente :} Introduction/intensification : \textbf{Cacao} (premier exportateur mondial 1913 ~ 12 000 t), \textbf{Café}, \textbf{Palmier à huile}, \textbf{Caoutchouc}. Plantations scientifiques (Victoria, Soppo, Edéa). Monétarisation économie.
\end{itemize}
\textbf{II. Une entreprise de domination violente et prédatrice (Antithèse)}
\begin{itemize}
    \item \textbf{Violence de la conquête :} « Pacification » = guerre totale. \textbf{~ 30 expéditions militaires majeures} (1889-1914). Villages rasés (Bafut, Bassa, Maka, Kirdi). Exécutions sommaires, prises d'otages, déportations (Manga Bell, Ngoso Din, Abumbi I). \textbf{Démographie :} Stagnation, voire déclin local (fuite, mortalité travail forcé, maladies).
    \item \textbf{Exploitation juridique et économique :} \textbf{Travail forcé (Zwangsarbeit)} : Code pénal 1896, décret 1900. Chemin de fer Nord : \textbf{~ 10 000 à 20 000 morts} (estimations historiens) sur 10 ans (maladies, accidents, épuisement). Portage à dos d'homme (avant rail).
    \item \textbf{Spoliation foncière :} Décrets 1896, 1899, 1901, 1913 (Douala). « Terres vacantes » = 90\% territoire concédé aux compagnies (Woermann, Jantzen, Plantations Kamerun). Déplacement populations (Bassa, Bakoko, Duala vers « Neue Duala » / New Bell).
    \item \textbf{Ségrégation raciale :} Code de l'indigénat (1896, durci 1904/1911). Justice à deux vitesses. Décret 1913 Douala : zones distinctes Européens/Indigènes (apartheid naissant).
\end{itemize}
\textbf{III. Un héritage ambivalent : les fondations d'un État sans nation (Synthèse)}
\begin{itemize}
    \item \textbf{Récupération par les Mandats :} Français (1919) et Britanniques (1922) reprennent \textbf{l'appareil d'État} (frontières, administration, justice, chemins de fer, cadastre, personnel africain formé). Continuité administrative forte.
    \item \textbf{Naissance du nationalisme :} La violence commune (travail forcé, impôt, répression) et l'école coloniale créent une \textbf{conscience nationale partagée} (élites lettrées : Manga Bell, Martin Paul Samba, Louis-Marie Pouka, plus tard Um Nyobé). Revendication : \textbf{respect du traité 1884, égalité droits, fin travail forcé}.
    \item \textbf{Paradoxe :} L'Allemagne a « inventé » le Cameroun territorial et étatique, mais a détruit les bases d'une légitimité partagée. L'héritage est \textbf{structurel (État) mais pas consensuel (Nation)}.
\end{itemize}
\textbf{Conclusion}
\begin{itemize}
    \item \textbf{Bilan :} Modernisation \textbf{par le haut, pour l'extraction}, payée par les corps et les terres des Camerounais. « Développement sans développement humain » (Amin).
    \item \textbf{Ouverture :} Ce legs contradictoire explique les tensions post-indépendance : État fort (héritage allemand/français) vs demande de démocratie/décentralisation (héritage résistances/chefferies). La réunification 1961 tente de résoudre le démantèlement du Kamerun allemand (partage 1916).
\end{itemize}
\end{exoresolu}

\begin{methode}[Analyser une source iconographique (Photo coloniale / Carte de propagande)]
\textbf{Objectif :} Déconstruire le regard colonial pour en extraire l'information historique.
\begin{enumerate}
    \item \textbf{Description objective (Ce que je vois) :} Type (photo, gravure, carte, timbre), date, auteur (photographe colonial, missionnaire, administration), support. Éléments visuels : personnages (tenues, attitudes, outils), décor (architecture, végétation, infrastructures), composition (cadrage, mise en scène).
    \item \textbf{Analyse de l'intention / Propagande :} Quel message l'auteur veut-il faire passer ?
    \begin{itemize}
        \item \textbf{Missionnaire :} « Civilisation » (école, hôpital, tenue européenne, conversion).
        \item \textbf{Administratif :} « Ordre / Progrès » (gare, route, chef en uniforme prussien, troupeau, plantation ordonnée).
        \item \textbf{Scientifique / Ethnographique :} « Types raciaux » (mesures anthropométriques, nudité imposée, classification).
    \end{itemize}
    \item \textbf{Critique interne (Ce que l'image cache / montre malgré elle) :}
    \begin{itemize}
        \item Regard des sujets (fierté, peur, soumission, jeu).
        \item Détails anachroniques ou contradictoires (chaînes aux pieds cachées, fusils tenus par askaris, cultures vivrières remplacées par rentes).
        \item Absents : Femmes (souvent invisibles), résistants, corps, funérailles.
    \end{itemize}
    \item \textbf{Mise en perspective historique :} Confronter l'image aux sources écrites (rapports admin, lettres chefs, statistiques). L'image \emph{illustre} mais ne \emph{prouve} pas seule.
    \item \textbf{Conclusion critique :} L'image est un \textbf{témoignage sur le colonisateur} (son regard, ses peurs, ses fantasmes) autant que sur le colonisé.
\end{enumerate}
\end{methode}

\begin{exoresolu}[Analyse d'une photographie : « Chef Bafut et administrateur allemand » (vers 1908)]
\textbf{Énoncé :} Vous observez une photographie en noir et blanc datée d'environ 1908. Elle montre un officier allemand en uniforme (casque colonial, épée), assis sur une chaise de camp, face à un chef africain vêtu d'un pagne prestigieux, coiffé d'un chapeau à plumes, entouré de notables. Au second plan, des askaris (soldats africains en uniforme kaki) et des porteurs. Analysez ce document comme source historique.
\tcblower
\textbf{Solution détaillée :}
\begin{enumerate}
    \item \textbf{Identification et description objective :}
    \begin{itemize}
        \item \textbf{Nature :} Photographie argentique (gelatinobromure), format plaque de verre probable. \textbf{Date :} \emph{c. 1908} (période « pacification » Grassfields, gouverneur von Puttkamer/Seitz). \textbf{Auteur probable :} Photographe de la \textbf{Schutztruppe} ou missionnaire (Basel Mission active à Bafut) ou administrateur amateur.
        \item \textbf{Scène :} Rencontre protocolaire (« Palaver ») devant une case à toit conique (architecture Grassfields) ou au poste militaire.
        \item \textbf{Personnages :}
            \begin{itemize}
                \item \textbf{Allemand :} Assis, position dominante (hauteur, chaise), uniforme kaki/casque colonial (symbole autorité impériale), épée au côté (pouvoir coercitif), carnet/notes (bureaucratie).
                \item \textbf{Chef africain (Fon de Bafut ?) :} Debout ou assis sur tabouret bas, pagne tissé (ndop ?), chapeau à plumes (insigne royauté), colliers de perles/ivoire (richesse, statut). Attitude : droite, regard direct ou baissé selon protocole.
                \item \textbf{Entourage :} Notables (respect), Askaris (fusils Mauser, baïonnettes = force), Porteurs (charges, sueur = travail forcé).
            \end{itemize}
    \end{itemize}
    \item \textbf{Analyse de la mise en scène (Intention propagandiste) :}
    \begin{itemize}
        \item \textbf{Composition :} Cadrage centré sur l'échange. Symétrie trompeuse : un siège (pouvoir) vs tabouret/pagne (tradition).
        \item \textbf{Message visé :} \textbf{« Règle indirecte harmonieuse ».} L'Allemagne respecte les « chefs naturels » (pagne, plumes). Collaboration élites locales. Ordre (askaris alignés). Modernité (chaise, carnet, uniforme) encadrant Tradition.
        \item \textbf{Cible :} Opinion publique allemande (cartes postales, rapports coloniaux), chefs rivaux (montrer soumission), administration centrale (Berlin) : « Pacification réussie ».
    \end{itemize}
    \item \textbf{Critique interne / Lecture « à contre-courant » (Ce que l'image ne dit pas) :}
    \begin{itemize}
        \item \textbf{Contexte réel :} 1908 = fin guerre Bafut (1901-1910). Le Fon Abumbi I est en résistance ou vient de se rendre sous contrainte. La « palaver » est une \textbf{reddition négociée} ou une soumission forcée.
        \item \textbf{Indices de contrainte :} Présence massive \textbf{askaris armés} (ratio 1 pour 10 civils ?). Porteurs chargés (corvée). Absence de femmes/enfants (guerre, peur). Regard du Fon : défi contenu ou résignation.
        \item \textbf{Invisibles :} Villages brûlés, greniers détruits, otages, morts de la colonne allemande.
    \end{itemize}
    \item \textbf{Mise en perspective historique :}
    \begin{itemize}
        \item Confronter au \textbf{rapport du Bezirksamtmann} (ex: von Pavel) : « Le chef soumis, peuple pacifié ».
        \item Confronter à la \textbf{tradition orale Bafut} : « Le Blanc a volé la parole, nous a trompés par le fusil ».
        \item Confronter aux \textbf{statistiques} : Chute production vivrière, fuite population vers Nigeria britannique (1908-1910).
    \end{itemize}
\end{enumerate}
\textbf{Conclusion :} Cette photographie est une \textbf{mise en scène de la domination}. Elle documente la \textbf{ritualisation de la soumission} (cadeaux, honneurs au chef) qui masque la \textbf{violence structurelle} (armée, travail forcé, spoliation). Elle prouve l'existence de la \textbf{règle indirecte} (chef maintenu comme intermédiaire) mais révèle par ses détails (askaris, porteurs, absence femmes) la réalité coercitive du « Schutzgebiet ». Source précieuse \emph{si} croisée avec archives écrites et mémoire orale.
\end{exoresolu}

\begin{attention}[Les 5 erreurs qui coûtent des points au Bac (Histoire - Période Allemande)]
\begin{enumerate}
    \item \textbf{Confondre « Traité de protectorat » et « Traité de cession / Annexion ».} $\rightarrow$ \textbf{Piège :} Dire « les chefs ont vendu leur pays ». \textbf{Juste :} « Les chefs ont signé un traité de \emph{protection} (art. 2 : droits maintenus), que l'Allemagne a unilatéralement transformé en annexion administrative. » Nuance juridique cruciale (Art. 2 vs pratique).
    \item \textbf{Oublier le \cle{Neukamerun} (1911) dans les frontières.} $\rightarrow$ \textbf{Piège :} Dessiner la carte du Cameroun actuel pour 1914. \textbf{Juste :} Le Kamerun 1914 est \textbf{plus grand} (Neukamerun : parties RCA, Congo, Tchad actuels). Citer le traité franco-allemand de 1911 (échange Maroc/Congo). Perte de 2 pts carte + 2 pts dissertation.
    \item \textbf{Réduire la résistance à la seule lutte armée (Bafut/Bamoun) et ignorer la résistance juridique (Manga Bell) et passive.} $\rightarrow$ \textbf{Piège :} « Les Camerounais se sont battus avec des arcs ». \textbf{Juste :} Typologie : Armée (Bafut, Kirdi), Juridique/Politique (Manga Bell, pétitions Reichstag), Passive (fuite, refus impôt, sabotage), Culturelle (écriture Shümom, églises indépendantes). Montrer la \textbf{modernité} de certaines résistances.
    \item \textbf{Attribuer aux Allemands des réalisations françaises (ou inversement) :} $\rightarrow$ \textbf{Pièges classiques :} « Les Allemands ont construit le barrage d'Edéa » (FAUX : études allemandes, construction \textbf{Français 1940s}). « Les Allemands ont fait le chemin de fer Transcamerounais » (FAUX : Nord + Centre seulement, Transcamerounais = Français 1960s-70s). « Douala capitale allemande » (FAUX : Buea 1901-1919). \textbf{Précision :} Chemins de fer : Nord (Douala-Ngaoundéré), Centre (Douala-Yaoundé), Sud (Douala-Kribi inachevé).
    \item \textbf{Rédiger une conclusion « moralisatrice » sans analyse historique.} $\rightarrow$ \textbf{Piège :} « C'était mal / bien ». \textbf{Juste :} Conclusion \textbf{nuancée et structurée} : « Héritage contradictoire : fondation de l'État territorial et administratif (frontières, rails, administration, éducation) par une violence extractive (travail forcé, spoliation, répression) qui a engendré le nationalisme camerounais. » Utiliser le vocabulaire du cours : \cle{Zwangsarbeit}, \cle{Schutztruppe}, \cle{Indirekte Herrschaft}, \cle{Neukamerun}, \cle{Traité 1884}.
\end{enumerate}
\end{attention}

\newpage

\coursec{Exercices}

\courssub{Je vérifie mes connaissances}

\begin{exercice}[QCM : Repères chronologiques et acteurs]
Parmi les propositions suivantes, une seule est exacte. Recopie la bonne réponse sur ta copie.
\begin{enumerate}
    \item La capitale du Kamerun a été transférée de Douala à Buea en :
    \begin{itemize}
        \item[A.] 1885
        \item[B.] 1901
        \item[C.] 1911
        \item[D.] 1914
    \end{itemize}
    \item Le traité de protection germano-douala du 12 juillet 1884 a été signé côté allemand par :
    \begin{itemize}
        \item[A.] Gustav Nachtigal
        \item[B.] Julius von Soden
        \item[C.] Jesko von Puttkamer
        \item[D.] Theodor Seitz
    \end{itemize}
    \item La résistance de Rudolf Duala Manga Bell (1914) s'est principalement opposée à :
    \begin{itemize}
        \item[A.] L'impôt de capitation
        \item[B.] Le projet d'expropriation des terres riveraines à Douala (ordonnances de 1912-1913)
        \item[C.] Le travail forcé sur les plantations
        \item[D.] L'interdiction du commerce de l'huile de palme
    \end{itemize}
    \item L'extension territoriale du Kamerun vers l'Est et le Sud (le « Nouveau Cameroun » ou Neukamerun) résulte de l'accord de :
    \begin{itemize}
        \item[A.] Berlin (1885)
        \item[B.] Fès (1911)
        \item[C.] Versailles (1919)
        \item[D.] Londres (1919)
    \end{itemize}
\end{enumerate}
\end{exercice}

\begin{exercice}[Vrai / Faux justifié]
Indique si chaque affirmation est VRAIE ou FAUSSE, puis justifie brièvement ta réponse par un fait historique précis.
\begin{enumerate}
    \item Le Kamerun a eu trois capitales successives sous l'administration allemande : Douala, Buea, Yaoundé.
    \item Le gouverneur Jesko von Puttkamer a mis fin à la politique de « domination indirecte » pour instaurer une administration directe sur l'ensemble du territoire.
    \item La résistance du peuple Bassa a été dirigée par le chef Madola et a utilisé la guérilla dans la forêt.
    \item Les chemins de fer du Kamerun ont été construits uniquement pour relier les ports aux capitales administratives.
\end{enumerate}
\end{exercice}

\begin{exercice}[Définitions et identification]
Donne une définition précise (deux lignes maximum) ou identifie le personnage/terme suivant :
\begin{enumerate}
    \item \textbf{Traité de protection} (contexte de 1884).
    \item \textbf{Indigénat} (régime juridique).
    \item \textbf{Charles Atangana}.
    \item \textbf{Accord de Fès (1911)}.
\end{enumerate}
\end{exercice}

\begin{exercice}[Calculs et données chiffrées]
\begin{enumerate}
    \item Calcule la durée exacte (en années) de la période allemande au Kamerun (de la signature du traité de 1884 à la chute de Mora en février 1916).
    \item Le Kamerun initial (1884) faisait environ 495 000 km\textsuperscript{2}. Après 1911, sa superficie passe à environ 790 000 km\textsuperscript{2}. Calcule le pourcentage d'augmentation de la surface territoriale (arrondi à l'entier).
    \item En 1913, le budget du Kamerun est équilibré à 25 millions de marks. Les recettes douanières représentent 60\,\% de ce budget. Calcule le montant des recettes douanières.
\end{enumerate}
\end{exercice}

\courssub{Je m'entraîne}

\begin{exercice}[Analyse du traité germano-douala du 12 juillet 1884]
À partir du texte simplifié ci-dessous (extrait des clauses principales), réponds aux questions.
\begin{quote}
\textit{« Sa Majesté le Roi des Douala... cède à l'Empire allemand tous ses droits souverains sur le territoire... L'Empire allemand garantit aux rois et chefs leurs droits de propriété et leurs usages coutumiers... Les sujets allemands jouissent de la liberté du commerce... L'administration allemande ne percevra pas d'impôts sans l'accord des chefs. »}
\end{quote}
\begin{enumerate}
    \item Identifie les deux parties contractantes et leurs représentants principaux.
    \item Relève deux avantages accordés aux chefs douala par ce traité.
    \item Relève deux avantages accordés à l'Allemagne.
    \item Explique pourquoi ce traité est considéré comme l'acte de naissance juridique du « Kamerun ».
\end{enumerate}
\end{exercice}

\begin{exercice}[L'organisation administrative : du protectorat à la colonie]
Complète le tableau suivant qui résume l'évolution de l'administration allemande (1884-1914).
\begin{center}
\begin{tabularx}{\linewidth}{|l|X|X|X|}\hline
\rowcolor{popblueL}\textbf{Période} & \textbf{Gouverneur / Administrateur} & \textbf{Capitale} & \textbf{Principale caractéristique administrative} \\ \hline
1884-1885 & & Douala & Administration consulaire, appui des firmes commerciales \\ \hline
1885-1891 & Julius von Soden & & Mise en place de l'administration d'État, premières lois coloniales \\ \hline
1895-1907 & & Buea & \textbf{Domination indirecte} : gouvernement par les chefs traditionnels \\ \hline
1907-1910 & Theodor Seitz & & Début de l'administration directe, répression des résistances \\ \hline
1910-1914 & & \textbf{Yaoundé} & Administration directe affermie, grands travaux publics \\ \hline
\end{tabularx}
\end{center}
\end{exercice}

\begin{exercice}[Économie du Kamerun : productions et infrastructures]
\begin{enumerate}
    \item Cite les trois principaux produits d'exportation du Kamerun vers 1913.
    \item Nomme les deux lignes de chemin de fer construites par les Allemands et précise pour chacune ses extrémités.
    \item Explique en deux phrases le rôle du « travail forcé » (Zwangsarbeit) dans la réalisation de ces infrastructures.
\end{enumerate}
\end{exercice}

\begin{exercice}[Cartographie : Le Kamerun en 1914]
Sur un croquis simplifié du Kamerun (à reproduire schématiquement), place et nomme :
\begin{enumerate}
    \item Les trois capitales successives : Douala, Buea, Yaoundé.
    \item Les deux lignes de chemin de fer (Nordbahn et Mittellandbahn).
    \item Deux foyers de résistance majeure : un dans la zone côtière et forestière (ex. Douala, Bassa), un dans la zone septentrionale (ex. Adamawa, pays Bafia).
    \item Un fleuve servant d'axe de pénétration (Wouri, Sanaga, Nyong, Bénoué ou Logone).
\end{enumerate}
\textit{Trace une légende organisée (signes ponctuels, linéaires, surfaciques) et donne un titre au croquis.}
\end{exercice}

\begin{popfigure}
\begin{center}
\begin{tikzpicture}[scale=0.9, every node/.style={font=\small}]
% Contour très simplifié du Kamerun
\draw[thick, popdark, fill=popcream] (0,0) -- (0.5,2.2) -- (1.2,4.5) -- (2.2,6.5) -- (3.2,7.2) -- (4.5,6.8) -- (5.2,5.5) -- (5.8,4.2) -- (5.2,2.8) -- (4.2,1.8) -- (3.0,1.2) -- (1.5,0.4) -- cycle;
% Neukamerun (sud-est)
\draw[thick, poporange, fill=poporangeL, opacity=0.6] (4.2,1.8) -- (5.2,2.8) -- (6.6,2.2) -- (6.8,0.6) -- (5.2,0.2) -- (3.0,1.2) -- cycle;
\node[poporange] at (5.6,1.3) {\textbf{Neukamerun}};
\node[poporange] at (5.6,0.9) {(1911)};
% Villes
\fill[popink] (0.7,2.4) circle (2pt) node[right=2pt, popink]{\textbf{Douala}};
\fill[popink] (0.5,3.0) circle (2pt) node[left=2pt, popink]{\textbf{Buea}};
\fill[popink] (2.6,2.6) circle (2pt) node[right=2pt, popink]{\textbf{Yaoundé}};
% Chemins de fer
\draw[very thick, poppurple] (0.7,2.4) -- (1.6,3.4) node[right=1pt, poppurple]{Nkongsamba};
\draw[very thick, popgreen] (0.7,2.4) -- (2.0,2.2) node[below=1pt, popgreen]{Éséka};
% Fleuves
\draw[thick, popblue] (0.4,2.0) .. controls (1.2,1.6) and (1.8,1.4) .. (2.6,1.5) node[right, popblue]{Sanaga};
\draw[thick, popblue] (3.4,6.6) .. controls (3.8,5.8) and (4.2,5.2) .. (4.6,4.6) node[right, popblue]{Bénoué};
% Résistances
\node[popdark] at (1.3,1.7) {$\times$};
\node[popdark, right=1pt] at (1.35,1.7) {Pays Bassa};
\node[popdark] at (3.6,5.6) {$\times$};
\node[popdark, right=1pt] at (3.65,5.6) {Adamawa};
% Nord
\draw[-{Stealth}, thick, popdark] (6.4,6.2) -- (6.4,7.2) node[above, popdark]{\textbf{N}};
\end{tikzpicture}
\end{center}
\legende{Croquis schématique du Kamerun allemand en 1914 : capitales (points), chemins de fer Nordbahn (violet) et Mittellandbahn (vert), foyers de résistance ($\times$), Neukamerun (zone orange). Contour approximatif, sans échelle.}
\end{popfigure}

\courssub{Je me prépare au Bac}

\begin{exercice}[Sujet type Probatoire : Dissertation guidée]
\textbf{Sujet :} « La période allemande (1884-1916) a-t-elle été bénéfique pour le Cameroun ? »
Rédige une dissertation structurée en respectant les étapes suivantes :
\begin{enumerate}
    \item \textbf{Introduction :} accroche (contexte de la colonisation européenne), définition des termes (bénéfique : pour qui ? quel bilan ?), problématique, annonce du plan.
    \item \textbf{Partie I : Les apports et les transformations structurelles (arguments « pour »).}
    Développe deux arguments majeurs (ex. unification territoriale, infrastructures, éducation, santé, introduction de cultures rentières) en t'appuyant sur des exemples précis (noms de gares, de plantes, d'écoles, de traités).
    \item \textbf{Partie II : Les coûts humains, l'exploitation et la violence coloniale (arguments « contre »).}
    Développe deux arguments majeurs (ex. travail forcé, expropriations, répression militaire, régime de l'indigénat, pertes démographiques) avec des exemples datés et localisés (résistances Bassa, Bafia, Douala, Adamawa).
    \item \textbf{Conclusion :} bilan nuancé (héritage ambivalent), ouverture sur la période du mandat (Société des Nations, partage franco-britannique de 1916-1919).
\end{enumerate}
\end{exercice}

\begin{methode}[Réussir la dissertation d'Histoire au Probatoire technique]
\begin{enumerate}
    \item \textbf{Analyser le sujet :} souligner les mots-clés (« bénéfique », « période allemande », « Cameroun ») et délimiter le cadre chronologique (1884-1916).
    \item \textbf{Poser une problématique :} transformer le sujet en question directrice (ex. « Quel est le bilan réel de la colonisation allemande pour les populations camerounaises ? »).
    \item \textbf{Mobiliser des exemples précis :} chaque argument doit être prouvé par un fait daté, localisé, avec un acteur nommé.
    \item \textbf{Nuancer :} éviter le jugement manichéen ; montrer que les « apports » servaient d'abord les intérêts coloniaux.
    \item \textbf{Soigner la langue et la structure :} introduction, deux parties équilibrées, conclusion avec ouverture.
\end{enumerate}
\end{methode}

\begin{exercice}[Sujet type Baccalauréat Technique : Étude de document + Question de synthèse]
\textbf{Document :} Extrait adapté d'un rapport administratif de l'époque du gouverneur Theodor Seitz (vers 1909) sur la « politique indigène ».
\begin{quote}
\textit{« L'expérience a montré que l'administration par l'intermédiaire des chefs traditionnels présente des dangers... Le chef abuse souvent de son pouvoir, perçoit des impôts pour son compte, vend la justice... Il faut substituer progressivement l'administration directe : fonctionnaires allemands résidents, tribunaux allemands, impôt perçu par l'État. Le chemin de fer et la plantation exigent une main-d'œuvre disciplinée et disponible, non soumise à l'arbitraire des chefs. »}
\end{quote}
\textbf{Questions :}
\begin{enumerate}
    \item \textbf{Présentation du document :} nature, auteur, date, contexte historique (2 lignes).
    \item \textbf{Analyse :}
    \begin{itemize}
        \item[a)] Quels sont les deux reproches principaux faits aux chefs traditionnels par l'administration ?
        \item[b)] Quelles sont les deux mesures concrètes proposées pour instaurer l'administration directe ?
        \item[c)] Quels intérêts économiques sous-tendent ce changement de politique (cite le texte) ?
    \end{itemize}
    \item \textbf{Mise en perspective (synthèse) :} en t'appuyant sur le document et tes connaissances, montre comment le passage de la « domination indirecte » à l'« administration directe » (1907-1914) a modifié les rapports entre l'État colonial et les sociétés camerounaises (résistances, collaborations, transformations sociales). Réponse attendue : une vingtaine de lignes structurée en deux paragraphes.
\end{enumerate}
\end{exercice}

\begin{exercice}[Sujet type Bac : Question de cours construite]
\textbf{Sujet :} « Les Camerounais face à l'administration allemande : entre résistances et collaborations. »
Rédige une réponse construite (environ 30 lignes) en t'appuyant sur \textbf{au moins quatre exemples précis} (noms de chefs, de peuples, de dates, de lieux) pour illustrer les deux attitudes suivantes :
\begin{enumerate}
    \item \textbf{Les résistances armées et politiques :} distingue les résistances primaires (à la conquête, 1884-1907) et les résistances secondaires (contre l'administration directe, 1907-1914). Cite au moins deux foyers distincts (ex. Bassa, Bafia, Douala, Adamawa).
    \item \textbf{Les collaborations et adaptations :} explique les motivations des chefs collaborateurs (ex. Charles Atangana, sultan Ibrahim Njoya, chefs douala signataires de 1884) et les formes de cette collaboration (recrutement de porteurs, auxiliaires, médiation foncière, conversion religieuse).
\end{enumerate}
Conclus sur la complexité des stratégies africaines face à la contrainte coloniale.
\end{exercice}



\newpage

\coursec{Corrigés des exercices}

\coursec{Corrigés des exercices — La période allemande (Première Technique)}

\begin{corrige}[Exercice 1 — QCM : Repères chronologiques et acteurs]
\begin{enumerate}
    \item \textbf{B. 1901.} Sous le gouverneur Jesko von Puttkamer, la capitale est transférée de Douala à Buea, au pied du mont Cameroun, pour des raisons climatiques et sanitaires (Douala est jugée malsaine). Les autres dates correspondent à d'autres événements : 1885 (début de l'administration Soden), 1911 (Neukamerun), 1914 (début de la guerre).
    \item \textbf{A. Gustav Nachtigal.} Consul général d'Allemagne pour l'Afrique occidentale, il signe le 12 juillet 1884 le traité avec les rois Bell et Akwa, puis hisse le drapeau allemand le 14 juillet 1884. Soden, Puttkamer et Seitz sont des gouverneurs postérieurs.
    \item \textbf{B.} Le roi Rudolf Duala Manga Bell proteste contre le projet d'expropriation et de déplacement des populations douala des bords du Wouri (Joss Plateau), décidé par les ordonnances de 1912-1913. Accusé de haute trahison, il est pendu le 8 août 1914 avec son secrétaire Adolf Ngoso Din.
    \item \textbf{B. Fès (1911).} Par l'accord Maroc-Congo du 4 novembre 1911, la France cède à l'Allemagne des territoires congolais (le Neukamerun) en échange de sa liberté d'action au Maroc. Berlin (1885) règle le partage de l'Afrique ; Versailles et Londres (1919) concernent le partage du Kamerun après la guerre.
\end{enumerate}
\textbf{Points de vigilance :} ne pas confondre la date du traité (12 juillet 1884) avec celle de la proclamation du protectorat (14 juillet 1884) ; bien distinguer les quatre gouverneurs cités et leur ordre chronologique.
\end{corrige}

\begin{corrige}[Exercice 2 — Vrai / Faux justifié]
\begin{enumerate}
    \item \textbf{VRAI.} Douala est la première capitale (1884-1901), puis Buea (1901-1909 environ), enfin Yaoundé à partir de 1909-1910, ville fondée comme poste scientifique et militaire en 1889.
    \item \textbf{FAUX.} C'est l'inverse : Puttkamer (1895-1907) est le partisan de la domination indirecte (gouvernement par les chefs traditionnels). C'est son successeur Theodor Seitz (1907-1910) qui engage le passage à l'administration directe.
    \item \textbf{VRAI.} Les Bassa, conduits par le chef Madola, mènent une guérilla forestière contre la pénétration allemande ; la résistance est matée militairement au début des années 1900.
    \item \textbf{FAUX.} Les lignes (Nordbahn : Bonabéri-Nkongsamba ; Mittellandbahn : Douala-Éséka) servent d'abord l'économie de plantation : évacuer le caoutchouc, le cacao et l'huile de palme vers le port de Douala, et pénétrer l'intérieur du pays.
\end{enumerate}
\textbf{Points de vigilance :} une justification est exigée pour chaque item, même quand l'affirmation est vraie ; pour les items faux, il faut corriger l'affirmation et pas seulement la nier.
\end{corrige}

\begin{corrige}[Exercice 3 — Définitions et identification]
\begin{enumerate}
    \item \textbf{Traité de protection :} accord signé entre une puissance européenne et des chefs africains par lequel ces derniers cèdent leurs droits de souveraineté (relations extérieures, commerce) en échange d'une « protection » militaire et du respect promis de leurs coutumes et propriétés.
    \item \textbf{Indigénat :} régime juridique d'exception réservé aux populations colonisées : elles sont des « sujets » et non des citoyens, passibles de sanctions administratives (corvées, amendes, châtiments) sans procès régulier.
    \item \textbf{Charles Atangana (vers 1880-1943) :} chef supérieur des Ewondo-Béné, interprète puis collaborateur de l'administration allemande ; nommé « chef suprême » des peuples du Centre, il illustre la stratégie d'adaptation et de collaboration.
    \item \textbf{Accord de Fès (4 novembre 1911) :} accord franco-allemand réglant la crise du Maroc ; la France cède à l'Allemagne des territoires du Congo rattachés au Kamerun (Neukamerun), portant la colonie à environ 790\,000 km\textsuperscript{2}.
\end{enumerate}
\textbf{Points de vigilance :} une définition doit comporter la nature du terme, son contexte et sa fonction ; pour un personnage, donner dates, statut et rôle historique.
\end{corrige}

\begin{corrige}[Exercice 4 — Calculs et données chiffrées]
\begin{enumerate}
    \item Durée : $1916 - 1884 = 32$ ans. La période allemande dure environ \textbf{32 ans} (juillet 1884 -- février 1916, chute de Mora).
    \item Augmentation absolue : $790 - 495 = 295$ milliers de km\textsuperscript{2}. Taux d'augmentation :
    \[ \frac{295}{495} \times 100 \approx 59,6\,\% \approx 60\,\%. \]
    La superficie du Kamerun augmente d'environ \textbf{60\,\%} grâce au Neukamerun.
    \item Recettes douanières : $25 \times \dfrac{60}{100} = 15$. Les recettes douanières s'élèvent à \textbf{15 millions de marks}.
\end{enumerate}
\textbf{Points de vigilance :} pour le taux d'augmentation, le diviseur est la valeur de départ (495) et non la valeur d'arrivée (790) ; toujours préciser l'unité du résultat (km\textsuperscript{2}, \%, millions de marks).
\end{corrige}

\begin{corrige}[Exercice 5 — Analyse du traité germano-douala du 12 juillet 1884]
\begin{enumerate}
    \item Les deux parties sont : d'une part les rois et chefs douala (roi Ndumb'a Lobe Bell, roi Akwa et les chefs signataires), d'autre part l'Empire allemand représenté par le consul Gustav Nachtigal (avec l'appui des firmes commerciales allemandes installées à Douala, comme la maison Woermann).
    \item Avantages pour les chefs douala : la garantie de leurs droits de propriété et de leurs usages coutumiers ; la promesse qu'aucun impôt ne sera perçu sans leur accord ; la « protection » allemande contre leurs rivaux.
    \item Avantages pour l'Allemagne : la cession des droits souverains sur le territoire ; la liberté du commerce pour les sujets allemands (accès direct au marché, suppression des monopoles des intermédiaires douala).
    \item Ce traité est l'acte de naissance juridique du Kamerun car il fonde, pour la première fois, une entité territoriale placée sous souveraineté allemande : le drapeau allemand est hissé le 14 juillet 1884 et la puissance coloniale étend ensuite ce protectorat de la côte vers l'intérieur, créant un espace politique nouveau appelé « Kamerun ».
\end{enumerate}
\textbf{Points de vigilance :} les réponses aux questions 2 et 3 doivent s'appuyer sur des citations du texte ; la question 4 exige une mise en perspective (du traité local à l'entité territoriale), pas une simple paraphrase.
\end{corrige}

\begin{corrige}[Exercice 6 — L'organisation administrative : du protectorat à la colonie]
\begin{center}
\begin{tabularx}{\linewidth}{|l|X|X|X|}\hline
\rowcolor{popblueL}\textbf{Période} & \textbf{Gouverneur / Administrateur} & \textbf{Capitale} & \textbf{Principale caractéristique administrative} \\ \hline
1884-1885 & Gustav Nachtigal (consul impérial) & Douala & Administration consulaire, appui des firmes commerciales \\ \hline
1885-1891 & Julius von Soden & Douala & Mise en place de l'administration d'État, premières lois coloniales \\ \hline
1895-1907 & Jesko von Puttkamer & Buea & \textbf{Domination indirecte} : gouvernement par les chefs traditionnels \\ \hline
1907-1910 & Theodor Seitz & Buea & Début de l'administration directe, répression des résistances \\ \hline
1910-1914 & Karl Ebermaier & \textbf{Yaoundé} & Administration directe affermie, grands travaux publics \\ \hline
\end{tabularx}
\end{center}
\textbf{Points de vigilance :} le transfert de capitale à Yaoundé accompagne le passage à l'administration directe ; ne pas attribuer la domination indirecte à Seitz, qui la combat au contraire. \emph{Note : la période 1891-1895 (gouverneurs Leist puis Zimmerer) marque la transition vers la « pacification » militaire.}
\end{corrige}

\begin{corrige}[Exercice 7 — Économie du Kamerun : productions et infrastructures]
\begin{enumerate}
    \item Vers 1913, les trois principaux produits d'exportation sont : le \textbf{cacao}, le \textbf{caoutchouc} et l'\textbf{huile de palme} (avec les palmistes), issus des grandes plantations de la zone côtière.
    \item Les deux lignes construites sont : la \textbf{Nordbahn} (ligne du Nord), de Bonabéri à Nkongsamba (environ 160 km, achevée en 1911), et la \textbf{Mittellandbahn} (ligne du Centre), de Douala à Éséka (environ 160 km, achevée en 1914), destinée à être prolongée vers Yaoundé.
    \item Le travail forcé (\emph{Zwangsarbeit}) fournit gratuitement ou à très bas prix la main-d'œuvre nécessaire au terrassement, au portage et à la pose des voies. Les villageois sont réquisitionnés par les chefs ou directement par l'administration ; les conditions de travail très dures (maladies, accidents, malnutrition) provoquent une forte mortalité parmi les travailleurs.
\end{enumerate}
\textbf{Points de vigilance :} bien orthographier les noms allemands des lignes et associer chaque ligne à ses extrémités exactes ; le rôle du travail forcé doit être expliqué (fonction économique + coût humain), pas seulement nommé.
\end{corrige}

\begin{corrige}[Exercice 8 — Cartographie : Le Kamerun en 1914]
Le croquis attendu doit comporter un titre (« Le Kamerun allemand en 1914 »), une flèche du nord et une légende organisée.

\begin{popfigure}
\begin{tikzpicture}[scale=0.75, every node/.style={font=\small}]
% Frontière simplifiée du Kamerun 1914 (avec Neukamerun)
\draw[thick, popdark] (0,0) coordinate (SW) -- (9,0) coordinate (SE) -- (10,6) coordinate (NE) -- (4,10) coordinate (NW) -- (0,7) -- cycle;
% Fleuve Wouri
\draw[popblue, thick] (1.5,0.5) .. controls (2,1.5) and (2.5,2) .. (2.5,2.5) node[above right, black] {Wouri};
% Fleuve Sanaga
\draw[popblue, thick] (3,1) .. controls (4,3) and (5,4) .. (5.5,5) node[above right, black] {Sanaga};
% Fleuve Bénoué (Nord)
\draw[popblue, thick] (7,2) .. controls (7.5,4) and (8,5) .. (8.5,6) node[above left, black] {Bénoué};
% Nordbahn (Bonabéri - Nkongsamba)
\draw[poporange, very thick, -{Stealth[length=3pt]}] (1.5,0.5) -- (3.5,2.5) node[above, midway, sloped, black, font=\tiny] {Nordbahn};
% Mittellandbahn (Douala - Eséka - Yaoundé)
\draw[poporange, very thick, -{Stealth[length=3pt]}] (1.5,0.5) -- (4.5,3.5) node[above, midway, sloped, black, font=\tiny] {Mittellandbahn};
% Villes
\node[circle, fill=popblue, inner sep=2pt, label=below:{Douala}] at (1.5,0.5) {};
\node[circle, fill=popblue, inner sep=2pt, label=left:{Buéa}] at (0.8,1.8) {};
\node[circle, fill=popblue, inner sep=2pt, label=right:{Yaoundé}] at (4.5,3.5) {};
\node[circle, fill=popblue, inner sep=2pt, label=above:{Ngaoundéré}] at (7,5) {};
\node[circle, fill=popblue, inner sep=2pt, label=above:{Garoua}] at (8,5.5) {};
% Zones de plantations (côte)
\fill[popgreen!20] (0.5,0.2) rectangle (3.5,2);
\node[popgreen!70!black, font=\tiny] at (2,1) {Plantations};
% Neukamerun (hachures)
\draw[pattern=north east lines, pattern color=popgreen!50!black] (4,10) -- (10,6) -- (10,10) -- (4,10) -- cycle;
\node[popgreen!50!black, font=\small, rotate=-30] at (7,8) {Neukamerun (1911)};
% Foyers de résistance
\node[draw=poppink, fill=poppink!10, rounded corners=2pt, inner sep=2pt, font=\tiny] at (2.5,1.2) {Bassa (Madola)};
\node[draw=poppink, fill=poppink!10, rounded corners=2pt, inner sep=2pt, font=\tiny] at (7.5,6.5) {Adamawa};
% Flèche Nord
\node at (9.5,9.5) {\Huge \textbf{N} $\uparrow$};
% Titre
\node[align=center, font=\bfseries\small] at (5,10.5) {Le Kamerun allemand en 1914};
\end{tikzpicture}
\legende{Croquis corrigé : capitales (points bleus), chemins de fer (traits orange), fleuves (bleu), Neukamerun (hachures vertes), foyers de résistance (rose).}
\end{popfigure}

\textbf{Barème indicatif (sur 10) :} titre et orientation 1 pt ; légende organisée en trois catégories de signes 2 pts ; capitales correctement placées 2 pts ; chemins de fer corrects 2 pts ; foyers de résistance localisés 2 pts ; fleuves tracés et nommés 1 pt.
\textbf{Points de vigilance :} un croquis sans légende ni titre est sanctionné ; les villes doivent être placées dans le bon ordre côte-intérieur (Douala et Buea à l'ouest, Yaoundé au centre).
\end{corrige}

\begin{corrige}[Exercice 9 — Sujet type Probatoire : Dissertation guidée]
\textbf{Corrigé structuré (plan rédigé).}

\emph{Introduction.} Accroche : dans le cadre du partage de l'Afrique acté à Berlin (1884-1885), l'Allemagne impose son protectorat au Kamerun le 12 juillet 1884. Définition : « bénéfique » suppose un bilan positif, mais pour qui : la métropole ou les populations camerounaises ? Problématique : quel est le bilan réel de la colonisation allemande pour le Cameroun ? Annonce du plan : après les transformations structurelles indéniables (I), il faut mesurer les coûts humains et la violence coloniale (II).

\emph{Partie I — Les apports et transformations.} Argument 1 : la constitution d'une entité territoriale unifiée, le « Kamerun », doté de frontières, d'une capitale (Yaoundé) et d'une administration d'État. Argument 2 : les infrastructures et l'économie d'exportation : chemins de fer Nordbahn (Bonabéri-Nkongsamba, 1911) et Mittellandbahn (Douala-Éséka, 1914), port de Douala, plantations de cacao et de caoutchouc, premières écoles et postes sanitaires.

\emph{Partie II — Les coûts humains et la violence.} Argument 1 : l'exploitation : travail forcé sur les chantiers ferroviaires, impôt de capitation, expropriations (ordonnances de 1912-1913 à Douala), régime de l'indigénat. Argument 2 : la répression : guérilla Bassa matée, soulèvements Bafia et de l'Adamawa écrasés, pendaison du roi Rudolf Duala Manga Bell et de Ngoso Din le 8 août 1914, fortes pertes démographiques.

\emph{Conclusion.} Le bilan est ambivalent : les « apports » servaient d'abord les intérêts coloniaux et furent payés par les populations. Ouverture : la défaite allemande (Mora, février 1916) ouvre la période du partage franco-britannique et des mandats de la Société des Nations.

\textbf{Barème indicatif (sur 20) :} introduction (accroche, problématique, plan) 4 pts ; partie I (2 arguments prouvés) 6 pts ; partie II (2 arguments prouvés) 6 pts ; conclusion avec ouverture 2 pts ; langue et organisation 2 pts.
\textbf{Points de vigilance :} chaque argument doit être étayé par un fait daté et localisé ; éviter le plan « catalogue » sans problématique ; la nuance (bénéfique pour qui ?) est attendue dans la conclusion.
\end{corrige}

\begin{corrige}[Exercice 10 — Sujet type Baccalauréat Technique : Étude de document + synthèse]
\begin{enumerate}
    \item \textbf{Présentation :} il s'agit d'un rapport administratif colonial, rédigé dans l'entourage du gouverneur Theodor Seitz vers 1909, au moment où l'Allemagne abandonne la domination indirecte de l'ère Puttkamer pour imposer l'administration directe.
    \item \textbf{Analyse :}
    \begin{itemize}
        \item[a)] Les deux reproches : les chefs \textbf{abusent de leur pouvoir} (ils « perçoivent des impôts pour leur compte ») et ils \textbf{corrompent la justice} (ils « vendent la justice »).
        \item[b)] Les deux mesures : l'envoi de \textbf{fonctionnaires allemands résidents} et de \textbf{tribunaux allemands} ; la perception de l'\textbf{impôt par l'État} colonial et non plus par les chefs.
        \item[c)] Le texte déclare : « Le chemin de fer et la plantation exigent une main-d'œuvre disciplinée et disponible » : il s'agit de soustraire la main-d'œuvre à l'autorité des chefs pour la mettre au service des chantiers ferroviaires et des plantations.
    \end{itemize}
    \item \textbf{Synthèse (corrigé rédigé) :} Sous Puttkamer (1895-1907), l'Allemagne gouverne le Kamerun par l'intermédiaire des chefs traditionnels, reconnus et rémunérés : c'est la domination indirecte, peu coûteuse mais qui laisse aux chefs une large autonomie. À partir de 1907, Seitz puis Ebermaier imposent l'administration directe : fonctionnaires allemands dans les circonscriptions, tribunaux coloniaux, impôt perçu par l'État, recrutement direct de la main-d'œuvre. Ce tournant dépossède les chefs de leurs attributions et soumet les populations à une présence coloniale quotidienne et coercitive.

    Les conséquences sont profondes. D'un côté, la pression fiscale, les corvées et les expropriations (comme à Douala en 1912-1913) provoquent de nouvelles résistances : révoltes des Bafia, des peuples de l'Adamawa, et surtout l'opposition politique du roi Rudolf Duala Manga Bell, pendu en août 1914. De l'autre, certaines élites choisissent l'adaptation : Charles Atangana devient le relais privilégié de l'administration au Centre, tandis que le sultan Ibrahim Njoya à Foumban négocie la préservation de son royaume. L'administration directe transforme donc durablement les sociétés camerounaises : affaiblissement des pouvoirs traditionnels, émergence d'élites salariées et scolarisées, mais aussi enracinement d'un État autoritaire dont hériteront les mandats français et britannique.
\end{enumerate}
\textbf{Barème indicatif (sur 20) :} présentation du document 3 pts ; analyse a) 3 pts, b) 3 pts, c) 3 pts ; synthèse structurée en deux paragraphes avec exemples 6 pts ; langue 2 pts.
\textbf{Points de vigilance :} en analyse de texte, citer le document entre guillemets pour prouver chaque réponse ; la synthèse doit croiser le document ET les connaissances de cours, et respecter la consigne des deux paragraphes.
\end{corrige}

\begin{corrige}[Exercice 11 — Sujet type Bac : Question de cours construite]
\textbf{Introduction.} Entre 1884 et 1916, les sociétés camerounaises ne réagissent pas uniformément à la présence allemande : à côté de résistances armées parfois longues et acharnées, on observe des collaborations calculées. L'attitude des Camerounais dépend des rapports de force locaux, des intérêts des chefs et du degré de pression coloniale.

\textbf{I. Les résistances.} Les résistances primaires accompagnent la conquête : dès 1884-1885, une partie des Douala (partisans du roi Lock Priso à Hickorytown/Bonabéri) refuse le protectorat et affronte les troupes allemandes ; dans les années 1890-1900, les Bassa du chef Madola pratiquent la guérilla forestière, les Bafia se soulèvent contre les expéditions militaires, et dans le Nord les lamidats de l'Adamawa et les peuples du Logone opposent une résistance armée à la pénétration. Les résistances secondaires répondent à l'administration directe : révoltes fiscales, refus des corvées ferroviaires, et surtout l'opposition politique du roi Rudolf Duala Manga Bell contre l'expropriation des riverains du Wouri (1912-1914), qui lui vaut d'être pendu le 8 août 1914 avec Ngoso Din.

\textbf{II. Les collaborations.} D'autres acteurs choisissent l'accommodement. Les rois Bell et Akwa signent le traité de 1884 pour préserver leur commerce et prendre l'avantage sur leurs rivaux. Charles Atangana, interprète devenu chef supérieur des Ewondo-Béné, sert d'intermédiaire à l'administration : recrutement de porteurs et de travailleurs, collecte de l'impôt, médiation foncière. Le sultan Ibrahim Njoya de Foumban coopère avec les Allemands pour protéger l'autonomie de son royaume Bamoun. Des auxiliaires (soldats, interprètes, catéchistes) et des conversions au christianisme complètent ces stratégies d'adaptation.

\textbf{Conclusion.} Résistance et collaboration ne s'excluent pas absolument : un même chef peut coopérer puis s'opposer selon les circonstances. Ces stratégies diverses témoignent de la capacité des sociétés camerounaises à négocier, autant que possible, leur survie face à la contrainte coloniale, avant que la guerre de 1914-1916 ne mette fin à la domination allemande.

\textbf{Barème indicatif (sur 20) :} introduction et problématique 3 pts ; résistances primaires et secondaires distinguées avec au moins deux foyers nommés et datés 6 pts ; collaborations expliquées (motivations + formes) avec au moins deux exemples nommés 6 pts ; conclusion nuancée 3 pts ; langue et cohérence 2 pts.
\textbf{Points de vigilance :} la distinction résistances primaires / secondaires est explicitement demandée et doit structurer la partie I ; il faut au moins quatre exemples précis au total (noms, dates, lieux) ; éviter d'opposer résistants et collaborateurs de manière manichéenne : la conclusion doit restituer la complexité des stratégies africaines.
\end{corrige}

\newpage

\coursec{Fiche bilan}

\courssub{Carte mentale de synthèse}
\begin{popfigure}
\begin{center}\tcbox[colback=poporange,colframe=poporange,coltext=white,arc=9pt,boxsep=4pt,left=6pt,right=6pt]{\bfseries\small La période allemande (1884--1916)}\end{center}
\vspace{-2pt}
\begin{tcbraster}[raster columns=3,raster equal height=rows,raster column skip=6pt,raster row skip=6pt,enhanced,blankest]
\begin{tcolorbox}[enhanced,colback=white,colframe=poporange,boxrule=0.9pt,arc=3pt,title={Fondation du Kamerun (1884--1911)},fonttitle=\bfseries\scriptsize,coltitle=white,colbacktitle=poporange,left=4pt,right=4pt,top=2pt,bottom=2pt,boxsep=1pt,toptitle=1.5pt,bottomtitle=1.5pt,halign=left]
\begin{itemize}[nosep,leftmargin=*,itemsep=1pt,font=\scriptsize]
\item Traité Douala 12 juil. 1884
\item Conf. Berlin 1884--1885
\item Traité Maroc 1911 (N Kamerun)
\end{itemize}
\end{tcolorbox}
\begin{tcolorbox}[enhanced,colback=white,colframe=popgreen,boxrule=0.9pt,arc=3pt,title={Organisation administrative},fonttitle=\bfseries\scriptsize,coltitle=white,colbacktitle=popgreen,left=4pt,right=4pt,top=2pt,bottom=2pt,boxsep=1pt,toptitle=1.5pt,bottomtitle=1.5pt,halign=left]
\begin{itemize}[nosep,leftmargin=*,itemsep=1pt,font=\scriptsize]
\item Gouverneur (Buea/Yaoundé)
\item Kreises \& Stations
\item Chefs coutumiers auxiliaires
\end{itemize}
\end{tcolorbox}
\begin{tcolorbox}[enhanced,colback=white,colframe=poppurple,boxrule=0.9pt,arc=3pt,title={Organisation économique},fonttitle=\bfseries\scriptsize,coltitle=white,colbacktitle=poppurple,left=4pt,right=4pt,top=2pt,bottom=2pt,boxsep=1pt,toptitle=1.5pt,bottomtitle=1.5pt,halign=left]
\begin{itemize}[nosep,leftmargin=*,itemsep=1pt,font=\scriptsize]
\item Plantations (cacao, banane)
\item Travail forcé \& corvées
\item Chemins de fer Nord/Sud
\end{itemize}
\end{tcolorbox}
\begin{tcolorbox}[enhanced,colback=white,colframe=poppink,boxrule=0.9pt,arc=3pt,title={Résistances camerounaises},fonttitle=\bfseries\scriptsize,coltitle=white,colbacktitle=poppink,left=4pt,right=4pt,top=2pt,bottom=2pt,boxsep=1pt,toptitle=1.5pt,bottomtitle=1.5pt,halign=left]
\begin{itemize}[nosep,leftmargin=*,itemsep=1pt,font=\scriptsize]
\item Rudolf Duala Manga Bell
\item Martin-Paul Samba
\item Madola, Bakweri, Bafia, Bamoun
\end{itemize}
\end{tcolorbox}
\begin{tcolorbox}[enhanced,colback=white,colframe=popgold,boxrule=0.9pt,arc=3pt,title={Collaborations},fonttitle=\bfseries\scriptsize,coltitle=white,colbacktitle=popgold,left=4pt,right=4pt,top=2pt,bottom=2pt,boxsep=1pt,toptitle=1.5pt,bottomtitle=1.5pt,halign=left]
\begin{itemize}[nosep,leftmargin=*,itemsep=1pt,font=\scriptsize]
\item Chefs ralliés (Fon, Sultan)
\item Évangélisation (Baptistes, Pallotins)
\item Écoles \& santé
\end{itemize}
\end{tcolorbox}
\begin{tcolorbox}[enhanced,colback=white,colframe=popdark,boxrule=0.9pt,arc=3pt,title={Fin de la période (1914--1916)},fonttitle=\bfseries\scriptsize,coltitle=white,colbacktitle=popdark,left=4pt,right=4pt,top=2pt,bottom=2pt,boxsep=1pt,toptitle=1.5pt,bottomtitle=1.5pt,halign=left]
\begin{itemize}[nosep,leftmargin=*,itemsep=1pt,font=\scriptsize]
\item 1ère Guerre Mondiale
\item Campagne alliée
\item Partage Fr/UK 1916
\end{itemize}
\end{tcolorbox}
\end{tcbraster}

\legende{Carte mentale récapitulative de la leçon n°5 : La période allemande (Kamerun, 1884--1916).}
\end{popfigure}

\begin{bilan}
\courssub{Repères chronologiques et définitions clés}
\begin{center}
\begin{tabularx}{\linewidth}{|>{\bfseries}l|X|}\hline
\rowcolor{popblueL}
\textbf{Terme / Date} & \textbf{Contenu essentiel} \\\hline
12 juillet 1884 & Traité de protectorat signé par les rois Bell et Akwa (Douala) avec Nachtigal. Naissance du \cle{Kamerun}.\\\hline
1884--1885 & Conférence de Berlin : reconnaissance internationale des claims allemands.\\\hline
1885--1891 & Administration par la \cle{Compagnie du Kamerun} (chargée de l'exploitation).\\\hline
1891 & Le Kamerun devient \cle{colonie de peuplement} (Reichskolonie), administration directe.\\\hline
1907 & Transfert de la capitale de Buea (Buea) à \cle{Jaunde} (Yaoundé).\\\hline
4 nov. 1911 & Traité de Fès (accord franco-allemand) : \cle{Nouveau Kamerun} (Neukamerun) à l'Est et au Sud.\\\hline
1914--1916 & Campagne du Kamerun (1ère GM). Reddition de Mora (fév. 1916). Partage franco-britannique (Simon-Milner 1919).\\\hline
\cle{Indirect Rule} (allemand) & Administration par l'intermédiaire des chefs traditionnels (Fon, Sultan, Chef de canton) transformés en auxiliaires.\\\hline
\cle{Travail forcé} & Corvées obligatoires (routes, chemins de fer, plantations), base de l'économie coloniale.\\\hline
\end{tabularx}
\end{center}

\courssub{Acteurs principaux}
\begin{itemize}
\item \textbf{Allemand :} Gustav Nachtigal (fondateur), Julius von Soden (1er gouverneur civil), Jesko von Puttkamer (gouverneur "pacificateur", 1895-1907), Theodor Seitz (dernier gouverneur).
\item \textbf{Camerounais résistants :} \cle{Rudolf Duala Manga Bell} (Douala, exécuté 1914), \cle{Martin-Paul Samba} (Bulu, exécuté 1914), Madola (Bakweri), chefs Bafia, Bamoun (Njoya - collaboration puis tension), Bassa, Maka.
\item \textbf{Collaborateurs :} Sultan Ibrahim Njoya (Bamoun, écriture Shümom), chefs ralliés pour conserver prérogatives.
\item \textbf{Missionnaires :} Alfred Saker (Baptistes, précurseur), Mission pallotine (catholique), Mission de Bâle (presbytérienne) : écoles, santé, alphabétisation.
\end{itemize}

\courssub{Bilan administratif et économique (Méthode express : Tableau comparatif)}
\begin{center}
\begin{tabularx}{\linewidth}{|>{\bfseries}X|X|X|}\hline
\rowcolor{popblueL}
\textbf{Domaine} & \textbf{Réalisations allemandes} & \textbf{Conséquences pour les populations} \\\hline
Administratif & Division en \cle{Kreises} (circonscriptions) + stations militaires ; justice allemande (codes) ; police indigène. & Perte de souveraineté ; chefferies instrumentalisées ; justice inégale (code indigène vs code allemand). \\\hline
Économique & Plantations (cacao, café, banane, caoutchouc, huile de palme) ; 3 lignes ferroviaires (Nord, Midi, Nord-Ouest) ; ports (Douala, Kribi, Victoria). & Expropriation terres ; travail forcé/corvées ; migration interne ; économie d'exportation mono-culturelle. \\\hline
Social / Culturel & Écoles missionnaires (alphabétisation langues locales + allemand) ; hôpitaux (recherche sommeil, lèpre) ; écriture Bamoun (Njoya). & Émergence élite lettrée (clerks, catéchistes) ; diffusion christianisme ; perturbation structures sociales traditionnelles. \\\hline
\end{tabularx}
\end{center}

\courssub{Typologie des résistances (Savoir-faire : Classer/Justifier)}
\begin{itemize}
\item \textbf{Résistances précoces (1884--1900) :} Souvent locales, défensives (Bakweri, Bassa, Maka). Échec face à la puissance de feu.
\item \textbf{Résistances politiques et juridiques (1910--1914) :} \cle{Duala Manga Bell} (pétitions au Reichstag, avocat allemand) et \cle{Martin-Paul Samba} (tentative d'alliance internationale, achat d'armes). Répression féroce (exécution 8 août 1914).
\item \textbf{Résistance par l'adaptation :} Sultan Njoya (Bamoun) : modernisation de l'État, invention écriture, religion syncrétique, collaboration tactique.
\end{itemize}
\end{bilan}

\begin{aretenir}[Checklist avant le Bac -- Leçon 5 : La période allemande]
\begin{itemize}
\item $\square$ Je sais définir \cle{Kamerun}, \cle{Traité de protectorat (1884)}, \cle{Conférence de Berlin}, \cle{Neukamerun (1911)}.
\item $\square$ Je sais situer sur une frise : 1884 (traité), 1891 (colonie directe), 1907 (capitale Yaoundé), 1911 (Neukamerun), 1914-1916 (guerre/fin).
\item $\square$ Je sais expliquer l'organisation administrative : Gouverneur $\to$ Kreises $\to$ Stations $\to$ Chefs auxiliaires (Indirect Rule).
\item $\square$ Je sais décrire l'économie de plantation (produits, main-d'œuvre forcée, infrastructures ferroviaires/portuaires).
\item $\square$ Je connais les figures de \cle{Rudolf Duala Manga Bell} et \cle{Martin-Paul Samba} (motifs, actions, sort).
\item $\square$ Je sais différencier \textbf{résistance armée}, \textbf{résistance juridique/politique} et \textbf{collaboration/adaptation} (ex: Njoya).
\item $\square$ Je sais analyser le rôle des missions (éducation, santé, alphabétisation) comme acteurs de la colonisation et vecteurs de modernité.
\item $\square$ Je sais expliquer les causes de la défaite allemande (1ère GM, encerclement allié, supériorité numérique/logistique).
\item $\square$ Je connais les conséquences du partage 1916 (mandats SDN : Cameroun français / British Cameroons).
\item $\square$ Je suis capable de \textbf{localiser sur carte muette} : Douala, Yaoundé, Buea, Kribi, Victoria, Garoua, Ngaoundéré, lignes de chemin de fer, limites Neukamerun.
\item $\square$ Je sais rédiger une introduction, un développement en 2 ou 3 parties équilibrées et une conclusion ouverte sur le devenir du territoire.
\end{itemize}
\end{aretenir}

\findecours

\end{document}
